% Énergie potentielle et énergie mécanique — Physique 1ère D — Cours signé OnBuch+
% Programme officiel MINESEC. Compiler avec : tectonic N05-energie-potentielle-mecanique.tex
\newcommand{\DOCMATIERE}{Physique}
\newcommand{\DOCNIVEAU}{1ère D}
\newcommand{\DOCMODULE}{Module 2 : Mouvements et interactions mécaniques}
\newcommand{\DOCLECON}{N05}
\newcommand{\DOCTITRE}{Énergie potentielle et énergie mécanique}
% OnBuch+ — préambule « cours signés » ENRICHI (compilé avec tectonic / XeLaTeX)
% Système visuel « Pop » d'OnBuch+ : fond crème (jamais blanc), bordures encre
% épaisses, ombres « dures » décalées, titres Archivo Black en capitales.
% Version MATHÉMATIQUES : graphes (pgfplots/tikz), boîtes pédagogiques
% (définition, propriété, méthode, attention, le savais-tu, exercice résolu,
% exercice, TP, bilan), page de garde et signature OnBuch+ ancrée en bas.
%
% Macros attendues dans main.tex AVANT \input{preamble} :
%   \DOCMATIERE  \DOCNIVEAU  \DOCMODULE  \DOCLECON (numéro)  \DOCTITRE
\documentclass[11pt]{article}

\usepackage{fontspec}
\usepackage[french]{babel}
\usepackage[a4paper,margin=2cm,top=3.4cm,bottom=2.8cm,headheight=1.4cm,footskip=1.3cm]{geometry}
\usepackage{amsmath,amssymb,amsfonts}
\usepackage{mathtools} % \xmapsto, \coloneqq... souvent utilisés par les modèles
\usepackage{cancel} % simplifications barrées (bilans de forces)
% \abs{z} (module, valeur absolue) et \norm{u} (norme), utilisés par les modèles
\providecommand{\abs}{}\let\abs\relax\DeclarePairedDelimiter{\abs}{\lvert}{\rvert}
\providecommand{\norm}{}\let\norm\relax\DeclarePairedDelimiter{\norm}{\lVert}{\rVert}
\usepackage[table]{xcolor} % [table] : fournit aussi \rowcolors (alternance de lignes)
\usepackage{pagecolor}
\usepackage{tikz}
\usetikzlibrary{arrows.meta,positioning,calc,shapes.geometric,shapes.misc,shapes.arrows,decorations.pathmorphing,decorations.pathreplacing,decorations.markings,patterns,shadows,mindmap,trees,fit,backgrounds,matrix,intersections,angles,quotes}
\usepackage{pgfplots}
\usepackage[version=4]{mhchem} % équations nucléaires \ce{^{235}_{92}U}
\usepackage[european,siunitx]{circuitikz} % schémas électriques
\pgfplotsset{compat=1.18}
\usepgfplotslibrary{fillbetween,groupplots} % aires entre courbes (calcul intégral)
\usepackage{enumitem}
\usepackage{fancyhdr}
% [most] charge toutes les bibliothèques tcolorbox (listings, minted, raster,
% documentation...) et gonfle inutilement la mémoire TeX sur un document long
% (~300 boîtes) : on ne charge que ce qui est réellement utilisé.
\usepackage[skins,breakable]{tcolorbox}
\usepackage{graphicx}
\usepackage{colortbl}
\usepackage{booktabs}
\usepackage{tabularx}
\usepackage{diagbox} % case d'en-tête diagonale des tableaux à double entrée
% Colonnes extensibles centrée / gauche / droite (C, L, R), souvent utilisées par les modèles
\newcolumntype{C}{>{\centering\arraybackslash}X}
\newcolumntype{L}{>{\raggedright\arraybackslash}X}
\newcolumntype{R}{>{\raggedleft\arraybackslash}X}
\usepackage{array}
\usepackage{multicol}
\usepackage{siunitx}
% parse-numbers=false : accepte aussi \SI{2,5 \times 10^{-3}}{...}
\sisetup{locale=FR,output-decimal-marker={,},group-minimum-digits=4,parse-numbers=false}
\DeclareSIUnit\ohm{\ensuremath{\symup{\Omega}}} % U+2126 absent de la police
\DeclareSIUnit\division{div} % réglages d'oscilloscope (ms/div, V/div)
\DeclareSIUnit\tr{tr} % tours (vitesse de rotation, ex. \SI{1500}{\tr\per\minute})
\DeclareSIUnit\ch{ch} % cheval-vapeur (puissance moteur, unité usuelle hors SI)
\DeclareSIUnit\franccfa{FCFA} % franc CFA (coûts, contexte camerounais)
\DeclareSIUnit\dioptre{\ensuremath{\delta}} % vergence d'une lentille (dioptrie)
\usepackage{qrcode}
% \degree : commande fréquemment utilisée par les modèles pour un angle ou une
% température en dehors de \SI{}{} (non fournie par les paquets chargés).
\providecommand{\degree}{^{\circ}}
% \qed (fin de démonstration), utilisé par les modèles sans amsthm
\providecommand{\qed}{\hfill$\square$}
% \barre : double barre des valeurs interdites dans un tableau de variations (tkz-tab)
\providecommand{\barre}{\ensuremath{\|}}
% environnement proof (amsthm non chargé) utilisé par les modèles
\ifdefined\proof\else\newenvironment{proof}[1][Démonstration]{\par\noindent\textit{#1.}\ }{\qed\par}\fi
\usepackage{lastpage}
\usepackage{hyperref}
\hypersetup{hidelinks}

% ============================================================
% Palette « Pop » OnBuch+ — mêmes hex que la classe P (plus_ui.dart)
% ============================================================
\definecolor{poporange}{HTML}{F59321}
\definecolor{popdark}{HTML}{E07A0C}
\definecolor{popcream}{HTML}{FFF6E8}
\definecolor{popink}{HTML}{1C1714}
\definecolor{poppurple}{HTML}{7A5AE0}
\definecolor{popgreen}{HTML}{2E9E6A}
\definecolor{poppink}{HTML}{E0457B}
\definecolor{popblue}{HTML}{2F7DE1}
\definecolor{popgold}{HTML}{C98A12}
\definecolor{popmuted}{HTML}{8A7F76}
\definecolor{popline}{HTML}{EADFD2}
\definecolor{popgray}{HTML}{8A7F76}
\colorlet{popgrey}{popgray}
% Alias tolérants (le modèle nomme parfois les couleurs différemment)
\colorlet{popwhite}{white}
\colorlet{popblack}{popink}
\colorlet{popred}{poppink}
\colorlet{popyellow}{popgold}
% Teintes claires pour fonds de tableaux / zones
\colorlet{popblueL}{popblue!10!white}
\colorlet{poporangeL}{poporange!14!white}
\colorlet{popgreenL}{popgreen!12!white}
\colorlet{poppurpleL}{poppurple!12!white}
\colorlet{poppinkL}{poppink!12!white}
\colorlet{popgoldL}{popgold!14!white}

\pagecolor{popcream}

% ============================================================
% Typographie
% ============================================================
\defaultfontfeatures{Ligatures=TeX}
\setmainfont{PlusJakartaSans-Medium.ttf}[Path=../../../fonts/,BoldFont=PlusJakartaSans-Bold.ttf]
% Maths en sans-serif (Fira Math) avec chiffres et lettres droites de Plus
% Jakarta, pour que les formules se fondent dans le texte.
\usepackage[math-style=ISO,bold-style=ISO]{unicode-math}
\setmathfont{FiraMath-Regular.otf}[Path=../../../fonts/]
\setmathfont{PlusJakartaSans-Medium.ttf}[Path=../../../fonts/,range={up/num,up/{Latin,latin},\mathup}]
\usepackage{icomma} % virgule décimale sans espace en mode math
% Caractères mathématiques Unicode absents des polices de texte (Plus Jakarta,
% Archivo Black) : rendus par leur équivalent mathématique, en texte comme en
% formule — sinon ils s'affichent en carré vide (ex. « Arithmétique dans ℤ »).
\usepackage{newunicodechar}
\newunicodechar{ℝ}{\ensuremath{\mathbb{R}}}
\newunicodechar{ℤ}{\ensuremath{\mathbb{Z}}}
\newunicodechar{ℂ}{\ensuremath{\mathbb{C}}}
\newunicodechar{ℕ}{\ensuremath{\mathbb{N}}}
\newunicodechar{ℚ}{\ensuremath{\mathbb{Q}}}
\newunicodechar{𝔻}{\ensuremath{\mathbb{D}}}
\newunicodechar{α}{\ensuremath{\alpha}}
\newunicodechar{β}{\ensuremath{\beta}}
\newunicodechar{γ}{\ensuremath{\gamma}}
\newunicodechar{δ}{\ensuremath{\delta}}
\newunicodechar{ε}{\ensuremath{\varepsilon}}
\newunicodechar{θ}{\ensuremath{\theta}}
\newunicodechar{λ}{\ensuremath{\lambda}}
\newunicodechar{μ}{\ensuremath{\mu}}
\newunicodechar{σ}{\ensuremath{\sigma}}
\newunicodechar{τ}{\ensuremath{\tau}}
\newunicodechar{φ}{\ensuremath{\varphi}}
\newunicodechar{ψ}{\ensuremath{\psi}}
\newunicodechar{ω}{\ensuremath{\omega}}
\newunicodechar{Σ}{\ensuremath{\Sigma}}
% majuscules produites par \MakeUppercase dans les titres (α-aminés -> Α)
\newunicodechar{Α}{\ensuremath{\alpha}}
\newunicodechar{Β}{\ensuremath{\beta}}
\newunicodechar{Γ}{\ensuremath{\Gamma}}
\newunicodechar{Δ}{\ensuremath{\Delta}}
\newunicodechar{Ε}{\ensuremath{\varepsilon}}
\newunicodechar{Θ}{\ensuremath{\Theta}}
\newunicodechar{Λ}{\ensuremath{\Lambda}}
\newunicodechar{Μ}{\ensuremath{\mu}}
\newunicodechar{Τ}{\ensuremath{\tau}}
\newunicodechar{Φ}{\ensuremath{\Phi}}
\newunicodechar{Ψ}{\ensuremath{\Psi}}
\newunicodechar{Ω}{\ensuremath{\Omega}}
\newunicodechar{↦}{\ensuremath{\mapsto}}
\newunicodechar{∈}{\ensuremath{\in}}
\newunicodechar{∉}{\ensuremath{\notin}}
\newunicodechar{∧}{\ensuremath{\wedge}}
\newunicodechar{⇔}{\ensuremath{\Leftrightarrow}}
\newunicodechar{⟺}{\ensuremath{\Longleftrightarrow}}
\newunicodechar{⇒}{\ensuremath{\Rightarrow}}
\newunicodechar{⟹}{\ensuremath{\Longrightarrow}}
\newunicodechar{∘}{\ensuremath{\circ}}
\newunicodechar{∪}{\ensuremath{\cup}}
\newunicodechar{∩}{\ensuremath{\cap}}
\newunicodechar{≡}{\ensuremath{\equiv}}
\newunicodechar{⊂}{\ensuremath{\subset}}
\newunicodechar{⊥}{\ensuremath{\perp}}
\newunicodechar{∃}{\ensuremath{\exists}}
\newunicodechar{∀}{\ensuremath{\forall}}
\newunicodechar{∛}{\ensuremath{\sqrt[3]{\,}}}
\newunicodechar{∜}{\ensuremath{\sqrt[4]{\,}}}
\newunicodechar{⃗}{\ensuremath{{}^{\to}}}
\newunicodechar{ⁿ}{\ifmmode^{n}\else\textsuperscript{\ensuremath{n}}\fi}
\newunicodechar{ˣ}{\ifmmode^{x}\else\textsuperscript{\ensuremath{x}}\fi}
\newunicodechar{ᵇ}{\ifmmode^{b}\else\textsuperscript{\ensuremath{b}}\fi}
\newunicodechar{ʳ}{\ifmmode^{r}\else\textsuperscript{\ensuremath{r}}\fi}
\newunicodechar{ᵘ}{\ifmmode^{u}\else\textsuperscript{\ensuremath{u}}\fi}
\newunicodechar{ʸ}{\ifmmode^{y}\else\textsuperscript{\ensuremath{y}}\fi}
\newunicodechar{ᵃ}{\ifmmode^{a}\else\textsuperscript{\ensuremath{a}}\fi}
\newunicodechar{ᵉ}{\ifmmode^{e}\else\textsuperscript{\ensuremath{e}}\fi}
\newunicodechar{ᵏ}{\ifmmode^{k}\else\textsuperscript{\ensuremath{k}}\fi}
\newunicodechar{ᵖ}{\ifmmode^{p}\else\textsuperscript{\ensuremath{p}}\fi}
\newunicodechar{ᴺ}{\ifmmode^{N}\else\textsuperscript{\ensuremath{N}}\fi}
\newunicodechar{ᶜ}{\ifmmode^{c}\else\textsuperscript{\ensuremath{c}}\fi}
\newunicodechar{ʰ}{\ifmmode^{h}\else\textsuperscript{\ensuremath{h}}\fi}
\newunicodechar{ᵠ}{\ifmmode^{q}\else\textsuperscript{\ensuremath{q}}\fi}
\newunicodechar{ₙ}{\ifmmode_{n}\else\textsubscript{\ensuremath{n}}\fi}
\newunicodechar{ₐ}{\ifmmode_{a}\else\textsubscript{\ensuremath{a}}\fi}
\newunicodechar{ᵢ}{\ifmmode_{i}\else\textsubscript{\ensuremath{i}}\fi}
\newunicodechar{ᵣ}{\ifmmode_{r}\else\textsubscript{\ensuremath{r}}\fi}
\newunicodechar{ₖ}{\ifmmode_{k}\else\textsubscript{\ensuremath{k}}\fi}
\newunicodechar{ᵦ}{\ifmmode_{\beta}\else\textsubscript{\ensuremath{\beta}}\fi}
\newunicodechar{ₓ}{\ifmmode_{x}\else\textsubscript{\ensuremath{x}}\fi}
\newfontfamily\popdisplay{ArchivoBlack-Regular.ttf}[Path=../../../fonts/]
\newfontfamily\poplogo{SpaceGrotesk-Bold.ttf}[Path=../../../fonts/]
\newfontfamily\popsemi{PlusJakartaSans-SemiBold.ttf}[Path=../../../fonts/]
\newfontfamily\popxbold{PlusJakartaSans-ExtraBold.ttf}[Path=../../../fonts/]
\newfontfamily\popxboldls{PlusJakartaSans-ExtraBold.ttf}[Path=../../../fonts/,LetterSpace=3.0]

\setlist[enumerate]{leftmargin=1.7em,itemsep=5pt,topsep=4pt}
\setlist[itemize]{leftmargin=1.7em,itemsep=4pt,topsep=4pt,label={\color{poporange}\textbullet}}
\setlength{\parindent}{0pt}
\setlength{\parskip}{5pt}
\renewcommand{\arraystretch}{1.25}

% pgfplots : style Pop par défaut
\pgfplotsset{
  every axis/.append style={axis line style={popink,line width=1pt},tick style={popink},
    grid style={popline,line width=0.5pt},label style={font=\small},tick label style={font=\footnotesize}},
  popcurve/.style={poporange,line width=1.8pt,no markers,smooth},
}
% Même style disponible dans un \draw TikZ simple (hors pgfplots)
\tikzset{popcurve/.style={poporange,line width=1.8pt}}
\tikzset{popgrid/.style={popline,very thin}}
\colorlet{popcurve}{poporange} % le nom du style est parfois utilisé comme couleur

% ============================================================
% Pill / squiggle
% ============================================================
\newcommand{\poppill}[3]{%
  \tikz[baseline=-0.55ex]\node[draw=popink,line width=1.2pt,rounded corners=2.6pt,%
    fill=#2,inner xsep=6.5pt,inner ysep=3pt]{%
    \popxboldls\fontsize{7}{7}\selectfont\color{#3}\MakeUppercase{#1}};%
}
\newcommand{\popsquiggle}[1]{%
  \tikz[baseline=-0.4ex]{
    \draw[#1,line width=2.6pt,line cap=round]
      plot [smooth] coordinates {(0,0.09) (0.34,0.01) (0.68,0.21) (1.02,0.01) (1.36,0.10)};
  }%
}
% Mot-clé surligné (marqueur orange sous le texte)
\newcommand{\cle}[1]{{\popxbold\color{popdark}#1}}

% ============================================================
% En-tête / pied
% ============================================================
\pagestyle{fancy}
\fancyhf{}
\renewcommand{\headrulewidth}{0pt}
\renewcommand{\footrulewidth}{0pt}
\lhead{\raisebox{-0.15ex}{\poplogo\fontsize{13}{13}\selectfont\color{popink}OnBuch\color{poporange}+}}
\rhead{\poppill{\DOCMATIERE\ \textbullet\ \DOCNIVEAU\ \textbullet\ Leçon \DOCLECON}{white}{popink}}
\lfoot{\popxbold\fontsize{7.6}{7.6}\selectfont\color{popmuted}\MakeUppercase{Cours signé OnBuch+}\ \textbullet\ {\popsemi\DOCTITRE}}
\rfoot{\tikz[baseline=-0.6ex]\node[rounded corners=8.5pt,fill=popink,minimum height=17pt,minimum width=17pt,inner xsep=5pt,inner sep=0]{\popxbold\fontsize{8}{8}\selectfont\color{popcream}\thepage\,/\,\pageref*{LastPage}};}

% ============================================================
% Titres
% ============================================================
\newcommand{\chaptitre}[2]{%
  \begin{tcolorbox}[enhanced,colback=poporange,colframe=popink,boxrule=2.6pt,arc=11pt,
      drop shadow=popink,left=15pt,right=15pt,top=12pt,bottom=11pt,before skip=3pt,after skip=18pt]
    \poppill{#2}{popink}{popcream}\par\vspace{9pt}
    {\raggedright\hyphenpenalty=10000\exhyphenpenalty=10000\popdisplay\fontsize{22}{24}\selectfont\color{popcream}\MakeUppercase{#1}\par}
  \end{tcolorbox}
}
% Section numérotée : pastille orange avec le numéro + titre Archivo
\newcounter{popsec}
\newcommand{\coursec}[1]{%
  \stepcounter{popsec}\setcounter{popsub}{0}%
  \par\vspace{16pt}\noindent
  \begin{minipage}{\linewidth}
  \tikz[baseline=-0.7ex]\node[draw=popink,line width=1.6pt,fill=poporange,rounded corners=4pt,minimum size=22pt,inner sep=0]{\popdisplay\fontsize{12}{12}\selectfont\color{popcream}\thepopsec};%
  \hspace{8pt}{\popdisplay\fontsize{15}{17}\selectfont\color{popink}\MakeUppercase{#1}}\par
  \vspace{4pt}\noindent{\color{poporange}\rule{36pt}{3.6pt}}
  \end{minipage}\par\nopagebreak\vspace{8pt}
}
\newcounter{popsub}
\newcommand{\courssub}[1]{%
  \stepcounter{popsub}%
  \par\vspace{9pt}\noindent{\popxbold\fontsize{11.5}{13}\selectfont\color{popink}{\color{poporange}\thepopsec.\thepopsub}\hspace{6pt}#1}\par\nopagebreak\vspace{4pt}
}

% ============================================================
% Boîtes pédagogiques — même recette : fond blanc, bordure épaisse colorée,
% ombre dure encre, badge plein. Titre optionnel [#1].
% ============================================================
\tcbset{popbase/.style={breakable,enhanced,colback=white,boxrule=2.3pt,arc=9pt,
  shadow={3pt}{-3pt}{0pt}{popink},left=14pt,right=14pt,top=11pt,bottom=11pt,before skip=11pt,after skip=15pt,
  fontupper=\popsemi\fontsize{10.6}{14.5}\selectfont\color{popink},
  fontlower=\popsemi\fontsize{10.6}{14.5}\selectfont\color{popink}}}
% Coche dessinée (le glyphe ✓ n'existe pas dans les polices du document)
\newcommand{\popcheck}[1]{\tikz[baseline=-0.5ex]\draw[#1,line width=1.6pt,line cap=round,line join=round](-0.13,0.0)--(-0.04,-0.09)--(0.13,0.1);}
\providecommand{\coche}{\popcheck{popgreen}}
\providecommand{\resultat}[1]{\par\smallskip\noindent\fcolorbox{popgreen}{popgreenL}{\parbox{\dimexpr\linewidth-2\fboxsep-2\fboxrule\relax}{\textbf{Résultat.} #1}}\par\smallskip}
\newcommand{\popboxhead}[3]{\poppill{#1}{#2}{white}\ifx\relax#3\relax\else\hspace{6pt}{\popxbold\fontsize{10.6}{12}\selectfont\color{popink}#3}\fi\par\vspace{7pt}}

\newtcolorbox{aretenir}[1][]{popbase,colframe=popgreen,before upper={\popboxhead{À retenir}{popgreen}{#1}}}
\newtcolorbox{exemplebox}[1][]{popbase,colframe=poporange,before upper={\popboxhead{Exemple}{poporange}{#1}}}
\newtcolorbox{definition}[1][]{popbase,colframe=popblue,before upper={\popboxhead{Définition}{popblue}{#1}}}
\newtcolorbox{propriete}[1][]{popbase,colframe=poppurple,before upper={\popboxhead{Propriété}{poppurple}{#1}}}
\newtcolorbox{methode}[1][]{popbase,colframe=popgold,before upper={\popboxhead{Méthode}{popgold}{#1}}}
\newtcolorbox{attention}[1][]{popbase,colframe=poppink,before upper={\popboxhead{Attention}{poppink}{#1}}}
\newtcolorbox{experience}[1][]{popbase,colframe=popblue,colback=popblueL,before upper={\popboxhead{Expérience}{popblue}{#1}}}
% « Le savais-tu ? » : carte violette pleine (contexte, culture, Cameroun)
\newtcolorbox{savaistu}[1][]{popbase,colback=poppurple,colframe=popink,
  fontupper=\popsemi\fontsize{10.4}{14.2}\selectfont\color{white},
  before upper={\poppill{Le savais-tu\,?}{white}{poppurple}\ifx\relax#1\relax\else\hspace{6pt}{\popxbold\color{white}#1}\fi\par\vspace{7pt}}}
% Exercice résolu : énoncé en haut, \tcblower puis la solution sur fond crème
\newcounter{popexo}\newcounter{popexr}
\newtcolorbox{exoresolu}[1][]{popbase,colframe=poporange,colbacklower=popcream,
  segmentation style={popink,line width=1.2pt,dashed},
  before upper={\stepcounter{popexr}\popboxhead{Exercice résolu \thepopexr}{poporange}{#1}},
  before lower={\poppill{Solution}{popink}{popcream}\par\vspace{6pt}}}
% Exercice (non résolu en place) — la correction est dans la partie Corrigés
\newtcolorbox{exercice}[1][]{popbase,colframe=popink,
  before upper={\stepcounter{popexo}\popboxhead{Exercice \thepopexo}{popink}{#1}}}
% Corrigé
\newtcolorbox{corrige}[1][]{popbase,colframe=popgreen,colback=popgreenL,
  before upper={\popboxhead{Corrigé}{popgreen}{#1}}}
% Objectifs / prérequis / bilan
\newtcolorbox{objectifs}{popbase,colframe=popink,colback=poporangeL,
  before upper={\popboxhead{Objectifs de la leçon}{popink}{}}}
\newtcolorbox{prerequis}{popbase,colframe=popink,colback=popblueL,
  before upper={\popboxhead{Prérequis}{popblue}{}}}
\newtcolorbox{bilan}{popbase,colframe=popink,colback=white,boxrule=2.8pt,
  before upper={\popboxhead{Fiche bilan}{poporange}{L'essentiel à connaître pour l'examen}}}
% Figure encadrée avec légende
\newtcolorbox{popfigure}[1][]{enhanced,breakable=false,colback=white,colframe=popline,boxrule=1.4pt,arc=8pt,
  left=10pt,right=10pt,top=10pt,bottom=8pt,before skip=10pt,after skip=12pt,halign=center,
  lower separated=false,halign lower=center,
  fontlower=\popsemi\fontsize{9}{11.5}\selectfont\color{popmuted},
  #1}
\newcommand{\legende}[1]{\par\vspace{6pt}{\popsemi\fontsize{9}{11.5}\selectfont\color{popmuted}#1}}

% ============================================================
% Signature OnBuch+
% ============================================================
\newcommand{\poplogoinline}[1]{{\poplogo\fontsize{#1}{#1}\selectfont\color{popink}OnBuch\color{poporange}+}}
\newcommand{\signatureonbuch}{%
  \begin{tcolorbox}[enhanced,colback=popink,colframe=popink,boxrule=2.2pt,arc=10pt,
      drop shadow=poporange,left=14pt,right=14pt,top=10pt,bottom=10pt,before skip=0pt,after skip=0pt]
    \tikz[baseline=-0.7ex]\node[circle,fill=popgreen,draw=popcream,line width=1.3pt,minimum size=22pt,inner sep=0]{\popcheck{white}};%
    \hspace{9pt}\begin{minipage}[c]{0.62\linewidth}
      {\popxbold\fontsize{12}{13}\selectfont\color{popcream}Signé\ \textbullet\ {\poplogo OnBuch\color{poporange}+}}\par\vspace{2pt}
      {\raggedright\popsemi\fontsize{8}{10.5}\selectfont\color{popline}\MakeUppercase{Équipe pédagogique OnBuch+}\\\MakeUppercase{Conforme au programme officiel MINESEC}\par}
    \end{minipage}\hfill
    {\poplogo\fontsize{22}{22}\selectfont\color{popcream}OnBuch\color{poporange}+}
  \end{tcolorbox}%
}

% ============================================================
% Page de garde
% #1 titre · #2 matière · #3 niveau · #4 module · #5 n° leçon · #6 durée ·
% #7 description / objectifs (texte ou itemize)
% ============================================================
\newcommand{\pagegarde}[7]{%
  \thispagestyle{empty}
  \newgeometry{a4paper,margin=1.7cm,top=1.6cm,bottom=1.4cm}%
  \begin{tikzpicture}[remember picture,overlay]
    \fill[poporange!18] ([xshift=-3.2cm,yshift=-3.6cm]current page.north east) circle (4.6cm);
    \fill[poppurple!14] ([xshift=2.4cm,yshift=7.6cm]current page.south west) circle (3.8cm);
  \end{tikzpicture}%
  {\poplogo\fontsize{30}{30}\selectfont\color{popink}OnBuch\color{poporange}+}\hfill
  \raisebox{0.6ex}{\poppill{Cours signé \textbullet\ Premium}{popink}{popcream}}\par
  \vspace{1pt}\popsquiggle{poppurple}\par
  \vspace{8pt}
  \begin{tcolorbox}[enhanced,colback=poporange,colframe=popink,boxrule=2.8pt,arc=13pt,
      drop shadow=popink,left=18pt,right=18pt,top=12pt,bottom=13pt,after skip=10pt]
    \poppill{#2 \textbullet\ #3}{popink}{popcream}\hspace{5pt}\poppill{#4}{white}{popink}\par\vspace{8pt}
    {\popdisplay\fontsize{14}{14}\selectfont\color{popink}LEÇON #5}\par\vspace{4pt}
    {\raggedright\hyphenpenalty=10000\exhyphenpenalty=10000\popdisplay\fontsize{25}{27}\selectfont\color{popcream}\MakeUppercase{#1}\par}
    \vspace{8pt}
    {\popxbold\fontsize{9.5}{9.5}\selectfont\color{popink}\MakeUppercase{Durée conseillée : #6}}
  \end{tcolorbox}
  \begin{tcolorbox}[enhanced,colback=white,colframe=popink,boxrule=2.2pt,arc=10pt,
      drop shadow=popink,left=16pt,right=16pt,top=10pt,bottom=10pt,after skip=8pt]
    \poppill{Dans cette leçon}{poppurple}{white}\par\vspace{5pt}
    \popsemi\fontsize{9.3}{12.6}\selectfont\color{popink}#7
  \end{tcolorbox}
  \noindent\begin{minipage}[t]{0.487\linewidth}
    \begin{tcolorbox}[enhanced,colback=popgreen,colframe=popink,boxrule=2.2pt,arc=10pt,
        drop shadow=popink,left=11pt,right=11pt,top=10pt,bottom=10pt,height=3.1cm,valign=center]
      \tcbox[colback=white,colframe=popink,boxrule=1.3pt,arc=5pt,boxsep=0pt,nobeforeafter,
        left=3pt,right=3pt,top=3pt,bottom=3pt,valign=center]{\qrcode[height=1.9cm]{https://play.google.com/store/apps/details?id=com.luvvix.onbuch&hl=fr}}%
      \hspace{8pt}\begin{minipage}[c]{0.5\linewidth}
        {\popdisplay\fontsize{11}{12.5}\selectfont\color{white}\MakeUppercase{Télécharge OnBuch}}\par\vspace{4pt}
        {\raggedright\popsemi\fontsize{8.4}{11}\selectfont\color{white}Gratuit \textbullet\ Google Play\\Tuteur IA, annales, concours, résultats\par}
      \end{minipage}
    \end{tcolorbox}
  \end{minipage}\hfill
  \begin{minipage}[t]{0.487\linewidth}
    \begin{tcolorbox}[enhanced,colback=poppurple,colframe=popink,boxrule=2.2pt,arc=10pt,
        drop shadow=popink,left=11pt,right=11pt,top=10pt,bottom=10pt,height=3.1cm,valign=center]
      \tikz[baseline=-0.7ex]\node[circle,fill=white,draw=popink,line width=1.3pt,minimum size=1.6cm,inner sep=0]{\popdisplay\fontsize{24}{24}\selectfont\color{poppurple}+};%
      \hspace{8pt}\begin{minipage}[c]{0.6\linewidth}
        {\popdisplay\fontsize{11}{12.5}\selectfont\color{white}\MakeUppercase{Passe à OnBuch+}}\par\vspace{4pt}
        {\raggedright\popsemi\fontsize{8.4}{11}\selectfont\color{white}Tous les cours signés, Tuteur IA illimité, fiches PDF — onglet OnBuch+ de l'app\par}
      \end{minipage}
    \end{tcolorbox}
  \end{minipage}
  \vspace{8pt}
  \popcontenu
  \vfill
  \signatureonbuch
  \restoregeometry
  \newpage
}

% Rangée « Ce cours contient » (page de garde)
\newcommand{\poptile}[3]{% #1 couleur #2 gros texte #3 légende
  \begin{tcolorbox}[enhanced,colback=#1,colframe=popink,boxrule=2pt,arc=9pt,shadow={2.5pt}{-2.5pt}{0pt}{popink},
      left=6pt,right=6pt,top=9pt,bottom=9pt,halign=center,valign=center,height=1.75cm,width=\linewidth,nobeforeafter]
    {\popdisplay\fontsize{12.5}{13}\selectfont\color{white}\MakeUppercase{#2}}\par\vspace{3pt}
    {\centering\popsemi\fontsize{7.8}{9.5}\selectfont\color{white}#3\par}
  \end{tcolorbox}}
\newcommand{\popcontenu}{%
  \par\noindent{\popxboldls\fontsize{7.5}{7.5}\selectfont\color{popmuted}CE COURS SIGNÉ CONTIENT}\par\vspace{7pt}
  \noindent\begin{minipage}[t]{0.235\linewidth}\poptile{popblue}{Cours}{complet, illustré, pas à pas}\end{minipage}\hfill
  \begin{minipage}[t]{0.235\linewidth}\poptile{poporange}{Méthodes}{et exercices résolus}\end{minipage}\hfill
  \begin{minipage}[t]{0.235\linewidth}\poptile{poppink}{Exercices}{type examen + corrigés}\end{minipage}\hfill
  \begin{minipage}[t]{0.235\linewidth}\poptile{popgreen}{Fiche bilan}{l'essentiel à retenir}\end{minipage}\par}

% Sommaire
\newtcolorbox{sommaire}{enhanced,colback=white,colframe=popink,boxrule=2.2pt,arc=10pt,
  drop shadow=popink,left=18pt,right=18pt,top=13pt,bottom=13pt,before skip=6pt,after skip=20pt,
  before upper={\poppill{Sommaire}{poppurple}{white}\par\vspace{9pt}\popsemi\fontsize{10.6}{14.5}\selectfont\color{popink}}}

% Signature de fin de cours — poussée tout en bas de la dernière page
\newcommand{\findecours}{%
  \par\vfill\nopagebreak
  \begin{center}\popsquiggle{poporange}\end{center}\vspace{4pt}
  \signatureonbuch
}
\AtBeginDocument{\def\llbracket{\mathopen{[\mkern-3.5mu[}}\def\rrbracket{\mathclose{]\mkern-3.5mu]}}} % intervalles d'entiers (glyphe ⟦ absent de la police)
% Glyphes absents de Fira Math : redéfinis à partir de glyphes disponibles
\AtBeginDocument{%
  \def\setminus{\mathbin{\backslash}}\let\smallsetminus\setminus
  \def\vdots{\mathord{\vbox{\baselineskip3.2pt\lineskiplimit0pt\kern4pt\hbox{.}\hbox{.}\hbox{.}}}}%
  \def\ddots{\mathinner{\mkern1mu\raise6.5pt\hbox{.}\mkern2mu\raise3.6pt\hbox{.}\mkern2mu\raise0.7pt\hbox{.}\mkern1mu}}%
  \def\triangle{\mathord{\scalebox{0.9}{$\symup{\Delta}$}}}%
  \def\nabla{\mathord{\raisebox{\depth}{\scalebox{1}[-1]{$\symup{\Delta}$}}}}%
  \def\checkmark{\popcheck{popdark}}%
  \def\star{\mathbin{\ast}}\def\bigstar{\mathord{\ast}}%
  \def\diamond{\mathbin{\scalebox{0.6}{$\Diamond$}}}%
  \def\bigcup{\mathop{\mathchoice{\raisebox{-0.3em}{\scalebox{1.7}{$\cup$}}}{\scalebox{1.25}{$\cup$}}{\cup}{\cup}}\displaylimits}%
  \def\bigcap{\mathop{\mathchoice{\raisebox{-0.3em}{\scalebox{1.7}{$\cap$}}}{\scalebox{1.25}{$\cap$}}{\cap}{\cap}}\displaylimits}%
  \def\lesseqqgtr{\mathrel{\lessgtr}}\def\bigoplus{\mathop{\oplus}\displaylimits}%
}
\providecommand{\ssi}{si et seulement si} % abréviation fréquente
\AtBeginDocument{\providecommand{\wideparen}{\overparen}} % arc de cercle

%% FIGFIX-BEGIN v5 — correctifs globaux des figures (docs/onbuch-generation/tools/figfix)
\usepackage{adjustbox}\usepackage{etoolbox}\tcbuselibrary{raster}
% toute figure TikZ/pgfplots plus large que la ligne est réduite à la largeur disponible
\BeforeBeginEnvironment{tikzpicture}{\begin{adjustbox}{max width=\linewidth}}
\AfterEndEnvironment{tikzpicture}{\end{adjustbox}}
% légendes pgfplots lisibles par-dessus les courbes
\pgfplotsset{every axis/.append style={legend style={fill=white,fill opacity=0.92,text opacity=1,draw=black!15,inner sep=3pt}}}
% étiquettes d'arêtes (syntaxe edge["..."]) avec fond blanc
\tikzset{every edge quotes/.append style={fill=white,inner sep=1.2pt}}
%% FIGFIX-END
\begin{document}

\pagegarde{Énergie potentielle et énergie mécanique}{Physique}{1ère D}{Module 2 : Mouvements et interactions mécaniques}{N05}{4 h}{Cette leçon introduit l'énergie potentielle (de pesanteur et élastique) et l'énergie mécanique d'un système. Elle établit le principe de conservation de l'énergie mécanique et montre comment le travail des forces non conservatrices explique les situations de non conservation (freinage, chocs).
\vspace{4pt}
\begin{multicols}{2}\small
\begin{itemize}
\item À la fin de cette leçon, je sais définir l'énergie potentielle et citer des exemples d'énergies potentielles
\item À la fin de cette leçon, je sais exprimer et calculer l'énergie potentielle de pesanteur d'un système au voisinage de la Terre
\item À la fin de cette leçon, je sais exprimer et calculer l'énergie potentielle élastique d'un ressort à réponse linéaire
\item À la fin de cette leçon, je sais exprimer et calculer l'énergie mécanique d'un système
\item À la fin de cette leçon, je sais énoncer le principe de conservation de l'énergie mécanique et identifier les situations de conservation
\item À la fin de cette leçon, je sais exploiter la relation entre les variations d'énergie cinétique et d'énergie potentielle
\end{itemize}\end{multicols}}

\begin{sommaire}
\begin{multicols}{2}
\begin{enumerate}
\item Notion d'énergie potentielle
\item Énergie potentielle de pesanteur
\item Énergie potentielle élastique
\item Énergie mécanique et sa conservation
\item Non conservation de l'énergie mécanique
\item Bilan et méthodes de résolution
\item Activité d'intégration
\item Méthodes et exercices résolus
\item Exercices
\item Corrigés des exercices
\item Fiche bilan
\end{enumerate}
\end{multicols}
\end{sommaire}

\begin{objectifs}
\begin{itemize}
\item À la fin de cette leçon, je sais définir l'énergie potentielle et citer des exemples d'énergies potentielles
\item À la fin de cette leçon, je sais exprimer et calculer l'énergie potentielle de pesanteur d'un système au voisinage de la Terre
\item À la fin de cette leçon, je sais exprimer et calculer l'énergie potentielle élastique d'un ressort à réponse linéaire
\item À la fin de cette leçon, je sais exprimer et calculer l'énergie mécanique d'un système
\item À la fin de cette leçon, je sais énoncer le principe de conservation de l'énergie mécanique et identifier les situations de conservation
\item À la fin de cette leçon, je sais exploiter la relation entre les variations d'énergie cinétique et d'énergie potentielle
\item À la fin de cette leçon, je sais exprimer la variation de l'énergie mécanique à l'aide du travail des forces non conservatrices
\item À la fin de cette leçon, je sais analyser une situation de freinage ou de choc du point de vue énergétique
\end{itemize}
\end{objectifs}

\begin{prerequis}
\begin{itemize}
\item Énergie cinétique d'un solide en translation : Ec = ½mv²
\item Théorème de l'énergie cinétique
\item Travail d'une force constante : W(F) = F·AB·cos α
\item Travail du poids : W(P) = mg(z\_A − z\_B)
\item Force de rappel d'un ressort : T = k·x (loi de Hooke)
\end{itemize}
\end{prerequis}

\newpage

\coursec{Notion d'énergie potentielle}

\courssub{Introduction : l'énergie invisible du mouvement}

Au stade omnisports de Yaoundé, un sauteur à la perche camerounais s'élance sur la piste. Il court, plante sa perche en fibre de verre qui se plie dangereusement sous son poids, puis — comme propulsé par un ressort invisible — s'envole à plus de 4 mètres de hauteur avant de retomber mollement sur le tapis. D'où vient l'énergie qui le propulse si haut alors qu'il ne fait que courir ? Où est-elle stockée dans la perche pliée ?

Quelques kilomètres plus loin, sur la route de Douala, un chauffeur de taxi-brousse freine brusquement pour éviter un mouton. La voiture s'immobilise, mais les disques de frein brûlent la main qu'on y approche. L'énergie cinétique du véhicule a disparu : où est-elle passée ?

Cette leçon répond à ces questions en introduisant deux nouvelles formes d'énergie : \textbf{l'énergie potentielle} (cette section) et \textbf{l'énergie mécanique} (sections suivantes). Nous allons découvrir qu'un objet immobile peut posséder de l'énergie, prête à se libérer.

\courssub{Observation : l'énergie au repos}

Regardons trois situations familières au Cameroun :

\begin{popfigure}
\begin{tikzpicture}[scale=0.9, every node/.style={font=\small}]
    % Styles
    \tikzset{
        ground/.style={pattern=north east lines, pattern color=popmuted, draw=popdark},
        tree/.style={brown!70!black, line width=2pt},
        leaf/.style={popgreen!70!black, fill=popgreen!20},
        mango/.style={poporange, fill=poporange!30},
        spring/.style={popblue, thick, decorate, decoration={coil, aspect=0.4, segment length=3mm, amplitude=2.5mm}},
        hand/.style={fill=poppink!30, draw=popdark, rounded corners=2pt},
        water/.style={popblue!40, fill=popblue!10},
        dam/.style={popdark, fill=popgray!30, thick},
        arrow/.style={->, >=Stealth, thick, poporange},
        labelEp/.style={poppurple, font=\bfseries\small}
    }

    % --- SCENE 1 : Mangue sur l'arbre ---
    \begin{scope}[xshift=0cm]
        % Sol
        \draw[ground] (-1.5,-0.2) rectangle (1.5,0);
        % Arbre
        \draw[tree] (0,0) -- (0,2.5);
        \draw[tree] (-0.1,2.5) -- (-0.8,3.2) -- (-0.5,3.5);
        \draw[tree] (0.1,2.5) -- (0.7,3.1) -- (0.4,3.4);
        % Feuilles
        \foreach \x/\y in {-0.6,3.3/0.3, 0.5,3.2/0.3, -0.3,3.5/0.2, 0.2,3.4/0.2}
            \draw[leaf] (\x,\y) ellipse (0.3 and 0.15);
        % Mangue
        \draw[mango] (0.2,1.8) circle (0.25) node[above right=2pt] {\textbf{Mangue}};
        % Vecteur poids
        \draw[arrow] (0.2,1.55) -- (0.2,0.5) node[midway, right] {$\vec{P}=m\vec{g}$};
        % Cote hauteur
        \draw[<->, thin, popdark] (1.7,0) -- (1.7,1.8) node[midway, right] {$h$};
        % Label Energie
        \node[labelEp, align=center] at (0, -1.2) {Énergie potentielle\\de pesanteur $E_p$};
        \node[align=center, popdark] at (0, 3.8) {\textbf{(a) Mangue sur l'arbre}};
    \end{scope}

    % --- SCENE 2 : Ressort comprimé ---
    \begin{scope}[xshift=5.5cm]
        % Mur
        \draw[popdark, fill=popgray!20, thick] (-1.5, -0.5) rectangle (-1.2, 1.5);
        % Ressort au repos (gris clair)
        \draw[popmuted, decorate, decoration={coil, aspect=0.4, segment length=3mm, amplitude=2mm}] (-1.2, 0.5) -- (1.5, 0.5);
        % Ressort comprimé (bleu)
        \draw[spring] (-1.2, 0.5) -- (0.2, 0.5);
        % Bloc / Main qui pousse
        \draw[hand] (0.2, 0.1) rectangle (1.0, 0.9) node[midway] {\textbf{Main}};
        % Flèche force
        \draw[arrow] (1.2, 0.5) -- (2.0, 0.5) node[midway, above] {$\vec{F}$};
        % Cote compression
        \draw[<->, thin, popdark] (-1.2, -0.8) -- (0.2, -0.8) node[midway, below] {$x$};
        % Label Energie
        \node[labelEp, align=center] at (0, -1.2) {Énergie potentielle\\élastique $E_p$};
        \node[align=center, popdark] at (0, 1.8) {\textbf{(b) Ressort comprimé}};
    \end{scope}

    % --- SCENE 3 : Barrage de la Sanaga ---
    \begin{scope}[xshift=11.5cm]
        % Barrage
        \draw[dam] (-1.5, -0.5) -- (-1.5, 2.5) -- (-0.8, 2.5) -- (-0.8, -0.5) -- cycle;
        % Eau retenue (haut)
        \draw[water] (-1.5, -0.5) -- (-1.5, 2.0) -- (-0.8, 2.0) -- (-0.8, -0.5) -- cycle;
        % Niveau aval (bas)
        \draw[water] (0.5, -0.5) -- (0.5, 0.5) -- (1.5, 0.5) -- (1.5, -0.5) -- cycle;
        % Tuyau / Penstock
        \draw[popdark, thick] (-1.15, 0.5) -- (-1.15, -0.5) -- (0.5, -0.5) -- (0.5, 0.5);
        % Turbine
        \draw[popdark, fill=popgray!20, thick] (0.5, -0.5) rectangle (1.5, 0.5) node[midway, rotate=90] {\textbf{Turbine}};
        % Lignes HT
        \draw[popgold, thick, decorate, decoration={snake, amplitude=1pt, segment length=4pt}] (1.5, 0) -- (2.5, 1.5) node[above] {Électricité};
        \draw[popgold, thick, decorate, decoration={snake, amplitude=1pt, segment length=4pt}] (1.5, 0) -- (2.5, -1.5) node[below] {Vers Yaoundé};
        % Flèche écoulement
        \draw[arrow, popblue] (-1.15, 1.5) -- (-1.15, 0.0) node[midway, left] {Écoulement};
        \draw[arrow, popblue] (-0.5, -0.5) -- (0.5, -0.5) node[midway, below] {Vitesse};
        % Cote hauteur de chute
        \draw[<->, thin, popdark] (1.8, -0.5) -- (1.8, 2.0) node[midway, right] {$H$};
        % Label Energie
        \node[labelEp, align=center] at (0, -1.2) {Énergie potentielle\\$\to$ Cinétique $\to$ Électrique};
        \node[align=center, popdark] at (0, 2.8) {\textbf{(c) Barrage de la Sanaga (Song Loulou / Edéa)}};
    \end{scope}
\end{tikzpicture}
\legende{Trois systèmes au repos possédant de l'énergie potentielle : (a) une mangue à hauteur $h$ (pesanteur), (b) un ressort comprimé de $x$ (élastique), (c) l'eau retenue par le barrage de la Sanaga (pesanteur $\to$ électricité).}
\end{popfigure}

Dans les trois cas, le système est \textbf{immobile} (vitesse nulle, donc énergie cinétique $E_c = 0$). Pourtant, chacun est capable de produire du mouvement, de la chaleur, de l'électricité ou de la déformation. Cette capacité à engendrer une action future, stockée dans la \textbf{position} (hauteur) ou l'\textbf{état de déformation} (compression), est ce que nous appelons \textbf{énergie potentielle}.

\begin{definition}[Énergie potentielle d'un système]
L'\textbf{énergie potentielle} (notée $E_p$) d'un système est l'énergie qu'il possède du fait de sa \textbf{position} dans un champ de forces (champ de pesanteur, champ électrique) ou de son \textbf{état de déformation} (ressort étiré ou comprimé, gaz comprimé). Elle représente une énergie « en réserve », susceptible d'être convertie en d'autres formes d'énergie (cinétique, thermique, électrique, etc.).
\begin{itemize}
    \item Symbole : $E_p$ (ou $U$ dans certaines notations anglo-saxonnes).
    \item Unité SI : le \textbf{joule} (symbole \textbf{J}).
    \item C'est une grandeur \textbf{scalaire} (un nombre, pas un vecteur).
\end{itemize}
\end{definition}

\courssub{Définition générale et choix de la référence}

Une propriété fondamentale distingue l'énergie potentielle de l'énergie cinétique : \textbf{l'énergie potentielle est définie à une constante additive près}.

\emph{Pourquoi ?} Parce que seule la \textbf{variation} d'énergie potentielle $\Delta E_p = E_{p,f} - E_{p,i}$ entre un état initial et un état final a une signification physique (elle est liée au travail des forces conservatrices, comme nous le verrons). La valeur absolue $E_p$ en un point n'a pas de sens intrinsèque ; on peut ajouter la même constante à toutes les valeurs sans changer la physique.

\begin{propriete}[Choix de l'origine de l'énergie potentielle]
Pour lever l'indétermination, on \textbf{choisit arbitrairement une position de référence} (souvent notée $O$ ou $z=0$) où l'on \textbf{décrète} que $E_p = 0$.
\begin{itemize}
    \item Pour la pesanteur : on choisit souvent le niveau du sol, le fond d'un trou, ou le niveau de la mer.
    \item Pour un ressort : on choisit sa position de repos (non déformé).
\end{itemize}
Une fois cette référence fixée, l'énergie potentielle en tout point est \textbf{déterminée de façon unique}.
\end{propriete}

\begin{attention}[Erreur fréquente : « L'énergie potentielle est toujours positive »]
\textbf{FAUX.} L'énergie potentielle peut être \textbf{négative} si le système se trouve \textbf{sous} le niveau de référence choisi (ex: une pierre au fond d'un puits par rapport au niveau du sol). Seule la \textbf{variation} $\Delta E_p$ a un sens physique absolu. Au Probatoire, \textbf{précisez toujours votre choix de référence} (ex: « Je choisis $E_p=0$ au niveau du sol »).
\end{attention}

\begin{popfigure}
\begin{tikzpicture}[scale=1.1, >=Stealth]
    % Axe vertical z
    \draw[->, thick, popdark] (0, -3.5) -- (0, 4) node[above left] {$z$ (m)};
    \draw[thick, popdark] (-0.1, 0) -- (0.1, 0) node[right] {$O$};
    \node[below right] at (0,0) {$E_p=0$ (référence)};

    % Graduations
    \foreach \y/\label in {-3/-3, -2/-2, -1/-1, 1/1, 2/2, 3/3} {
        \draw[thin, popdark] (-0.1, \y) -- (0.1, \y);
        \node[left=2pt] at (0, \y) {\small $\label$};
    }

    % Point A (haut)
    \draw[popblue, very thick] (-0.8, 2.5) -- (0.8, 2.5);
    \node[popblue, right] at (0.8, 2.5) {$A\ (z_A > 0)$};
    \node[popblue, left] at (-0.8, 2.5) {$E_{p,A} > 0$};

    % Point B (bas)
    \draw[poporange, very thick] (-0.8, -2.0) -- (0.8, -2.0);
    \node[poporange, right] at (0.8, -2.0) {$B\ (z_B < 0)$};
    \node[poporange, left] at (-0.8, -2.0) {$E_{p,B} < 0$};

    % Vecteur g
    \draw[->, very thick, popred] (1.5, 1) -- (1.5, -1) node[midway, right] {$\vec{g}$};
    \node[popred, right] at (1.5, 1.3) {Champ de pesanteur};

    % Flèche variation
    \draw[<->, thick, poppurple, dashed] (1.0, -2.0) -- (1.0, 2.5) node[midway, right, poppurple] {$\Delta z = z_A - z_B > 0$};
    \node[poppurple, align=center] at (2.5, 0) {$\Delta E_p = E_{p,A} - E_{p,B} > 0$\\Énergie gagnée en montant};
\end{tikzpicture}
\legende{Choix de l'origine des altitudes pour l'énergie potentielle de pesanteur. L'origine $O$ ($z=0$) est choisie arbitrairement. Au-dessus, $E_p > 0$ ; en dessous, $E_p < 0$. Seule la différence $\Delta E_p$ est physiquement significative.}
\end{popfigure}

\courssub{Principales formes d'énergie potentielle}

Le programme de Première C exige la maîtrise de deux formes principales. D'autres existent et seront rencontrées en Terminale ou à l'université.

\begin{aretenir}[Formes d'énergie potentielle au programme de 1ère C]
\begin{enumerate}
    \item \textbf{Énergie potentielle de pesanteur ($E_{pp}$)} : Energie due à la position dans le champ de pesanteur terrestre (hauteur). \emph{Étudiée en détail Section 2.}
    \item \textbf{Énergie potentielle élastique ($E_{pe}$)} : Energie stockée dans un ressort (ou tout solide élastique à réponse linéaire) du fait de sa déformation (allongement ou compression). \emph{Étudiée en détail Section 3.}
\end{enumerate}
\end{aretenir}

\begin{savaistu}[Le barrage de Song Loulou : l'énergie de la Sanaga au service du Cameroun]
La rivière Sanaga, plus long fleuve entièrement camerounais, alimente les barrages hydroélectriques d'\textbf{Edéa} (mis en service 1956), de \textbf{Song Loulou} (1984) et de \textbf{Lom Pangar} (2018, régulation). L'eau stockée en amont (lac de retenue) possède une immense \textbf{énergie potentielle de pesanteur} ($E_p = mgh$). En tombant dans les conduites forcées (penstocks), cette énergie se convertit en énergie cinétique, qui fait tourner les turbines couplées à des alternateurs produisant l'électricité distribuée par \textbf{ENEO} (ex-SONEL) vers Yaoundé, Douala et tout le réseau interconnecté Sud (RIS). C'est une application géante de la conversion $E_p \to E_c \to E_{\text{électrique}}$ !
\end{savaistu}

\emph{Autres formes (culture générale, non exigées au Probatoire 1ère C) :}
\begin{itemize}
    \item \textbf{Énergie potentielle électrique / électrostatique} : Energie d'une charge dans un champ électrique (condensateur, orage).
    \item \textbf{Énergie potentielle chimique} : Energie stockée dans les liaisons chimiques (carburant, batterie, nourriture, bois). C'est la source de l'énergie thermique des freins du taxi-brousse (via combustion essence $\to$ chaleur $\to$ freinage).
    \item \textbf{Énergie potentielle nucléaire} : Energie de liaison des nucléons (centrales nucléaires, soleil, radiothérapie à l'Hôpital Général de Yaoundé).
\end{itemize}

\begin{exemplebox}[La mangue qui tombe : conversion $E_p \to E_c$]
Une mangue de masse $m = 200\,\text{g} = 0,200\,\text{kg}$ est cueillie à une hauteur $h = 5,0\,\text{m}$ au-dessus du sol. On choisit le niveau de référence $E_p = 0$ au niveau du sol ($z=0$).
\begin{itemize}
    \item \textbf{État initial (sur l'arbre) :} $v_i = 0 \Rightarrow E_{c,i} = 0$. $z_i = 5,0\,\text{m}$. L'énergie potentielle de pesanteur vaut $E_{p,i} = m g z_i$ (formule à venir). Elle est \textbf{maximale}.
    \item \textbf{État final (au sol, juste avant l'impact) :} $z_f = 0 \Rightarrow E_{p,f} = 0$. La mangue a acquis une vitesse $v_f$. Son énergie cinétique $E_{c,f} = \frac{1}{2}m v_f^2$ est \textbf{maximale}.
    \item \textbf{Bilan :} L'énergie potentielle de pesanteur s'est transformée en énergie cinétique. $E_{p,i} \approx E_{c,f}$ (si on néglige les frottements de l'air).
\end{itemize}
C'est exactement le principe du sauteur à la perche : son énergie cinétique de course $\to$ énergie potentielle élastique de la perche pliée $\to$ énergie potentielle de pesanteur au sommet du saut $\to$ énergie cinétique de la retombée.
\end{exemplebox}

\begin{attention}[Piège classique : « Objet immobile = Énergie nulle »]
\textbf{Ne confondez jamais « énergie cinétique nulle » et « énergie totale nulle ».}
\begin{itemize}
    \item Un livre posé sur une étagère a $E_c = 0$ mais $E_p > 0$ (par rapport au sol).
    \item Une pile neuve dans une lampe éteinte a $E_c = 0$ mais possède de l'énergie potentielle chimique.
    \item Un ressort comprimé maintenu par un doigt a $E_c = 0$ mais $E_{pe} > 0$.
\end{itemize}
\textbf{Astuce :} Dès qu'un système est susceptible d'évoluer spontanément (tomber, se détendre, réagir), il possède de l'énergie potentielle.
\end{attention}

\courssub{Exemple résolu : Choix de la référence et variation d'énergie potentielle}

\begin{exoresolu}[Variation d'énergie potentielle d'un bloc glissant sur un plan incliné]
\textbf{Énoncé :} Un bloc de masse $m = 2,0\,\text{kg}$ est lâché sans vitesse initiale au point $A$ d'un plan incliné sans frottement. Il glisse jusqu'au point $B$ situé $1,5\,\text{m}$ plus bas verticalement. On note $g = 9,81\,\text{N}\cdot\text{kg}^{-1}$.
\begin{enumerate}
    \item Choisir un référentiel et une origine pour l'énergie potentielle de pesanteur. Exprimer $E_{p,A}$ et $E_{p,B}$.
    \item Calculer la variation $\Delta E_p = E_{p,B} - E_{p,A}$.
    \item Interpréter le signe de $\Delta E_p$.
\end{enumerate}

\tcblower

\textbf{Solution détaillée :}

\textbf{1. Choix du référentiel et de l'origine.}
\begin{itemize}
    \item On choisit un repère terrestre (galiléen au voisinage de la Terre).
    \item On choisit un axe vertical $z$ orienté \textbf{vers le haut}.
    \item \textbf{Origine de l'énergie potentielle ($E_p=0$) :} On la place au niveau du point $B$ (le point le plus bas atteint). C'est un choix commode car $E_{p,B}=0$.
\end{itemize}
Dans ce repère :
\begin{itemize}
    \item $z_B = 0\,\text{m} \implies E_{p,B} = 0\,\text{J}$.
    \item $z_A = +1,5\,\text{m}$ (car $A$ est $1,5\,\text{m}$ \emph{au-dessus} de $B$).
    \item $E_{p,A} = m g z_A = 2,0 \times 9,81 \times 1,5 = \mathbf{29,43\,J}$.
\end{itemize}

\textbf{2. Calcul de la variation $\Delta E_p$.}
\[
\Delta E_p = E_{p,B} - E_{p,A} = 0 - 29,43 = \mathbf{-29,43\,J} \approx \mathbf{-29,4\,J}
\]

\textbf{3. Interprétation.}
La variation est \textbf{négative}. L'énergie potentielle de pesanteur \textbf{diminue} de $29,4\,\text{J}$ lors de la descente de $A$ vers $B$. Cette énergie « perdue » sous forme potentielle réapparaît sous forme d'énergie cinétique (le bloc accélère). Nous verrons dans la section 4 que $\Delta E_c = -\Delta E_p = +29,4\,\text{J}$ (si pas de frottement).

\begin{methode}[Bien poser son problème d'énergie potentielle]
\begin{enumerate}
    \item \textbf{Dessiner} un schéma simple avec l'axe vertical $z$ orienté vers le haut.
    \item \textbf{Choisir et annoncer} explicitement l'origine où $E_p = 0$ (souvent le point le plus bas de la trajectoire ou le niveau du sol).
    \item \textbf{Repérer} les altitudes $z_i$ et $z_f$ des positions initiale et finale par rapport à cette origine (positives si au-dessus, négatives si en-dessous).
    \item \textbf{Appliquer} la formule $E_p = mgz$ (pesanteur) ou $E_p = \frac{1}{2}kx^2$ (élastique).
    \item \textbf{Calculer} $\Delta E_p = E_{p,f} - E_{p,i}$ et \textbf{commenter le signe}.
\end{enumerate}
\end{methode}
\end{exoresolu}

\courssub{Bilan de la section}

\begin{aretenir}[À retenir sur la notion d'énergie potentielle]
\begin{itemize}
    \item L'énergie potentielle $E_p$ (en \textbf{J}) est une énergie de \textbf{position} ou de \textbf{configuration}. Un objet immobile peut en posséder.
    \item Elle est définie \textbf{à une constante près} : on doit \textbf{choisir une référence} ($E_p=0$) pour calculer des valeurs numériques.
    \item Seules les \textbf{variations} $\Delta E_p$ ont une signification physique absolue.
    \item Deux formes sont au programme de 1ère C :
    \begin{itemize}
        \item \textbf{Pesanteur :} $E_{pp} = mgz$ (près de la Terre, $z$ altitude par rapport à la référence).
        \item \textbf{Élastique :} $E_{pe} = \frac{1}{2}k(\Delta l)^2$ ($k$ raideur, $\Delta l$ déformation par rapport au repos).
    \end{itemize}
    \item D'autres formes existent (chimique, électrique, nucléaire) : même concept, formules différentes.
\end{itemize}
\end{aretenir}

Nous possédons maintenant l'outil « énergie potentielle ». Dans la section suivante, nous en déduirons les expressions précises pour la pesanteur et l'élastique, puis nous construirons l'\textbf{énergie mécanique} $E_m = E_c + E_p$ et découvrirons son extraordinaire propriété de \textbf{conservation} (ou non) selon les forces en présence.

\coursec{Énergie potentielle de pesanteur}

Dans la section précédente, nous avons introduit la notion générale d'\textbf{énergie potentielle} : une énergie de \emph{position} (ou de configuration) qui permet de rendre compte du travail d'une force conservative sans avoir à recalculer une intégrale curviligne à chaque changement de trajectoire. Nous allons maintenant appliquer ce concept fondamental à la force la plus familière de notre environnement : le \textbf{poids}. Que ce soit un sac de ciment hissé sur un chantier à Douala, l'eau retenue par le barrage de \textbf{Lom Pangar} ou de \textbf{Nachtigal}, ou encore un ballon de football qui s'élève après un coup de pied puissant lors de la finale de la Coupe d'Afrique des Nations, c'est toujours la même physique qui est en jeu : la conversion mutuelle entre énergie cinétique et énergie potentielle de pesanteur.

\courssub{Système étudié et rappel sur le travail du poids}

Considérons un \textbf{système} réduit à un \textbf{solide de masse $m$} (son centre d'inertie $G$ est confondu avec le point d'application du poids), placé \textbf{au voisinage de la Terre}. Le champ de pesanteur $\vec{g}$ est supposé \textbf{uniforme} (direction verticale descendante, norme constante $g \approx 9,81\,\text{N/kg}$ ou $10\,\text{N/kg}$ pour les estimations rapides).

Choisissons un \textbf{repère terrestre} supposé galiléen, muni d'un axe vertical \textbf{orienté vers le haut} (noté $z$). L'origine $O$ de cet axe est choisie arbitrairement (niveau de la mer, sol de la cour, fond du barrage, etc.). L'altitude du centre d'inertie $G$ est la coordonnée $z_G$ sur cet axe.

\begin{popfigure}
\begin{tikzpicture}[scale=1.1, >=Stealth]
% Axes
\draw[->, thick] (-3,0) -- (3,0) node[right] {$x$};
\draw[->, thick] (0,-0.5) -- (0,4) node[above] {$z$ (vers le haut)};
% Points A et B
\coordinate (A) at (-1.5, 1);
\coordinate (B) at (1.5, 3);
% Chemin 1 : vertical puis horizontal (escalier)
\draw[popblue, thick, ->] (A) -- (-1.5, 3) -- (B);
% Chemin 2 : courbe parabolique
\draw[poporange, thick, ->] (A) .. controls (0, 2.5) .. (B);
% Points
\fill[popblue] (A) circle (2.5pt) node[left] {$A\,(z_A)$};
\fill[poporange] (B) circle (2.5pt) node[right] {$B\,(z_B)$};
% Vecteur poids en A
\draw[popred, very thick, ->] (A) -- ++(0, -0.8) node[below right] {$\vec{P}=m\vec{g}$};
% Cotes
\draw[dashed, thin] (A) -- (-3,1) node[left] {$z_A$};
\draw[dashed, thin] (B) -- (-3,3) node[left] {$z_B$};
\draw[<->, thin] (-2.5,1) -- (-2.5,3) node[midway, left] {$\Delta z = z_B - z_A$};
% Legende
\node[align=left, anchor=north west] at (-2.8, 3.5) {\small \textbf{Chemin bleu :} vertical puis horizontal\\ \small \textbf{Chemin orange :} courbe quelconque};
\end{tikzpicture}
\legende{Travail du poids entre $A$ et $B$ : indépendant du chemin suivi. Seule la différence d'altitude $\Delta z = z_B - z_A$ compte.}
\end{popfigure}

Rappelons le résultat fondamental établi en classe de Seconde et approfondi en Première : le \textbf{travail du poids} $\vec{P} = m\vec{g}$ entre deux points $A$ et $B$ ne dépend \textbf{que des altitudes} $z_A$ et $z_B$ des points initial et final, et \textbf{pas du tout de la trajectoire} suivie.
\[
W_{A \to B}(\vec{P}) = \int_A^B \vec{P} \cdot d\vec{l} = \int_A^B (-mg)\,\mathrm{d}z = mg(z_A - z_B).
\]
Cette indépendance vis-à-vis du chemin est la signature caractéristique d'une \textbf{force conservative}. Elle autorise la définition d'une énergie potentielle associée.

\courssub{Définition de l'énergie potentielle de pesanteur}

\begin{definition}[Énergie potentielle de pesanteur]
L'\textbf{énergie potentielle de pesanteur} d'un solide de masse $m$ dont le centre d'inertie est d'altitude $z$ (sur un axe vertical orienté vers le haut) est définie par :
\[
E_p(z) = mgz + C,
\]
où $C$ est une \textbf{constante arbitraire} (choisie une fois pour toutes pour le problème étudié).
\end{definition}

\begin{attention}[Choix de l'origine des altitudes — Piège classique]
La valeur \textbf{absolue} de $E_p$ n'a pas de sens physique ; seule sa \textbf{variation} $\Delta E_p$ est mesurable et pertinente. Cependant, pour éviter les erreurs de signe, \textbf{vous devez toujours, dès la première ligne de votre copie} :
\begin{enumerate}
\item Dessiner un petit repère avec l'axe $z$ \textbf{orienté vers le haut} ($\uparrow$).
\item Placer l'origine $O$ ($z=0$) à un endroit commode (souvent le point le plus bas du mouvement, ou le sol).
\item Noter explicitement : \og Axe $z$ orienté vers le haut, origine au niveau du sol \fg (ou au point bas, etc.).
\end{enumerate}
Ne jamais écrire $E_p = mgh$ sans avoir défini ce qu'est $h$ (souvent confondu avec une hauteur positive vers le bas).
\end{attention}

\begin{propriete}[Choix usuel de la constante — Énergie potentielle nulle à l'origine]
Par commodité, on choisit presque toujours la constante $C$ de telle sorte que \textbf{$E_p = 0$ pour $z = 0$}. Cela donne l'expression simplifiée, \textbf{à connaître par cœur} :
\[
\boxed{E_p(z) = mgz}
\]
avec $z$ l'altitude du centre d'inertie sur un axe vertical \textbf{orienté vers le haut}, $m$ en $\text{kg}$, $g$ en $\text{N/kg}$ (ou $\text{m/s}^2$), $E_p$ en \textbf{joules (J)}.
\end{propriete}

\begin{methode}[Protocole obligatoire avant tout calcul d'énergie potentielle]
\begin{enumerate}
\item \textbf{Schématiser} : tracer l'axe $z$ vertical, flèche vers le haut.
\item \textbf{Fixer l'origine} $O$ ($z=0$) : point le plus bas atteint, niveau du sol, surface de l'eau, etc.
\item \textbf{Repérer les altitudes} $z_A, z_B, \dots$ des positions successives du centre d'inertie.
\item \textbf{Écrire} $E_p = mgz$ pour chaque position.
\item \textbf{Calculer la variation} $\Delta E_p = E_{p,f} - E_{p,i} = mg(z_f - z_i)$.
\end{enumerate}
\end{methode}

\begin{exemplebox}[Sac de ciment hissé sur un chantier — Application numérique]
Un maçon hisse un sac de ciment de masse $m = 50\,\text{kg}$ depuis le sol ($z=0$) jusqu'à un échafaudage situé à $z = 3,0\,\text{m}$ au-dessus du sol. On prend $g = 10\,\text{N/kg}$.
\begin{enumerate}
\item \textbf{Choix de l'axe} : $z$ vertical vers le haut, origine $O$ au niveau du sol $\Rightarrow z_{\text{sol}} = 0$, $z_{\text{échaf}} = 3,0\,\text{m}$.
\item \textbf{Énergie potentielle initiale} (au sol) : $E_{p,i} = mg \times 0 = 0\,\text{J}$.
\item \textbf{Énergie potentielle finale} (sur l'échafaudage) :
\[
E_{p,f} = mgz = 50 \times 10 \times 3,0 = 1\,500\,\text{J}.
\]
\item \textbf{Variation} : $\Delta E_p = E_{p,f} - E_{p,i} = 1\,500\,\text{J}$.
\end{enumerate}
Le maçon a fourni un travail d'au moins $1\,500\,\text{J}$ (en réalité plus à cause des frottements et de son rendement musculaire) pour augmenter l'énergie potentielle du sac. Cette énergie est \textbf{stockée} : si le sac tombe, elle se convertira en énergie cinétique.
\end{exemplebox}

\courssub{Représentation graphique : $E_p$ en fonction de $z$}

Puisque $E_p(z) = mgz$ (avec $C=0$), c'est une \textbf{fonction affine (droite) passant par l'origine}, de pente $mg$.

\begin{popfigure}
\begin{tikzpicture}
\begin{axis}[
    width=\linewidth, height=7cm,
    axis lines=middle,
    xlabel={$z$ (m)}, ylabel={$E_p$ (J)},
    xmin=-0.5, xmax=5,
    ymin=-200, ymax=2200,
    xtick={0,1,2,3,4}, ytick={0,500,1000,1500,2000},
    domain=0:4, samples=2,
    tick label style={font=\small},
    label style={font=\small},
    grid=major, grid style={popline},
]
\addplot[popcurve, very thick] {50*10*x};
\node[popblue, anchor=south west, font=\small] at (axis cs:1,500) {$E_p = mgz$};
\addplot[only marks, mark=*, mark size=2.5pt, popred] coordinates {(3,1500)};
\node[popred, anchor=south west, font=\small] at (axis cs:3.2,1500) {$(3\,\text{m}, 1500\,\text{J})$};
\end{axis}
\end{tikzpicture}
\legende{Énergie potentielle de pesanteur $E_p = mgz$ pour $m=50\,\text{kg}$, $g=10\,\text{N/kg}$. Droite de pente $mg = 500\,\text{N}$. L'énergie potentielle est proportionnelle à l'altitude.}
\end{popfigure}

\begin{aretenir}[Analyse dimensionnelle — Vérification systématique]
\[
[E_p] = [m][g][z] = \text{kg} \cdot \text{N/kg} \cdot \text{m} = \text{N} \cdot \text{m} = \text{J}.
\]
La pente de la droite $E_p(z)$ est $mg$, homogène à une force (newton). Cela confirme que $\dfrac{\mathrm{d}E_p}{\mathrm{d}z} = mg = -P_z$ (composante du poids sur l'axe $z$ orienté vers le haut).
\end{aretenir}

\courssub{Relation fondamentale : travail du poids et variation d'énergie potentielle}

C'est la relation clé qui permet de remplacer le calcul d'intégrale par une simple différence de valeurs d'une fonction d'état.

\begin{propriete}[Travail du poids et variation d'énergie potentielle de pesanteur]
Pour tout déplacement d'un point $A$ (altitude $z_A$) à un point $B$ (altitude $z_B$) :
\[
\boxed{W_{A \to B}(\vec{P}) = -\Delta E_p = E_p(A) - E_p(B) = mg(z_A - z_B)}.
\]
\end{propriete}

\begin{experience}[Démonstration guidée — À savoir refaire seul(e) au tableau]
\textbf{Objectif :} Établir $W_{A \to B}(\vec{P}) = -\Delta E_p$ à partir de la définition du travail et de $E_p = mgz$.

\begin{enumerate}
\item \textbf{Écrire le travail du poids} par définition intégrale sur un chemin quelconque $A \to B$ :
      $W_{A \to B}(\vec{P}) = \int_A^B \vec{P} \cdot d\vec{l}$.
\item \textbf{Projeter sur l'axe $z$ (vers le haut)} : $\vec{P} = -mg\,\vec{u}_z$, $d\vec{l} = \mathrm{d}z\,\vec{u}_z + \dots$ (composantes horizontales).
      $\vec{P} \cdot d\vec{l} = (-mg\,\vec{u}_z) \cdot (\mathrm{d}z\,\vec{u}_z) = -mg\,\mathrm{d}z$.
\item \textbf{Intégrer} : $W_{A \to B}(\vec{P}) = \int_{z_A}^{z_B} (-mg)\,\mathrm{d}z = -mg \int_{z_A}^{z_B} \mathrm{d}z = -mg(z_B - z_A) = mg(z_A - z_B)$.
\item \textbf{Calculer $\Delta E_p$} avec $E_p(z) = mgz$ (choix $C=0$) :
      $\Delta E_p = E_p(B) - E_p(A) = mgz_B - mgz_A = mg(z_B - z_A)$.
\item \textbf{Comparer} : $W_{A \to B}(\vec{P}) = -\Delta E_p$. \qed
\end{enumerate}
\end{experience}

\courssub{Interprétation physique : moteur ou résistif ?}

Le signe du travail du poids (et donc le signe de $-\Delta E_p$) dicte le rôle énergétique du poids :

\begin{center}
\begin{tabularx}{\linewidth}{|l|X|X|}\hline
\textbf{Situation} & \textbf{Mouvement} & \textbf{Conséquence énergétique} \\ \hline
\textbf{Descente} & $z_B < z_A \Rightarrow \Delta z < 0$ & $\Delta E_p = mg\Delta z < 0$ \quad (Énergie potentielle \textbf{diminue}) \\
& & $W(\vec{P}) = -\Delta E_p > 0$ \quad (Travail \textbf{moteur}) \\
& & Le poids \textbf{donne} de l'énergie au système $\to$ $E_c$ augmente. \\ \hline
\textbf{Montée} & $z_B > z_A \Rightarrow \Delta z > 0$ & $\Delta E_p = mg\Delta z > 0$ \quad (Énergie potentielle \textbf{augmente}) \\
& & $W(\vec{P}) = -\Delta E_p < 0$ \quad (Travail \textbf{résistif}) \\
& & Le poids \textbf{retire} de l'énergie au système $\to$ $E_c$ diminue. \\ \hline
\textbf{Horizontal} & $z_B = z_A \Rightarrow \Delta z = 0$ & $\Delta E_p = 0$ \quad $W(\vec{P}) = 0$ \\
& & Le poids ne travaille pas (force $\perp$ déplacement). \\ \hline
\end{tabularx}
\end{center}

\begin{attention}[Erreur fréquente au Probatoire — Signe de $\Delta E_p$]
\begin{itemize}
\item \textbf{Énoncé} : \og Un bloc glisse sans frottement sur un plan incliné. Calculer la variation de son énergie potentielle. \fg
\item \textbf{Erreur} : L'élève écrit $\Delta E_p = mg\Delta h$ avec $\Delta h > 0$ (la \og hauteur \fg descendue) et trouve $\Delta E_p > 0$.
\item \textbf{Correction} : $\Delta E_p = E_{p,f} - E_{p,i}$. Si le bloc \textbf{descend}, $z_f < z_i$, donc $\Delta E_p < 0$. L'énergie potentielle \textbf{diminue} toujours lors d'une descente.
\item \textbf{Astuce} : Toujours calculer $E_{p,f}$ et $E_{p,i}$ séparément avec $E_p = mgz$ ($z$ altitude sur axe $\uparrow$), puis faire la soustraction \textbf{finale moins initiale}. Ne jamais deviner le signe.
\end{itemize}
\end{attention}

\courssub{Exemple complet : Lancer vertical et diagramme énergétique}

\begin{exoresolu}[Balle de tennis lancée verticalement — Conversion $E_c \leftrightarrow E_p$]
Une balle de tennis de masse $m = 57\,\text{g} = 0,057\,\text{kg}$ est lancée verticalement vers le haut depuis le niveau du sol ($z=0$) avec une vitesse initiale $v_0 = 20\,\text{m/s}$. On néglige les frottements de l'air. $g = 10\,\text{N/kg}$.
\tcblower
\textbf{1. Choix du repère et énergies initiales.}
Axe $z$ vertical vers le haut, origine au point de lancer ($z=0$).
\begin{itemize}
\item $t=0$ : $z_0 = 0 \Rightarrow E_{p,0} = 0\,\text{J}$.
\item $v_0 = 20\,\text{m/s} \Rightarrow E_{c,0} = \frac{1}{2}mv_0^2 = 0,5 \times 0,057 \times 400 = 11,4\,\text{J}$.
\item Énergie mécanique initiale : $E_{m,0} = E_{c,0} + E_{p,0} = 11,4\,\text{J}$.
\end{itemize}

\textbf{2. Hauteur maximale $H$ atteinte.}
Au point haut, $v_H = 0 \Rightarrow E_{c,H} = 0$. Par conservation de l'énergie mécanique (section 4) : $E_{m,H} = E_{m,0}$.
\[
E_{p,H} = E_{m,0} = 11,4\,\text{J} \quad \Rightarrow \quad mgH = 11,4 \quad \Rightarrow \quad H = \frac{11,4}{0,057 \times 10} = 20\,\text{m}.
\]

\textbf{3. Vitesse au passage à $z = 10\,\text{m}$ (montée et descente).}
À $z = 10\,\text{m}$ : $E_p = mgz = 0,057 \times 10 \times 10 = 5,7\,\text{J}$.
Conservation : $E_c + E_p = 11,4 \Rightarrow E_c = 11,4 - 5,7 = 5,7\,\text{J}$.
\[
\frac{1}{2}mv^2 = 5,7 \quad \Rightarrow \quad v^2 = \frac{2 \times 5,7}{0,057} = 200 \quad \Rightarrow \quad v = \sqrt{200} \approx 14,1\,\text{m/s}.
\]
La vitesse a la même norme en montée et en descente à la même altitude (symétrie temporelle).
\end{exoresolu}

\begin{popfigure}
\begin{tikzpicture}[scale=0.9]
% t=0
\begin{scope}[shift={(0,0)}]
\draw[popblue, fill=popblue!30] (0,0) rectangle (1.5, 2.28) node[midway] {$E_c = 11,4\,\text{J}$};
\draw[poporange, fill=poporange!30] (1.5,0) rectangle (3, 0) node[midway] {$E_p = 0$};
\node[below] at (1.5, -0.3) {$t=0$ ($z=0$)};
\draw[<->] (0,2.4) -- (3,2.4) node[midway, above] {$E_m = 11,4\,\text{J}$};
\end{scope}
% z=10m
\begin{scope}[shift={(4,0)}]
\draw[popblue, fill=popblue!30] (0,0) rectangle (1.5, 1.14) node[midway] {$E_c = 5,7\,\text{J}$};
\draw[poporange, fill=poporange!30] (1.5,0) rectangle (3, 1.14) node[midway] {$E_p = 5,7\,\text{J}$};
\node[below] at (1.5, -0.3) {$z=10\,\text{m}$};
\draw[<->] (0,2.4) -- (3,2.4) node[midway, above] {$E_m = 11,4\,\text{J}$};
\end{scope}
% z=20m (sommet)
\begin{scope}[shift={(8,0)}]
\draw[popblue, fill=popblue!30] (0,0) rectangle (1.5, 0) node[midway] {$E_c = 0$};
\draw[poporange, fill=poporange!30] (1.5,0) rectangle (3, 2.28) node[midway] {$E_p = 11,4\,\text{J}$};
\node[below] at (1.5, -0.3) {$z=H=20\,\text{m}$};
\draw[<->] (0,2.4) -- (3,2.4) node[midway, above] {$E_m = 11,4\,\text{J}$};
\end{scope}
\end{tikzpicture}
\legende{Diagramme énergétique : l'énergie mécanique $E_m$ reste constante ; $E_c$ et $E_p$ s'échangent. Bleu = $E_c$, Orange = $E_p$.}
\end{popfigure}

\courssub{Application concrète : L'hydroélectricité au Cameroun}

\begin{savaistu}[Du barrage de Lom Pangar à la prise de courant — L'énergie potentielle à l'échelle nationale]
Les barrages de \textbf{Lom Pangar} (capacité $\approx 6\,\text{milliards m}^3$), \textbf{Nachtigal} (au fil de l'eau, $420\,\text{MW}$) et \textbf{Edéa} stockent de l'eau en altitude. L'eau du réservoir possède une \textbf{énergie potentielle de pesanteur} colossale : $E_p = m_{\text{eau}} g \bar{z}$, où $\bar{z}$ est l'altitude moyenne du centre de masse de l'eau par rapport aux turbines.
\begin{itemize}
\item Lorsque les vannes s'ouvrent, l'eau \textbf{descend} dans les conduites forcées : $\Delta E_p < 0$, le poids de l'eau fait un travail \textbf{moteur} ($W>0$).
\item Ce travail est converti en énergie cinétique de l'eau (jet rapide) puis en énergie mécanique de rotation des \textbf{turbines} (alternateurs).
\item Les alternateurs transforment cette énergie mécanique en \textbf{énergie électrique} transportée par le réseau \textbf{SONATREL/ENEO} vers Yaoundé, Douala, Bafoussam...
\end{itemize}
C'est une application industrielle directe de la relation $W(\vec{P}) = -\Delta E_p$. La \textbf{puissance} disponible dépend du \textbf{débit massique} $\dot{m}$ et de la \textbf{chute de hauteur} $\Delta z$ : $\mathcal{P} \approx \dot{m} g \Delta z$ (en négligeant les pertes). C'est pourquoi les barrages sont construits là où la topographie offre de grandes dénivelées.
\end{savaistu}

\begin{exercice}[Entraînement — Niveau Probatoire]
\textbf{Exercice 1.} Un mototaxi (masse totale $m = 200\,\text{kg}$ avec passager) gravit une côte rectiligne de longueur $L = 500\,\text{m}$ inclinée à $\alpha = 10^\circ$ par rapport à l'horizontale. On prend $g = 10\,\text{N/kg}$.
\begin{enumerate}
\item Choisir un repère adapté et exprimer l'énergie potentielle de pesanteur en fonction de l'abscisse $x$ le long de la pente (origine au départ).
\item Calculer la variation d'énergie potentielle $\Delta E_p$ entre le bas et le haut de la côte.
\item Quel est le travail du poids pendant cette montée ? Interpréter le signe.
\end{enumerate}

\textbf{Exercice 2.} Une pierre de masse $m = 200\,\text{g}$ est lâchée sans vitesse initiale du sommet d'un immeuble de hauteur $H = 45\,\text{m}$ (niveau du sol : $z=0$). On néglige l'air. $g=10\,\text{N/kg}$.
\begin{enumerate}
\item Écrire $E_p(z)$ et $E_c(z)$ en fonction de $z$ pendant la chute.
\item Vérifier que $E_m$ est constante.
\item Déterminer la vitesse d'impact au sol.
\end{enumerate}
\end{exercice}

\begin{corrige}[Exercice 1]
\begin{enumerate}
\item Axe $x$ le long de la pente, orienté vers le haut, origine au départ. Altitude $z = x\sin\alpha$.
$E_p(x) = mgz = mgx\sin\alpha = 200 \times 10 \times x \times \sin 10^\circ \approx 347,3\,x$ (en J, $x$ en m).
\item $\Delta E_p = E_p(L) - E_p(0) = mgL\sin\alpha = 2000 \times 500 \times 0,1736 \approx 1,74 \times 10^5\,\text{J}$.
\item $W(\vec{P}) = -\Delta E_p \approx -1,74 \times 10^5\,\text{J}$. Travail \textbf{résistif} (négatif) : le poids s'oppose à la montée, le moteur doit fournir au moins cette énergie (plus les frottements).
\end{enumerate}
\end{corrige}

\begin{corrige}[Exercice 2]
\begin{enumerate}
\item $E_p(z) = mgz = 0,2 \times 10 \times z = 2z$ (J).
$E_c(z) = E_m - E_p(z)$. Au départ $z=H=45$, $v=0 \Rightarrow E_c=0 \Rightarrow E_m = E_p(H) = 2 \times 45 = 90\,\text{J}$.
Donc $E_c(z) = 90 - 2z$.
\item $E_m = E_c(z) + E_p(z) = (90 - 2z) + 2z = 90\,\text{J}$ = constante. \qed
\item Au sol $z=0$ : $E_c = 90\,\text{J} = \frac{1}{2}mv^2 \Rightarrow v^2 = \frac{180}{0,2} = 900 \Rightarrow v = 30\,\text{m/s}$ (soit $108\,\text{km/h}$).
\end{enumerate}
\end{corrige}

\courssub{Bilan de la section}

\begin{aretenir}[Résumé : Énergie potentielle de pesanteur — Ce qu'il faut retenir pour le Probatoire]
\begin{itemize}
\item \textbf{Système} : Solide masse $m$ près de la Terre, champ $\vec{g}$ uniforme.
\item \textbf{Repère} : Axe $z$ \textbf{vertical orienté vers le haut}, origine choisie commodément ($z=0$).
\item \textbf{Expression} : $E_p(z) = mgz$ (choix $E_p=0$ pour $z=0$). Unité : \textbf{Joule (J)}.
\item \textbf{Relation fondamentale} : $W_{A\to B}(\vec{P}) = -\Delta E_p = mg(z_A - z_B)$.
\item \textbf{Interprétation} :
      \begin{itemize}
      \item Monte ($z \uparrow$) $\Rightarrow$ $\Delta E_p > 0$, $W(\vec{P}) < 0$ (résistif).
      \item Descend ($z \downarrow$) $\Rightarrow$ $\Delta E_p < 0$, $W(\vec{P}) > 0$ (moteur).
      \end{itemize}
\item \textbf{Méthode} : Toujours dessiner l'axe $z\uparrow$, noter $z_A, z_B$, calculer $E_{p,A}, E_{p,B}$, puis $\Delta E_p = E_{p,B} - E_{p,A}$.
\end{itemize}
\end{aretenir}

Nous possédons maintenant l'outil principal pour quantifier l'énergie de position dans le champ de pesanteur. Dans la section suivante, nous étudierons la deuxième grande forme d'énergie potentielle au programme : l'\textbf{énergie potentielle élastique} d'un ressort, avant de les réunir toutes les deux dans la définition de l'\textbf{énergie mécanique} et son principe de conservation.

\coursec{Énergie potentielle élastique}

Après avoir étudié l'énergie potentielle de pesanteur, qui dépend de la hauteur d'un objet dans le champ de pesanteur terrestre, nous nous intéressons maintenant à une autre forme d'énergie potentielle : celle qui est stockée dans un solide déformable, plus précisément dans un \cle{ressort à spires non jointives} (ou un élastique, une perche de saut, une lame de ressort). Contrairement à la force de pesanteur qui est constante au voisinage de la Terre, la force exercée par un ressort \emph{varie} avec la déformation. Cela nous obligera à calculer un travail non plus comme un simple produit $F \times d$, mais comme une aire sous une courbe.

\courssub{Système étudié et force de rappel}

Considérons un solide $S$ (une masse, un chariot, un projectile) accroché à un ressort horizontal de raideur $k$ (en $\si{\newton\per\meter}$), lui-même fixé à un mur. Le système est placé sur un plan horizontal parfait (sans frottement). Le ressort est dit \cle{linéaire} (ou de \cle{réponse linéaire}) : sa force de rappel $\vec{T}$ est proportionnelle à l'allongement algébrique $x$ par rapport à sa position d'équilibre (ressort non déformé).

\begin{popfigure}
\begin{tikzpicture}[scale=0.9, every node/.style={font=\small}]
    % Styles
    \tikzset{spring/.style={decorate, decoration={zigzag, pre length=0.1cm, post length=0.1cm, segment length=0.3cm, amplitude=0.3cm}, thick, popblue}}
    \tikzset{block/.style={draw, thick, minimum width=1.5cm, minimum height=1cm, fill=popblueL, rounded corners=2pt}}
    \tikzset{wall/.style={fill=popdark, pattern=north east lines, pattern color=popmuted}}
    
    % Repère commun
    \draw[->, thick] (-0.5, -3.5) -- (5.5, -3.5) node[right] {$x$};
    \draw[dashed, thin, popmuted] (2, -4) -- (2, 2.5) node[above, pos=0.1] {Équilibre ($x=0$)};
    
    % 1. Repos
    \draw[wall] (0, 1.5) rectangle (0.5, 2.5);
    \draw[spring] (0.5, 2) -- (2, 2);
    \draw[block] (2, 1.5) rectangle (3.5, 2.5) node[midway] {$S$};
    \node[above] at (1.25, 2.6) {État de repos : $x=0$};
    \node[below] at (1.25, 1.2) {$T=0$};
    
    % 2. Étiré
    \draw[wall] (0, -0.5) rectangle (0.5, 0.5);
    \draw[spring] (0.5, 0) -- (3, 0);
    \draw[block] (3, -0.5) rectangle (4.5, 0.5) node[midway] {$S$};
    \draw[<->, thick, poporange] (2, -0.7) -- (3, -0.7) node[midway, below] {$x>0$};
    \draw[->, thick, poporange] (4.5, 0) -- (5.2, 0) node[right] {$\vec{T}$ (vers la gauche)};
    \node[above] at (1.75, 0.6) {Étiré : $x>0$};
    
    % 3. Comprimé
    \draw[wall] (0, -2.5) rectangle (0.5, -1.5);
    \draw[spring] (0.5, -2) -- (1.2, -2);
    \draw[block] (1.2, -2.5) rectangle (2.7, -1.5) node[midway] {$S$};
    \draw[<->, thick, popgreen] (1.2, -2.7) -- (2, -2.7) node[midway, below] {$x<0$};
    \draw[->, thick, popgreen] (0.5, -2) -- (-0.2, -2) node[left] {$\vec{T}$ (vers la droite)};
    \node[above] at (0.85, -1.4) {Comprimé : $x<0$};
    
    % Force de rappel formula
    \node[align=center, fill=popcream, draw=poporange, rounded corners, inner sep=4pt] at (5, -2) {Loi de Hooke :\\ $\vec{T} = -k\,x\,\vec{u}_x$\\ $T = k|x|$};
\end{tikzpicture}
\legende{Un ressort à spires non jointives dans ses trois états : au repos ($x=0$), étiré ($x>0$) et comprimé ($x<0$). La force de rappel $\vec{T}$ s'oppose toujours à la déformation.}
\end{popfigure}

\begin{propriete}[Force de rappel d'un ressort linéaire (Loi de Hooke)]
Pour un ressort de raideur $k$ (en $\si{\newton\per\meter}$), la force de rappel $\vec{T}$ exercée sur le solide à la liaison est :
\[
\vec{T} = -k\,x\,\vec{u}_x
\]
où $x$ est l'allongement algébrique (positif si étirement, négatif si compression) et $\vec{u}_x$ le vecteur unitaire orienté dans le sens de l'étirement. La norme de la force est $T = k|x|$.
\end{propriete}

\begin{attention}[Rigidité vs Raideur]
Ne confonds pas la \cle{raideur} $k$ (en $\si{\newton\per\meter}$) qui caractérise le ressort, et la \cle{rigidité} qui est une propriété du matériau (module de Young). Un ressort souple a un $k$ faible ; un ressort raide a un $k$ grand. Vérifie toujours l'unité : $k$ s'exprime en newtons par mètre ($\si{\newton\per\meter}$ ou $\si{N/m}$).
\end{attention}

\courssub{Travail de la force de rappel et énergie potentielle élastique}

La force de rappel n'étant pas constante (elle vaut $0$ à l'équilibre et augmente linéairement avec $|x|$), on ne peut pas utiliser la formule $W = F \cdot d \cdot \cos\alpha$. On doit calculer le travail par l'intégrale (ou l'aire sous la courbe) de la force en fonction de la position.

\begin{experience}[Travail de la force de rappel — Méthode graphique]
\textbf{Matériel :} Dynamomètre, ressort de raideur connue $k$, règle graduée, masse $m$ sur rail à air (ou plan horizontal sans frottement).
\textbf{Protocole :}
\begin{enumerate}
    \item Placer le solide à l'équilibre ($x=0$).
    \item L'étirer lentement (quasi-statiquement) jusqu'à une position $x_f > 0$ en mesurant la force $F_{\text{main}}$ appliquée pour maintenir l'équilibre à chaque instant. Comme le mouvement est lent, $F_{\text{main}} = -\vec{T} = kx$.
    \item Tracer la courbe $F_{\text{main}} = f(x)$ (ou $T = f(x)$).
\end{enumerate}
\textbf{Observation :} La courbe $T = kx$ est une droite passant par l'origine de pente $k$. Le travail de la force extérieure $W(F_{\text{main}})$ pour amener le solide de $0$ à $x_f$ est l'aire du triangle rectangle de base $x_f$ et de hauteur $k x_f$.
\textbf{Interprétation :} Ce travail a été stocké sous forme d'énergie potentielle élastique dans le ressort. Par définition, la variation d'énergie potentielle est l'opposé du travail de la force intérieure (la force de rappel) : $\Delta E_p = -W(T)$. Comme $W(T) = -W(F_{\text{main}})$, on trouve l'expression de l'énergie potentielle élastique.
\end{experience}

\begin{popfigure}
\begin{tikzpicture}
    \begin{axis}[
        width=\linewidth, height=8cm,
        axis lines=middle,
        xlabel={$x$ (m)}, ylabel={$T$ (N)},
        xmin=-0.05, xmax=0.55,
        ymin=-0.5, ymax=5.5,
        xtick={0,0.2,0.4}, ytick={0,2,4},
        domain=0:0.5, samples=2,
        clip=false,
        tick label style={font=\small},
        label style={font=\small},
        enlargelimits=0.05
    ]
    % Courbe T = kx (k=10 N/m)
    \addplot[popcurve, thick, popblue, domain=0:0.5] {10*x};
    \node[popblue, right, font=\small] at (axis cs:0.5, 5) {$T = kx$};
    
    % Aire hachurée pour x_f = 0.4
    \addplot[poporangeL, fill=poporange!30, draw=poporange, domain=0:0.4] {10*x} \closedcycle;
    \draw[dashed, thin, popmuted] (axis cs:0.4,0) -- (axis cs:0.4,4);
    \draw[dashed, thin, popmuted] (axis cs:0,4) -- (axis cs:0.4,4);
    \node[below, font=\small] at (axis cs:0.4,0) {$x_f$};
    \node[left, font=\small] at (axis cs:0,4) {$k x_f$};
    
    % Flèches travail
    \draw[->, thick, poporange] (axis cs:0.2, 1.5) -- (axis cs:0.2, 2.5) node[above, poporange, font=\small] {$W(F_{\text{ext}}) = \frac{1}{2}k x_f^2$};
    \draw[->, thick, poppurple] (axis cs:0.2, 0.5) -- (axis cs:0.2, -0.5) node[below, poppurple, font=\small] {$W(T) = -\frac{1}{2}k x_f^2$};
    \end{axis}
\end{tikzpicture}
\legende{Travail de la force de rappel pour un étirement de $0$ à $x_f$. L'aire du triangle hachuré représente le travail de la force extérieure $W(F_{\text{ext}}) = \frac{1}{2}k x_f^2$. Le travail de la force de rappel est l'opposé : $W(T) = -\frac{1}{2}k x_f^2$.}
\end{popfigure}

\courssub{Définition de l'énergie potentielle élastique}

Choisissons l'origine de l'énergie potentielle élastique pour le ressort non déformé ($x=0$) : $E_p(0) = 0$. Pour un allongement algébrique $x$ (positif ou négatif), le travail de la force de rappel pour amener le solide de $0$ à $x$ est $W_{0 \to x}(T) = -\frac{1}{2}k x^2$ (l'aire du triangle est toujours positive, le travail de $\vec{T}$ est négatif car $\vec{T}$ s'oppose au déplacement si on étire, et dans le même sens si on lâche depuis l'étirement, mais la définition $\Delta E_p = -W(T)$ donne un résultat unique).

\begin{definition}[Énergie potentielle élastique]
L'énergie potentielle élastique $E_p$ d'un système constitué d'un ressort linéaire de raideur $k$ (en $\si{\newton\per\meter}$) et d'un solide, pour un allongement algébrique $x$ (en $\si{\meter}$) par rapport à la position d'équilibre, est :
\[
\boxed{E_p(x) = \frac{1}{2}\,k\,x^2}
\]
Elle s'exprime en joules ($\si{\joule}$). Elle est \emph{toujours positive ou nulle}, que le ressort soit étiré ($x>0$) ou comprimé ($x<0$), car elle dépend de $x^2$.
\end{definition}

\begin{propriete}[Relation travail — variation d'énergie potentielle élastique]
Entre deux positions $x_A$ et $x_B$, le travail de la force de rappel $\vec{T}$ est égal à l'opposé de la variation de l'énergie potentielle élastique :
\[
\boxed{W_{A \to B}(T) = -\Delta E_p = -\left( \frac{1}{2}k x_B^2 - \frac{1}{2}k x_A^2 \right) = \frac{1}{2}k x_A^2 - \frac{1}{2}k x_B^2}
\]
Cette relation est valable quel que soit le chemin suivi (la force de rappel est \cle{conservatrice}).
\end{propriete}

\begin{aretenir}[Analyse dimensionnelle]
Vérifions l'homogénéité de $E_p = \frac{1}{2}k x^2$ :
$[k] = \si{\newton\per\meter} = \si{\kilogram\per\second\squared}$.
$[x^2] = \si{\meter\squared}$.
$[k x^2] = \si{\kilogram\meter\squared\per\second\squared} = \si{\joule}$.
L'expression est homogène. Le facteur $\frac{1}{2}$ est sans dimension.
\end{aretenir}

\courssub{Évolution de l'énergie potentielle élastique}

Contrairement à l'énergie potentielle de pesanteur ($E_p = mgz$) qui est une fonction affine (droite) de l'altitude, l'énergie potentielle élastique est une fonction \cle{quadratique} (parabole) de l'allongement.

\begin{popfigure}
\begin{tikzpicture}
    \begin{axis}[
        width=\linewidth, height=8cm,
        axis lines=middle,
        xlabel={$x$ (m)}, ylabel={$E_p$ (J)},
        xmin=-0.55, xmax=0.55,
        ymin=-0.2, ymax=1.4,
        xtick={-0.4,-0.2,0,0.2,0.4}, ytick={0,0.4,0.8,1.2},
        domain=-0.5:0.5, samples=100,
        tick label style={font=\small},
        label style={font=\small},
        enlargelimits=0.05
    ]
    % Courbe Ep = 0.5 * k * x^2 (k=10 N/m)
    \addplot[popcurve, thick, poppurple] {0.5*10*x^2};
    \node[poppurple, above right, font=\small] at (axis cs:0.4, 0.8) {$E_p = \frac{1}{2}k x^2$};
    
    % Points remarquables
    \addplot[only marks, mark=*, mark size=2pt, poppurple] coordinates {(0,0) (-0.4, 0.8) (0.4, 0.8)};
    \node[below left, font=\small] at (axis cs:0,0) {$O$};
    \node[above left, font=\small] at (axis cs:-0.4, 0.8) {$x_A$};
    \node[above right, font=\small] at (axis cs:0.4, 0.8) {$x_B$};
    
    % Symétrie
    \draw[dashed, thin, popmuted] (axis cs:0,0) -- (axis cs:0,1.2);
    \node[above, font=\tiny, popmuted] at (axis cs:0, 1.25) {Axe de symétrie $x=0$};
    \end{axis}
\end{tikzpicture}
\legende{Courbe de l'énergie potentielle élastique $E_p = \frac{1}{2}k x^2$ (parabole). Elle est symétrique par rapport à l'axe $x=0$ (position d'équilibre). Minimum en $x=0$ ($E_p=0$).}
\end{popfigure}

\courssub{Exemples chiffrés et applications}

\begin{exemplebox}[Énergie stockée dans une perche de saut à la perche]
Un perchiste de masse $m = \SI{80}{kg}$ utilise une perche en fibre de carbone qui se comporte comme un ressort de raideur $k = \SI{1200}{N/m}$. Au moment de l'impulsion, la perche fléchit de $x = \SI{1,50}{m}$ par rapport à sa position droite.
Calculons l'énergie potentielle élastique stockée dans la perche à cet instant.
\[
E_p = \frac{1}{2} k x^2 = \frac{1}{2} \times \num{1200} \times (\num{1,50})^2 = \num{600} \times \num{2,25} = \SI{1350}{J}
\]
Cette énergie (provenant de l'énergie cinétique du perchiste) sera restituée pour le propulser vers le haut. Si toute cette énergie se convertit en énergie potentielle de pesanteur au sommet du saut ($\Delta E_p^{\text{pes}} = mg\Delta h$), on peut estimer le gain de hauteur $\Delta h$ :
\[
\Delta h = \frac{E_p}{mg} = \frac{\num{1350}}{\num{80} \times \num{9,81}} \approx \SI{1,72}{m}
\]
C'est une estimation théorique maximale ; dans la réalité, des pertes par frottements internes de la perche et frottements de l'air réduisent ce gain.
\end{exemplebox}

\begin{exemplebox}[Ressort de suspension d'un benskin (moto-taxi)]
À Yaoundé ou Douala, les \cle{benskines} (motos-taxis) affrontent des routes dégradées. La fourche avant contient un ressort (souvent associé à un amortisseur hydraulique). Supposons un ressort de raideur $k = \SI{25000}{N/m}$ (valeur typique pour une moto chargée). Lorsque la roue passe dans un trou, la fourche se comprime de $x = -\SI{5}{cm} = \SI{-0,05}{m}$.
L'énergie élastique stockée temporairement est :
\[
E_p = \frac{1}{2} \times \num{25000} \times (\num{0,05})^2 = \num{12500} \times \num{0,0025} = \SI{31,25}{J}
\]
Cette énergie est restituée pour repousser la roue vers le bas et maintenir le contact au sol. L'amortisseur (force non conservative) dissipe une partie de cette énergie en chaleur pour éviter les oscillations parasites (\guillemotleft pompage \guillemotright).
\end{exemplebox}

\begin{methode}[Calculer une énergie potentielle élastique]
\begin{enumerate}
    \item \textbf{Identifier} le système : ressort + solide. Noter la raideur $k$.
    \item \textbf{Repérer} la position d'équilibre ($x=0$) et mesurer l'allongement algébrique $x$ (en \textbf{mètres} !).
    \item \textbf{Convertir} les unités : $k$ en $\si{\newton\per\meter}$, $x$ en $\si{\meter}$.
    \item \textbf{Appliquer} la formule $E_p = \frac{1}{2} k x^2$.
    \item \textbf{Vérifier} : $E_p \ge 0$, unité en $\si{\joule}$, ordre de grandeur cohérent.
\end{enumerate}
\end{methode}

\begin{attention}[Pièges classiques au Probatoire]
\begin{itemize}
    \item \textbf{Oublier le facteur $\frac{1}{2}$ :} Écrire $E_p = k x^2$ est une erreur majeure. Rappelle-toi : c'est l'aire d'un \emph{triangle} (base $\times$ hauteur / 2), pas d'un rectangle.
    \item \textbf{Utiliser $x$ en centimètres :} Si $x = \SI{10}{cm}$, il faut utiliser $x = \SI{0,10}{m}$. Sinon $E_p$ sera fausse d'un facteur $10^4$ !
    \item \textbf{Confondre $x$ (allongement) et $L$ (longueur) :} $x = L - L_0$ où $L_0$ est la longueur au repos. Lis bien l'énoncé : \guillemotleft allongement de 5 cm \guillemotright $\neq$ \guillemotleft longueur de 5 cm \guillemotright.
    \item \textbf{Signe de $x$ :} La formule utilise $x^2$, donc le signe de $x$ n'importe pas pour $E_p$. Mais il est crucial pour le sens de la force $\vec{T} = -kx\vec{u}_x$.
\end{itemize}
\end{attention}

\begin{exoresolu}[Bloc projeté sur un ressort horizontal]
\textbf{Énoncé :} Un bloc $B$ de masse $m = \SI{0,50}{kg}$ glisse sans frottement sur un rail horizontal. Il heurte un ressort de raideur $k = \SI{200}{N/m}$ fixé à un mur. Au contact du ressort ($x=0$), la vitesse du bloc est $v_0 = \SI{2,0}{m/s}$. Le bloc comprime le ressort jusqu'à l'arrêt momentané.
\begin{enumerate}
    \item Calculer l'énergie cinétique initiale $E_{c0}$ du bloc.
    \item En supposant que l'énergie mécanique est conservée (pas de frottement), déterminer la compression maximale $x_{\text{max}}$ du ressort.
    \item Calculer la force maximale exercée par le ressort sur le bloc.
\end{enumerate}
\textbf{Solution détaillée :}
\begin{enumerate}
    \item \textbf{Énergie cinétique initiale :}
    \[
    E_{c0} = \frac{1}{2} m v_0^2 = \frac{1}{2} \times \num{0,50} \times (\num{2,0})^2 = \num{0,25} \times \num{4,0} = \SI{1,00}{J}
    \]
    \item \textbf{Conservation de l'énergie mécanique :} Le système \{Bloc + Ressort + Terre\} est isolé (poids et normale ne travaillent pas, pas de frottement). L'énergie mécanique $E_m = E_c + E_p^{\text{élas}}$ est constante.
    Au départ ($x=0$) : $E_m = E_{c0} + 0 = \SI{1,00}{J}$.
    À l'arrêt maximal ($v=0$, $x=x_{\text{max}}$) : $E_m = 0 + \frac{1}{2} k x_{\text{max}}^2$.
    \[
    \frac{1}{2} k x_{\text{max}}^2 = \num{1,00} \implies x_{\text{max}}^2 = \frac{2 \times \num{1,00}}{\num{200}} = \num{0,0100} \implies x_{\text{max}} = \SI{0,100}{m} = \SI{10,0}{cm}
    \]
    (On prend la racine positive pour la grandeur de la compression).
    \item \textbf{Force maximale :} Elle a lieu pour $x = x_{\text{max}}$.
    \[
    F_{\text{max}} = k x_{\text{max}} = \num{200} \times \num{0,100} = \SI{20,0}{N}
    \]
    Cette force est la force de rappel qui freine le bloc (direction opposée à la vitesse).
\end{enumerate}
\end{exoresolu}

\begin{savaistu}[Le benskin et l'énergie élastique]
Au Cameroun, le \cle{benskin} (moto-taxi) est le roi de la mobilité urbaine. Sa suspension avant (fourche télescopique) combine un ressort hélicoïdal (énergie potentielle élastique) et un amortisseur hydraulique (force de frottement visqueux, non conservative).
\begin{itemize}
    \item \textbf{Sur une bosse :} La roue monte, le ressort se comprime ($x<0$). L'énergie cinétique verticale de la moto + passager se convertit en $E_p^{\text{élas}}$. Le confort dépend de la raideur $k$ : trop raide $\to$ secousses violentes ; trop souple $\to$ \guillemotleft talonnage \guillemotright (butée en fin de course).
    \item \textbf{Dans un trou :} La roue descend, le ressort se détend ($x>0$). L'énergie élastique stockée est restituée en énergie cinétique pour repousser la roue vers le bas et garder l'adhérence.
    \item \textbf{Rôle de l'amortisseur :} Sans lui, la moto rebondirait indéfiniment (oscillations). L'amortisseur dissipe l'énergie mécanique en chaleur ($W_{\text{amort}} < 0$), ce qui assure la \cle{stabilité} et le \cle{confort}. C'est une application directe du bilan énergétique avec forces non conservatrices que nous verrons en Section 5.
\end{itemize}
\end{savaistu}

\courssub{Bilan de la section}

\begin{aretenir}[Résumé : Énergie potentielle élastique]
\begin{itemize}
    \item \textbf{Système :} Ressort linéaire (raideur $k$) + solide mobile.
    \item \textbf{Variable d'état :} Allongement algébrique $x$ (en $\si{\meter}$) par rapport à la position d'équilibre ($x=0$).
    \item \textbf{Force de rappel :} $\vec{T} = -k x \vec{u}_x$ (conservatrice).
    \item \textbf{Énergie potentielle élastique :} $\displaystyle E_p(x) = \frac{1}{2} k x^2 \ge 0$ (en $\si{\joule}$).
    \item \textbf{Travail de la force de rappel :} $W_{A \to B}(T) = -\Delta E_p = \frac{1}{2}k x_A^2 - \frac{1}{2}k x_B^2$.
    \item \textbf{Graphique :} $E_p(x)$ est une parabole symétrique, minimum en $x=0$.
    \item \textbf{Unités obligatoires :} $k$ en $\si{\newton\per\meter}$, $x$ en $\si{\meter}$ $\implies$ $E_p$ en $\si{\joule}$.
\end{itemize}
\end{aretenir}

Nous savons maintenant exprimer les deux principales énergies potentielles du programme : la pesanteur ($mgz$) et l'élastique ($\frac{1}{2}kx^2$). Dans la section suivante, nous les réunirons au sein de l'\cle{énergie mécanique} et énoncerons le principe fondamental de sa conservation.

\coursec{Énergie potentielle et énergie mécanique}

\courssub{Notion d'énergie potentielle}

Lorsqu'un footballeur tire un penalty, le ballon monte, ralentit, s'arrête un instant au sommet de sa trajectoire, puis redescend en accélérant. Son énergie cinétique a diminué puis augmenté : elle n'a pas disparu, elle a été « stockée » quelque part. Ce « quelque part », c'est l'\textbf{énergie potentielle}. Elle est associée à la \textbf{configuration} du système dans un champ de forces (pesanteur, élastique, électrique...).

\begin{experience}[Mise en évidence de l'énergie potentielle : le pendule de Galilée]
\textbf{Matériel :} Un fil inextensible, une masse (bille de métal), un support, une butée (clou) placée sur la trajectoire.
\textbf{Protocole :} On lâche la bille depuis une hauteur $h$ (position $A$). Elle passe par le point bas $B$ et remonte de l'autre côté.
\begin{enumerate}
    \item Sans butée : la bille remonte jusqu'à la hauteur initiale $h$ (position $C$). $v_C = 0$.
    \item Avec butée : la longueur effective du pendule diminue après $B$. La bille remonte \textbf{exactement à la même hauteur $h$} (position $C'$), mais avec un angle plus grand.
\end{enumerate}
\textbf{Interprétation :} Au point bas $B$, l'énergie cinétique est maximale. En $C$ et $C'$, elle est nulle. L'énergie cinétique a été convertie en une énergie de position : l'énergie potentielle de pesanteur. Le fait que la hauteur finale ne dépende \textbf{pas} du chemin (longueur du fil) montre que le travail du poids ne dépend que des positions initiale et finale.
\end{experience}

\begin{definition}[Force conservative et énergie potentielle]
Une force $\vec{F}$ est dite \textbf{conservative} si son travail entre deux points $A$ et $B$ ne dépend \textbf{que} de ces deux points, et non du chemin suivi. Équivalentement, son travail le long de tout chemin fermé est nul : $\oint \vec{F} \cdot d\vec{l} = 0$.

À toute force conservative, on peut associer une grandeur scalaire, l'\textbf{énergie potentielle} $E_p$, telle que la variation d'énergie potentielle entre $A$ et $B$ soit l'opposé du travail de la force :
\[
\Delta E_{p, A\to B} = E_{p,B} - E_{p,A} = -W_{A\to B}(\vec{F})
\]
L'énergie potentielle est définie à une constante additive près (choix de l'origine).
\end{definition}

\begin{propriete}[Critère de conservativité (niveau 1ère)]
En Première, on retient que les forces suivantes sont \textbf{conservatives} :
\begin{itemize}
    \item Le \textbf{poids} $\vec{P} = m\vec{g}$ (champ de pesanteur uniforme).
    \item La \textbf{force de rappel élastique} $\vec{F}_r = -k\vec{x}$ (ressort idéal, loi de Hooke).
    \item La force électrostatique (programme Terminale), la force de gravitation universelle.
\end{itemize}
Les forces suivantes sont \textbf{non conservatives} (ou dissipatives) :
\begin{itemize}
    \item Les \textbf{forces de frottement} (solides ou fluides) : travail toujours négatif (ou nul), proportionnel à la longueur du chemin.
    \item Les forces de \textbf{choc} (plastique), les forces \textbf{motrices} (moteur, muscle).
\end{itemize}
\end{propriete}

\begin{attention}[Énergie potentielle : variations vs valeurs absolues]
L'énergie potentielle n'a de sens physique que par ses \textbf{variations} $\Delta E_p$. On choisit arbitrairement une position de référence où $E_p = 0$ (origine). \textbf{Dans un même problème, l'origine doit être fixée une fois pour toutes.} Changer d'origine en cours de calcul conduit à des erreurs grossières.
\end{attention}

\courssub{Énergie potentielle de pesanteur}

\courssub{Expression et variation}
Considérons un solide de masse $m$ au voisinage de la Terre, soumis à son poids $\vec{P} = m\vec{g}$. On choisit un repère terrestre $(O, \vec{z})$ vertical ascendant. L'origine $z=0$ est choisie arbitrairement (souvent le sol ou le point le plus bas).

\begin{propriete}[Énergie potentielle de pesanteur]
L'énergie potentielle de pesanteur du système \textbf{\{solide–Terre\}} est :
\[
E_{p,\text{pes}} = mgz
\]
où $z$ est l'altitude (coordonnée verticale) du centre de gravité du solide par rapport à l'origine choisie.
Sa variation entre deux positions $A(z_A)$ et $B(z_B)$ est :
\[
\Delta E_{p,\text{pes}} = mg(z_B - z_A) = mg\Delta z
\]
\textbf{Analyse dimensionnelle :} $[m] = \text{kg}$, $[g] = \text{N/kg} = \text{m/s}^2$, $[z] = \text{m}$ $\Rightarrow$ $[E_p] = \text{kg} \cdot \text{m}^2/\text{s}^2 = \text{J}$.
\end{propriete}

\begin{exemplebox}[Barrage de Lom Pangar : l'énergie de l'eau stockée]
Le barrage de Lom Pangar (région de l'Est) retient un volume d'eau $V \approx 6 \times 10^9\,\text{m}^3$ (capacité utile). La masse d'eau est $m = \rho V \approx 1000 \times 6 \times 10^9 = 6 \times 10^{12}\,\text{kg}$. Le centre de gravité de la retenue est approximativement à $h = 30\,\text{m}$ au-dessus de la turbine (origine $z=0$).
L'énergie potentielle de pesanteur stockée (réservoir plein) vaut :
\[
E_p = mgh \approx 6 \times 10^{12} \times 9,81 \times 30 \approx \SI{1,77e15}{J}
\]
Cette énergie potentielle sera convertie en énergie cinétique de l'eau dans les conduites forcées, puis en énergie électrique via les alternateurs (puissance installée $\approx 30\,\text{MW}$). C'est la conversion $E_p \to E_c \to E_{\text{électrique}}$.
\end{exemplebox}

\begin{popfigure}
\begin{tikzpicture}[scale=1.2]
% Axes
\draw[->, thick, popdark] (0,0) -- (5,0) node[right] {$z$ (m)};
\draw[->, thick, popdark] (0,0) -- (0,4) node[above] {$E_p$ (J)};
% Origine
\node[below left] at (0,0) {$O$};
% Droite Ep = mgz (m=1kg, g=10)
\draw[popgreen, very thick, domain=0:4] plot (\x, {10*\x}) node[above right, popgreen] {$E_p = mgz$};
% Points
\draw[poporange, thick] (2,0) node[below] {$z_A$} -- (2,20) node[left] {$E_{p,A}$};
\draw[poporange, thick] (3.5,0) node[below] {$z_B$} -- (3.5,35) node[left] {$E_{p,B}$};
% Delta z
\draw[<->, poporange, thick] (2, -0.3) -- (3.5, -0.3) node[midway, below] {$\Delta z > 0$};
% Delta Ep
\draw[<->, poporange, thick] (4.2, 20) -- (4.2, 35) node[midway, right] {$\Delta E_p > 0$};
\end{tikzpicture}
\legende{Énergie potentielle de pesanteur en fonction de l'altitude $z$ pour $m=1\,\text{kg}$, $g=10\,\text{N/kg}$. La pente de la droite est $mg$.}
\end{popfigure}

\courssub{Énergie potentielle élastique}

\courssub{Expression et variation}
Considérons un ressort idéal (masse négligeable, loi de Hooke linéaire) de raideur $k$ (en $\text{N/m}$), de longueur au repos $L_0$. On fixe un repère $(O, \vec{x})$ le long de l'axe du ressort, origine $O$ à l'extrémité mobile au repos ($x=0$ correspond à $L=L_0$). L'allongement (ou compression) est $x = L - L_0$. La force de rappel exercée par le ressort sur le mobile est $\vec{F}_r = -kx\,\vec{u}_x$.

\begin{propriete}[Énergie potentielle élastique]
L'énergie potentielle élastique du système \textbf{\{masse–ressort\}} est :
\[
E_{p,\text{él}} = \frac{1}{2}kx^2
\]
Sa variation entre deux positions $A(x_A)$ et $B(x_B)$ est :
\[
\Delta E_{p,\text{él}} = \frac{1}{2}k(x_B^2 - x_A^2)
\]
\textbf{Analyse dimensionnelle :} $[k] = \text{N/m} = \text{kg/s}^2$, $[x^2] = \text{m}^2$ $\Rightarrow$ $[E_p] = \text{kg} \cdot \text{m}^2/\text{s}^2 = \text{J}$.
\end{propriete}

\begin{attention}[Variable $x$ : allongement algébrique]
$x$ est l'allongement \textbf{algébrique} : $x > 0$ si le ressort est étiré, $x < 0$ s'il est compressé. Comme $E_{p,\text{él}} \propto x^2$, l'énergie potentielle est toujours positive (ou nulle) et ne dépend que de l'\textbf{écart} à la longueur de repos, pas du sens de la déformation. L'origine $x=0$ (ressort détendu) est le minimum d'énergie potentielle.
\end{attention}

\begin{popfigure}
\begin{tikzpicture}[scale=1.1]
% Ressort au repos
\draw[popdark, thick] (0,0) -- (0.5,0);
\draw[popblue, very thick, decorate, decoration={coil, aspect=0.4, segment length=3mm, amplitude=3mm}] (0.5,0) -- (2.5,0);
\draw[popdark, thick] (2.5,0) -- (3,0);
\node[below] at (1.5, -0.8) {Repos : $x=0$, $E_p=0$};
% Masse
\draw[poporange, fill=poporange!30, thick] (2.5,-0.5) rectangle (3.5,0.5);
\node at (3,0) {$m$};

% Ressort étiré (décalé vers le bas)
\begin{scope}[yshift=-2.5cm]
\draw[popdark, thick] (0,0) -- (0.5,0);
\draw[popblue, very thick, decorate, decoration={coil, aspect=0.4, segment length=5mm, amplitude=3mm}] (0.5,0) -- (3.5,0);
\draw[popdark, thick] (3.5,0) -- (4,0);
\draw[<->, poporange, thick] (2.5, -0.8) -- (3.5, -0.8) node[midway, below] {$x > 0$};
\draw[poporange, fill=poporange!30, thick] (3.5,-0.5) rectangle (4.5,0.5);
\node at (4,0) {$m$};
\node[below] at (2, -1.5) {Étiré : $E_p = \frac{1}{2}kx^2$};
\end{scope}

% Ressort compressé (décalé vers le haut)
\begin{scope}[yshift=2.5cm]
\draw[popdark, thick] (0,0) -- (0.5,0);
\draw[popblue, very thick, decorate, decoration={coil, aspect=0.4, segment length=1.5mm, amplitude=3mm}] (0.5,0) -- (1.5,0);
\draw[popdark, thick] (1.5,0) -- (2,0);
\draw[<->, poporange, thick] (1.5, 0.8) -- (2.5, 0.8) node[midway, above] {$x < 0$};
\draw[poporange, fill=poporange!30, thick] (1.5,-0.5) rectangle (2.5,0.5);
\node at (2,0) {$m$};
\node[above] at (1, 1.5) {Compressé : $E_p = \frac{1}{2}kx^2$};
\end{scope}
\end{tikzpicture}
\legende{Énergie potentielle élastique : symétrique par rapport à la position de repos ($x=0$). Elle est stockée aussi bien en traction qu'en compression.}
\end{popfigure}

\begin{exemplebox}[Suspension d'un véhicule : le « dos d'âne »]
Une voiture de masse $m = 1200\,\text{kg}$ passe sur un dos d'âne. Chaque roue est soutenue par un ressort de raideur $k = 35\,000\,\text{N/m}$ (4 ressorts au total, $k_{\text{eq}} = 140\,000\,\text{N/m}$). La voiture s'enfonce de $x = 4\,\text{cm} = 0,04\,\text{m}$ par rapport à sa position d'équilibre statique.
L'énergie potentielle élastique stockée dans les 4 ressorts à cet instant est :
\[
E_{p,\text{él}} = \frac{1}{2}k_{\text{eq}}x^2 = \frac{1}{2} \times 140\,000 \times (0,04)^2 = \SI{112}{J}
\]
Cette énergie sera restituée pour repousser la caisse vers le haut (confort, tenue de route). Si les amortisseurs sont défaillants, cette énergie provoque des oscillations parasites (rebonds).
\end{exemplebox}

\courssub{Énergie mécanique et sa conservation}

Après avoir introduit les deux grandes formes d'énergie potentielle que nous rencontrerons en Première — l'énergie potentielle de pesanteur et l'énergie potentielle élastique —, nous sommes maintenant en mesure de définir la grandeur centrale de ce chapitre : l'énergie mécanique. C'est elle qui va nous permettre de résoudre avec une élégance redoutable des problèmes de dynamique qui seraient autrement très complexes (équations différentielles du second ordre, forces variables, trajectoires courbes). Comme le disait le prix Nobel Richard Feynman : « La nature utilise seulement les fils les plus longs pour tisser ses motifs, de sorte que chaque petit morceau de son tissu révèle l'organisation de la tapisserie entière. » L'énergie mécanique est l'un de ces fils conducteurs.

\courssub{Définition de l'énergie mécanique}

Rappelons-nous : l'énergie cinétique $E_c = \frac{1}{2}mv^2$ quantifie l'énergie du mouvement. L'énergie potentielle $E_p$ (de pesanteur $mgz$ ou élastique $\frac{1}{2}kx^2$) quantifie l'énergie « stockée » par la configuration du système dans un champ de forces conservatif. Il est naturel de les additionner.

\begin{definition}[Énergie mécanique d'un système]
L'\textbf{énergie mécanique} d'un système, notée $E_m$, est la somme de son énergie cinétique et de l'ensemble de ses énergies potentielles :
\[
E_m = E_c + \sum_i E_{p,i}
\]
Pour les systèmes que nous étudierons (solide en translation sans rotation, soumis à la pesanteur et/ou à une force de rappel élastique), cela s'écrit concrètement :
\begin{itemize}
    \item Système \{solide–Terre\} (pesanteur seule) : $E_m = \frac{1}{2}mv^2 + mgz$.
    \item Système \{masse–ressort\} horizontal (ressort seul) : $E_m = \frac{1}{2}mv^2 + \frac{1}{2}kx^2$.
    \item Système \{masse–ressort\} vertical (pesanteur + ressort) : $E_m = \frac{1}{2}mv^2 + mgz + \frac{1}{2}kx^2$.
\end{itemize}
$E_m$ s'exprime en \textbf{joules (J)}. C'est une grandeur scalaire, additive et dépendante du référentiel d'étude (via $v$ et $z$).
\end{definition}

\begin{attention}[Choix de l'origine des énergies potentielles]
L'énergie mécanique n'est définie qu'à une constante additive près, car chaque énergie potentielle possède sa propre origine arbitraire ($z=0$ pour la pesanteur, $x=0$ pour le ressort). \textbf{Dans un même problème, il faut fixer une fois pour toutes les origines de $z$ et $x$ et ne plus les changer.} Seules les \emph{variations} d'énergie mécanique $\Delta E_m$ ont un sens physique absolu.
\end{attention}

\courssub{Principe de conservation de l'énergie mécanique}

C'est le résultat le plus puissant de la mécanique newtonienne pour les forces conservatives. Nous allons le démontrer rigoureusement à partir du théorème de l'énergie cinétique, que vous maîtrisez depuis la leçon précédente.

\begin{experience}[Démonstration guidée : du théorème de l'énergie cinétique à la conservation de $E_m$]
\textbf{Objectif :} Prouver que si les seules forces qui travaillent sont conservatives, $E_m$ est constante.

\textbf{Étape 1 : Théorème de l'énergie cinétique.} Pour un système (solide indéformable en translation) entre deux positions $A$ et $B$ :
\[
\Delta E_c = E_{c,B} - E_{c,A} = W_{A\to B}(\text{forces extérieures})
\]

\textbf{Étape 2 : Travail des forces conservatives.} Par définition d'une force conservative (poids $\vec{P}$, force de rappel $\vec{F}_r = -k\vec{x}$), le travail de cette force est l'opposé de la variation de l'énergie potentielle associée :
\[
W_{A\to B}(\vec{P}) = -\Delta E_{p,\text{pes}} = -(mgz_B - mgz_A)
\]
\[
W_{A\to B}(\vec{F}_r) = -\Delta E_{p,\text{él}} = -\left(\frac{1}{2}kx_B^2 - \frac{1}{2}kx_A^2\right)
\]

\textbf{Étape 3 : Hypothèse de conservation.} Supposons que \textbf{toutes} les forces extérieures qui travaillent sont conservatives (ou que les forces non conservatives — frottements, moteur, main qui pousse — ont un travail nul). Alors :
\[
W_{A\to B}(\text{ext}) = \sum W_{\text{conservatives}} = -\sum \Delta E_{p,i} = -\Delta E_p
\]
où $\Delta E_p$ est la variation de la \emph{somme} des énergies potentielles.

\textbf{Étape 4 : Conclusion.}
\[
\Delta E_c = -\Delta E_p \quad \Longrightarrow \quad \Delta E_c + \Delta E_p = 0 \quad \Longrightarrow \quad \Delta E_m = 0
\]
L'énergie mécanique ne varie pas : elle est \textbf{conservée}.
\end{experience}

\begin{propriete}[Principe de conservation de l'énergie mécanique]
Si un système n'est soumis qu'à des forces \textbf{conservatives} (poids, force de rappel élastique) et/ou à des forces dont le \textbf{travail est nul} (réaction de liaison lisse sans frottement, force magnétique, force de Coriolis, tension d'un fil inextensible...), alors son énergie mécanique reste constante au cours du temps :
\[
E_m = \text{constante} \quad \text{ou} \quad \Delta E_m = 0 \quad \text{ou} \quad E_{m,A} = E_{m,B} \quad \forall A,B
\]
\textbf{Conséquence immédiate (transformations mutuelles) :}
\[
\Delta E_c = -\Delta E_p
\]
Toute perte d'énergie potentielle se convertit intégralement en gain d'énergie cinétique, et réciproquement.
\end{propriete}

\begin{savaistu}[Le mot « conservative » : une question de chemin]
Une force est dite \emph{conservative} si son travail ne dépend \textbf{que} des positions initiale et finale, \textbf{pas} du chemin suivi. Le poids et la force de rappel élastique ont cette propriété. Les forces de frottement, elles, sont \emph{non conservatives} : plus le chemin est long, plus elles dissipent d'énergie en chaleur. C'est pour cela qu'une bille qui descend une piste sinueuse sans frottement arrive avec la même vitesse que si elle tombait verticalement de la même hauteur : seule la différence d'altitude compte pour le travail du poids.
\end{savaistu}

\courssub{Illustration : diagramme énergétique et évolution graphique}

Pour visualiser cette conservation, rien ne vaut un diagramme à barres (représentation instantanée) et une courbe d'évolution (représentation dynamique).

\begin{popfigure}
\begin{tikzpicture}[scale=0.9, every node/.style={font=\small}]
% --- Styles ---
\tikzset{
    bar/.style={draw=popdark, thick, fill=popblue!60, minimum width=1.2cm},
    barec/.style={bar, fill=poporange!70},
    barep/.style={bar, fill=popgreen!70},
    barem/.style={bar, fill=poppurple!70, pattern=north east lines, pattern color=poppurple},
    ground/.style={draw=popdark, thick, fill=popmuted!30},
    ball/.style={circle, draw=popdark, thick, fill=poppink, minimum size=0.6cm},
    label/.style={anchor=north, yshift=-0.2cm, font=\small\bfseries},
    arr/.style={->, thick, popdark},
}
% --- Positions : A (haut), B (milieu), C (bas) ---
\def\xA{1.5} \def\yA{3.5}
\def\xB{5.5} \def\yB{2.0}
\def\xC{9.5} \def\yC{0.5}
\def\groundY{0}

% --- Piste (courbe lisse) ---
\draw[popdark, very thick] (\xA,\yA) .. controls (3,4) and (4,1.5) .. (\xB,\yB) .. controls (7,1) and (8.5,0) .. (\xC,\yC);
% Sol
\draw[ground] (0,\groundY) -- (11,\groundY) -- (11,-0.5) -- (0,-0.5) -- cycle;
\node[label] at (5.5,-0.25) {Référentiel terrestre (sol)};

% --- Billes et vecteurs vitesse ---
\foreach \x/\y/\vx/\vy/\pos in {
    \xA/\yA/0.8/0.2/A,
    \xB/\yB/1.2/-0.3/B,
    \xC/\yC/1.0/-0.5/C
} {
    \node[ball] (b\pos) at (\x,\y) {};
    \draw[arr, poporange] (b\pos.center) -- ++(\vx,\vy) node[midway, above right, poporange] {$\vec{v}_\pos$};
    % Hauteur z
    \draw[dashed, popmuted] (\x,\y) -- (\x,\groundY);
    \draw[<->, popmuted, thin] (\x+0.3,\y) -- (\x+0.3,\groundY) node[midway, right] {$z_\pos$};
}

% --- Diagrammes à barres (à droite) ---
\def\barX{11.5}
\def\barW{1.2}
\def\maxH{3.0} % hauteur max pour Em = 10 unités arbitraires
% Valeurs normalisées : Em = 10. A: Ep=10, Ec=0. B: Ep=4, Ec=6. C: Ep=1, Ec=9.
\foreach \pos/\hEp/\hEc/\yBase in {
    A/3.0/0.0/0.5,
    B/1.2/1.8/2.5,
    C/0.3/2.7/4.5
} {
    % Em total (fond)
    \draw[barem] (\barX,\yBase) rectangle ++(\barW,\maxH) node[midway, rotate=90, white, font=\tiny] {$E_m = \text{cste}$};
    % Ep
    \draw[barep] (\barX,\yBase) rectangle ++(\barW,\hEp) node[midway, right=0.1cm, popgreen] {$E_{p,\pos}$};
    % Ec (au-dessus de Ep)
    \draw[barec] (\barX,\yBase+\hEp) rectangle ++(\barW,\hEc) node[midway, right=0.1cm, poporange] {$E_{c,\pos}$};
    % Légende position
    \node[anchor=west] at (\barX+\barW+0.3,\yBase+\maxH/2) {\textbf{Position \pos}};
}

% Légendes globales (yBase du dernier groupe = 4.5, fixé en dur car \foreach ne conserve pas sa variable hors boucle)
\node[anchor=west, poporange] at (\barX+0.1,4.5+\maxH+0.5) {\textbf{$E_c$ (cinétique)}};
\node[anchor=west, popgreen] at (\barX+0.1,4.5+\maxH+0.9) {\textbf{$E_p$ (potentielle)}};
\node[anchor=west, poppurple] at (\barX+0.1,4.5+\maxH+1.3) {\textbf{$E_m = E_c + E_p$ (constante)}};
\end{tikzpicture}
\legende{Diagramme énergétique pour une bille glissant sans frottement sur une piste courbe. En A (haut) : $E_p$ max, $E_c=0$. En B : conversion partielle. En C (bas) : $E_p$ min, $E_c$ max. La hauteur totale des barres hachurées ($E_m$) ne change pas.}
\end{popfigure}

\begin{popfigure}
\begin{tikzpicture}
\begin{axis}[
    width=\linewidth, height=7cm,
    axis lines=middle,
    xlabel={$z$ (m) -- Hauteur}, ylabel={Énergie (J)},
    xmin=0, xmax=5.5, ymin=0, ymax=55,
    xtick={0,1,2,3,4,5}, ytick={0,10,20,30,40,50},
    legend style={at={(0.98,0.98)}, anchor=north east, fill=popcream, draw=popdark},
    domain=0:5, samples=100,
    tick label style={font=\small}, label style={font=\small},
    grid=major, grid style={popline},
]
% Constantes : m=1kg, g=10 N/kg, h0=5m => Em = 50 J
% Ep = mgz = 10z
% Ec = Em - Ep = 50 - 10z
\addplot[popgreen, thick] {10*x} node[above, pos=0.8, popgreen] {$E_p = mgz$};
\addplot[poporange, thick] {50 - 10*x} node[below, pos=0.2, poporange] {$E_c = E_m - E_p$};
\addplot[poppurple, very thick, dashed] {50} node[above, pos=0.5, poppurple] {$E_m = 50\,\text{J}$};
\legend{$E_p$, $E_c$, $E_m$}
% Point de départ (z=5, v=0)
\addplot[only marks, mark=*, mark size=3, popdark] coordinates {(5,0)} node[above right] {Lâcher ($z=h$)};
\addplot[only marks, mark=*, mark size=3, popdark] coordinates {(5,50)};
% Point d'arrivée (z=0)
\addplot[only marks, mark=*, mark size=3, popdark] coordinates {(0,50)} node[below right] {Arrivée ($z=0$)};
\addplot[only marks, mark=*, mark size=3, popdark] coordinates {(0,0)};
\end{axis}
\end{tikzpicture}
\legende{Évolution de $E_c$, $E_p$ et $E_m$ en fonction de l'altitude $z$ pour une chute libre sans frottements ($m=1\,\text{kg}$, $h=5\,\text{m}$, $g=10\,\text{N/kg}$). $E_m$ est une ligne horizontale. $E_p$ et $E_c$ sont des droites de pentes opposées.}
\end{popfigure}

\courssub{Exemples classiques de conservation de l'énergie mécanique}

Trois situations « paradigmes » du programme illustrent parfaitement ce principe. Dans chaque cas, on identifie le système, on liste les forces, on vérifie l'hypothèse (forces non conservatives négligeables ou travail nul), puis on écrit $E_{m,A} = E_{m,B}$.

\begin{exemplebox}[Chute libre sans frottements (revisité)]
Un objet de masse $m$ est lâché sans vitesse initiale ($v_A=0$) d'une hauteur $h$ au-dessus du sol ($z_A=h$, $z_B=0$). Système : \{objet–Terre\}. Forces : poids (conservative), réaction de l'air (négligée).
\[
E_{m,A} = E_{m,B} \implies \underbrace{\frac{1}{2}m v_A^2}_{0} + mgh = \frac{1}{2}m v_B^2 + \underbrace{mg\cdot0}_{0}
\]
\[
v_B = \sqrt{2gh}
\]
Remarquons que la masse $m$ se simplifie : \textbf{tous les corps tombent avec la même vitesse} (principe d'équivalence galiléen), ce que Galilée avait deviné et que l'expérience de la tour de Pise (ou le marteau et la plume sur la Lune) a confirmé.
\end{exemplebox}

\begin{exemplebox}[Pendule simple sans frottements]
Une masse $m$ attachée à un fil inextensible de longueur $L$ (masse négligeable) oscille. Système : \{masse–Terre\}. Forces : poids (conservative), tension du fil $\vec{T}$ (travail nul car $\vec{T} \perp \vec{v}$ à tout instant).
Choisissons l'origine de $z$ au point bas ($z=0$). En position extrême $A$ : $v_A=0$, $z_A = L(1-\cos\alpha_0)$. En position générique $B$ : $z_B = L(1-\cos\alpha)$.
\[
mgL(1-\cos\alpha_0) = \frac{1}{2}mv^2 + mgL(1-\cos\alpha)
\]
\[
v = \sqrt{2gL(\cos\alpha - \cos\alpha_0)}
\]
Au point bas ($\alpha=0$) : $v_{\text{max}} = \sqrt{2gL(1-\cos\alpha_0)}$. Pour de petits angles ($\cos\alpha \approx 1-\alpha^2/2$), on retrouve le mouvement harmonique simple.
\end{exemplebox}

\begin{exemplebox}[Oscillateur masse–ressort horizontal non amorti]
Un solide de masse $m$ sur un rail horizontal lisse (pas de frottement), attaché à un ressort de raideur $k$ (masse négligeable). Système : \{masse–ressort\}. Forces : force de rappel $\vec{F}_r = -k\vec{x}$ (conservative), poids et réaction du rail (travail nul).
Origine $x=0$ à la position d'équilibre (ressort détendu). Lâché de $x_A = A$ (amplitude) avec $v_A=0$.
\[
E_{m,A} = E_{m,B} \implies \frac{1}{2}kA^2 + 0 = \frac{1}{2}kx^2 + \frac{1}{2}mv^2
\]
\[
v(x) = \sqrt{\frac{k}{m}(A^2 - x^2)} = \omega_0\sqrt{A^2 - x^2} \quad \text{avec } \omega_0 = \sqrt{k/m}
\]
C'est l'équation de l'énergie de l'oscillateur harmonique. La vitesse est maximale au passage par l'équilibre ($x=0$) : $v_{\text{max}} = \omega_0 A$.
\end{exemplebox}

\courssub{Méthode de résolution : exploiter la conservation}

\begin{methode}[Appliquer la conservation de l'énergie mécanique]
\textbf{Étape 1 : Définir le système.} Inclure la Terre si le poids travaille, inclure le ressort si la force de rappel travaille.
\textbf{Étape 2 : Dresser le bilan des forces.} Identifier les forces conservatives (poids, rappel élastique) et les forces dont le travail est nul (liaisons lisses, force magnétique, tension de fil inextensible). \textbf{Vérifier qu'aucune force non conservative (frottements, moteur, choc plastique) ne travaille.}
\textbf{Étape 3 : Choisir deux positions pertinentes $A$ et $B$.} Typiquement : position initiale (souvent $v=0$ ou connue) et position finale (demandée ou inconnue). Une position intermédiaire si on cherche une vitesse en cours de route.
\textbf{Étape 4 : Fixer les origines des énergies potentielles.} $z=0$ pour la pesanteur (souvent le point le plus bas), $x=0$ pour le ressort (position détendue). \textbf{Ne pas changer d'origine entre $A$ et $B$.}
\textbf{Étape 5 : Écrire $E_{m,A} = E_{m,B}$.} Développer : $\frac{1}{2}mv_A^2 + mgz_A + \frac{1}{2}kx_A^2 = \frac{1}{2}mv_B^2 + mgz_B + \frac{1}{2}kx_B^2$.
\textbf{Étape 6 : Résoudre l'inconnue (vitesse, hauteur, compression, raideur...).} Vérifier l'homogénéité et l'ordre de grandeur.
\end{methode}

\courssub{Exemple chiffré détaillé : la mangue qui tombe}

Au Cameroun, en saison des pluies, les manguiers offrent des exemples parfaits de chute libre (quasi sans frottement pour une mangue dense).

\begin{exoresolu}[Vitesse d'impact d'une mangue]
\textbf{Énoncé :} Une mangue de masse $m = 300\,\text{g}$ se détache d'une branche située à $h = 12\,\text{m}$ au-dessus du sol. On néglige les frottements de l'air. On prend $g = 9,81\,\text{N/kg}$.
\begin{enumerate}
    \item Calculer son énergie mécanique au départ (position $A$).
    \item Déterminer sa vitesse $v_B$ juste avant l'impact au sol (position $B$).
    \item Calculer son énergie cinétique à $z = 6\,\text{m}$.
\end{enumerate}

\textbf{Solution détaillée}

\textbf{1. Système et hypothèses.} Système : \{mangue–Terre\}. Forces : $\vec{P}=m\vec{g}$ (conservative), frottements air négligés. $\Rightarrow$ Conservation de $E_m$ applicable.

\textbf{2. Choix des origines.} Origine des hauteurs $z=0$ au niveau du sol. Position $A$ : $z_A = h = 12\,\text{m}$, $v_A = 0$. Position $B$ : $z_B = 0$.

\textbf{3. Énergie mécanique en $A$.}
\[
E_{m,A} = E_{c,A} + E_{p,A} = 0 + mgh
\]
\[
m = 0,300\,\text{kg},\quad g = 9,81\,\text{N/kg},\quad h = 12\,\text{m}
\]
\[
E_{m,A} = 0,300 \times 9,81 \times 12 = \num{35,316}\,\text{J} \approx \SI{35,3}{J}
\]
Par conservation, $E_m = \SI{35,3}{J}$ tout au long de la chute.

\textbf{4. Vitesse à l'impact $B$ ($z=0$).}
\[
E_{m,B} = E_{m,A} \implies \frac{1}{2}mv_B^2 + 0 = mgh
\]
\[
v_B = \sqrt{2gh} = \sqrt{2 \times 9,81 \times 12} = \sqrt{235,44} \approx \SI{15,3}{m/s}
\]
(Soit environ $\SI{55}{km/h}$ : attention aux passants !)

\textbf{5. Énergie cinétique à $z = 6\,\text{m}$ (mi-hauteur).}
À $z_C = 6\,\text{m}$, $E_{p,C} = mgz_C = 0,300 \times 9,81 \times 6 = \SI{17,66}{J}$.
\[
E_{c,C} = E_m - E_{p,C} = 35,316 - 17,658 = \SI{17,66}{J}
\]
Remarque : à mi-hauteur, $E_c = E_p = E_m/2$. La vitesse y est $v_C = \sqrt{2g(h/2)} = \sqrt{gh} \approx \SI{10,8}{m/s}$.
\end{exoresolu}

\courssub{Pièges à éviter et vigilance épistémologique}

\begin{attention}[Le piège n°1 : Appliquer la conservation... quand il y a frottements !]
C'est l'erreur la plus fréquente au Probatoire. \textbf{La conservation de $E_m$ ne s'applique QUE si le travail des forces non conservatives est nul.} Si un énoncé mentionne « frottements non négligeables », « amortissement », « choc inélastique », « freinage », « moteur », alors $\Delta E_m \neq 0$. Il faut alors utiliser le \textbf{théorème de l'énergie cinétique} sous sa forme générale : $\Delta E_c = W(\vec{P}) + W(\vec{F}_r) + W(\vec{f}_{\text{frott}})$, ou la relation de la section suivante : $\Delta E_m = W_{\text{non cons}}$.
\end{attention}

\begin{attention}[Le piège n°2 : Mélanger les origines de $z$ ou $x$]
Si vous prenez $z=0$ au sol pour le point $A$ et $z=0$ au point le plus bas de la trajectoire pour le point $B$, votre équation $E_{m,A}=E_{m,B}$ sera fausse. \textbf{Fixez l'origine une fois pour toutes au début du problème.} Dessinez un petit schéma avec l'origine marquée.
\end{attention}

\begin{attention}[Le piège n°3 : Oublier une énergie potentielle]
Si le système comporte \textbf{à la fois} un ressort vertical \textbf{et} la pesanteur, $E_m = \frac{1}{2}mv^2 + mgz + \frac{1}{2}kx^2$. Oublier $mgz$ ou $\frac{1}{2}kx^2$ conduit à une équation fausse. Vérifiez systématiquement : « Y a-t-il un ressort qui se déforme ? Y a-t-il un changement d'altitude ? »
\end{attention}

\begin{attention}[Le piège n°4 : Confondre « force conservative » et « force constante »]
Une force constante n'est pas nécessairement conservative (ex : force de frottement cinétique de magnitude constante mais direction opposée au mouvement $\to$ travail non nul sur un aller-retour). Une force conservative n'est pas nécessairement constante (ex : force de rappel $F=-kx$, force de gravitation universelle $F \propto 1/r^2$). Le critère est : \textbf{le travail dépend-il seulement des positions initiale et finale ?}
\end{attention}

\courssub{Non conservation de l'énergie mécanique}

Dans la réalité — sur les routes de Douala ou Yaoundé, dans les machines de l'usine SOSUCAM, dans les lignes à haute tension de la SONATREL — les frottements existent. L'énergie mécanique ne se conserve plus : elle \textbf{diminue} (freinage, amortissement) ou \textbf{augmente} (moteur, propulsion). Le principe de conservation de l'énergie \textbf{total} (mécanique + thermique + chimique + ...) reste valide, mais l'énergie mécanique seule n'est plus constante.

\begin{propriete}[Théorème de l'énergie mécanique — Forme générale]
La variation de l'énergie mécanique d'un système entre deux instants (ou positions) $A$ et $B$ est égale au travail des forces \textbf{non conservatives} (ou dont le travail n'est pas nul et non déjà comptabilisé dans $E_p$) :
\[
\Delta E_m = E_{m,B} - E_{m,A} = W_{A\to B}(\text{forces non conservatives})
\]
Si $W_{\text{non cons}} < 0$ (frottements, freinage) : $E_m$ diminue (dissipation en chaleur).
Si $W_{\text{non cons}} > 0$ (moteur, propulsion) : $E_m$ augmente (apport d'énergie externe).
\end{propriete}

\begin{experience}[Démonstration : $\Delta E_m = W_{\text{non cons}}$]
On reprend la démonstration de la conservation, mais sans l'hypothèse « forces non conservatives = 0 ».
\begin{align*}
\Delta E_c &= W_{\text{cons}} + W_{\text{non cons}} \\
\Delta E_c &= -\Delta E_p + W_{\text{non cons}} \quad (\text{car } W_{\text{cons}} = -\Delta E_p) \\
\Delta E_c + \Delta E_p &= W_{\text{non cons}} \\
\Delta E_m &= W_{\text{non cons}}
\end{align*}
C'est la relation fondamentale pour tout problème où $E_m$ n'est pas conservée.
\end{experience}

\begin{exemplebox}[Freinage d'un mototaxi à Yaoundé]
Un mototaxi (masse totale $m = 200\,\text{kg}$) roule à $v_A = \SI{72}{km/h} = \SI{20}{m/s}$. Le pilote freine brutalement ; la moto s'arrête en $d = 15\,\text{m}$ (frottements pneus/route + freins). On néglige la résistance de l'air.
\begin{enumerate}
    \item Calculer le travail des forces de frottement (freinage).
    \item En déduire la force de frottement moyenne $\bar{f}$ (supposée constante).
\end{enumerate}
\textbf{Solution :}
Système : \{moto + pilote\}. Forces : $\vec{P}, \vec{R}$ (travail nul), $\vec{f}$ (frottement, non conservative, travail négatif).
$E_{m,A} = \frac{1}{2}mv_A^2 = \frac{1}{2} \times 200 \times 20^2 = \SI{40\,000}{J}$. $E_{m,B} = 0$ (arrêt).
\[
\Delta E_m = 0 - 40\,000 = W(\vec{f}) = -\bar{f} \times d \quad (\vec{f} \text{ opposé au déplacement})
\]
\[
\bar{f} = \frac{40\,000}{15} \approx \SI{2667}{N}
\]
L'énergie mécanique initiale ($\SI{40}{kJ}$) a été entièrement dissipée en chaleur dans les freins et au contact pneu/route.
\end{exemplebox}

\begin{exemplebox}[Chute avec frottements de l'air : parachutiste]
Un parachutiste de masse $m = 80\,\text{kg}$ saute d'un avion. Au bout d'un certain temps, il atteint sa \textbf{vitesse limite} (vitesse terminale) $v_{\text{lim}} = \SI{50}{m/s}$ (environ $\SI{180}{km/h}$). Il a alors perdu une altitude $\Delta z = 500\,\text{m}$.
Calculer le travail des frottements de l'air $W_{\text{air}}$ pendant cette phase.
\textbf{Solution :}
$\Delta E_c = \frac{1}{2}m(v_{\text{lim}}^2 - 0) = \frac{1}{2} \times 80 \times 2500 = \SI{100\,000}{J} = \SI{100}{kJ}$.
$\Delta E_p = mg\Delta z = 80 \times 9,81 \times (-500) = -\SI{392\,400}{J} \approx -\SI{392,4}{kJ}$.
$\Delta E_m = \Delta E_c + \Delta E_p = 100 - 392,4 = -\SI{292,4}{kJ}$.
$W_{\text{air}} = \Delta E_m = -\SI{292,4}{kJ}$.
L'air a dissipé $\approx \SI{292}{kJ}$ en chaleur/turbulence. Seule une fraction ($\approx 25\%$) de l'énergie potentielle perdue est devenue énergie cinétique ; le reste a chauffé l'atmosphère.
\end{exemplebox}

\courssub{Bilan et méthodes de résolution}

\begin{aretenir}[Tableau récapitulatif : Énergies et Conservation]
\begin{center}
\begin{tabularx}{\linewidth}{|l|X|X|X|}\hline
\textbf{Situation} & \textbf{Forces qui travaillent} & \textbf{Énergie mécanique $E_m$} & \textbf{Relation à utiliser} \\ \hline
Conservation de $E_m$ & Uniquement conservatrices (Poids, Rappel) + Forces travail nul & \textbf{Constante} : $E_{m,A} = E_{m,B}$ & $\Delta E_c = -\Delta E_p$ \\ \hline
Non conservation ($E_m \downarrow$) & Forces non conservatives dissipatives (Frottements, chocs plastiques) & \textbf{Diminue} : $\Delta E_m < 0$ & $\Delta E_m = W_{\text{frott}} < 0$ \\ \hline
Non conservation ($E_m \uparrow$) & Forces motrices non conservatives (Moteur, propulsion) & \textbf{Augmente} : $\Delta E_m > 0$ & $\Delta E_m = W_{\text{moteur}} > 0$ \\ \hline
Cas général & Mélange & Variable & $\Delta E_m = W_{\text{non cons}}$ \\ \hline
\end{tabularx}
\end{center}
\end{aretenir}

\begin{methode}[Résolution complète : Conservation OU Non-conservation]
\textbf{1. Système \& Référentiel.} Définir le système (inclure Terre/Ressort). Repère galiléen (terrestre).
\textbf{2. Bilan des forces \& Travail.} Lister toutes les forces. Classer : Conservatives ($W = -\Delta E_p$), Travail nul ($\perp \vec{v}$), Non conservatives ($W \neq 0, \neq -\Delta E_p$).
\textbf{3. Choix de la relation.}
\begin{itemize}
    \item Si $W_{\text{non cons}} = 0$ : \textbf{Conservation} $E_{m,A} = E_{m,B}$.
    \item Sinon : \textbf{Théorème $E_m$} $\Delta E_m = W_{\text{non cons}}$.
\end{itemize}
\textbf{4. Origines.} Fixer $z=0$ (pesanteur) et $x=0$ (ressort) \textbf{une seule fois}. Schéma obligatoire.
\textbf{5. Écriture de l'équation.} Expliciter $E_c, E_{p,\text{pes}}, E_{p,\text{él}}$ en $A$ et $B$.
\textbf{6. Calcul \& Vérification.} Unités (SI), cohérence des signes ($W_{\text{frott}} < 0$), ordre de grandeur.
\end{methode}

\courssub{Exercices d'application}

\begin{exercice}[Glissade sur piste de skateboard]
Un adolescent de masse $m = 55\,\text{kg}$ sur un skateboard (masse négligeable, roulements parfaits) s'engage sans vitesse initiale en haut d'une rampe en quart de cercle de rayon $R = 3,0\,\text{m}$, lisse. On prend $g = 10\,\text{N/kg}$.
\begin{enumerate}
    \item En choisissant l'origine de l'énergie potentielle de pesanteur au point bas de la rampe, exprimer $E_m$ en fonction de l'angle $\theta$ que fait le rayon avec la verticale.
    \item Calculer sa vitesse au point bas ($\theta=0$).
    \item Quelle hauteur atteint-il de l'autre côté de la rampe (symétrique) ?
    \item Si on place un ressort de raideur $k = 2000\,\text{N/m}$ au point bas (horizontal), quelle est sa compression maximale $x_{\text{max}}$ ? (On néglige le frottement lors de la compression).
\end{enumerate}
\end{exercice}

\begin{corrige}[Exercice : Glissade sur piste de skateboard]
\textbf{1.} Système : \{ado+skate–Terre\}. Forces : $\vec{P}$ (conservative), $\vec{R}$ (réaction normale, travail nul car $\perp \vec{v}$). Conservation de $E_m$ valide.
Origine $z=0$ au point bas. Hauteur $z = R(1-\cos\theta)$.
$E_c = \frac{1}{2}mv^2$, $E_p = mgR(1-\cos\theta)$.
$E_m = \frac{1}{2}mv^2 + mgR(1-\cos\theta) = \text{cste}$.

\textbf{2.} Point haut $A$ : $\theta_A = 90^\circ \Rightarrow z_A = R$, $v_A=0$. $E_{m,A} = mgR$.
Point bas $B$ : $\theta_B = 0 \Rightarrow z_B = 0$. $E_{m,B} = \frac{1}{2}mv_B^2$.
$mgR = \frac{1}{2}mv_B^2 \Rightarrow v_B = \sqrt{2gR} = \sqrt{2\times10\times3} = \sqrt{60} \approx \SI{7,75}{m/s}$.

\textbf{3.} De l'autre côté, point $C$ symétrique : même hauteur $z_C = R$. Si pas de frottement, $E_{m,C} = E_{m,A} \Rightarrow mgR = mgz_C \Rightarrow z_C = R$. Il remonte exactement à la même hauteur (réversibilité).

\textbf{4.} Au point bas, il a $E_c = mgR = 55\times10\times3 = \SI{1650}{J}$. Il compresse le ressort horizontal.
Système : \{ado+skate–Terre–ressort\}. Forces : $\vec{P}, \vec{R}$ (travail nul), $\vec{F}_r$ (conservative). Conservation de $E_m$ valide.
État initial (contact ressort, $x=0$) : $E_m = \frac{1}{2}mv_B^2 = mgR$.
État final (compression max $x_{\text{max}}$, $v=0$) : $E_m = \frac{1}{2}kx_{\text{max}}^2$.
\[
mgR = \frac{1}{2}kx_{\text{max}}^2 \implies x_{\text{max}} = \sqrt{\frac{2mgR}{k}} = \sqrt{\frac{2\times55\times10\times3}{2000}} = \sqrt{1,65} \approx \SI{1,28}{m}
\]
\end{corrige}

\begin{exercice}[Freinage d'un camion sur la route Nationale 1]
Un camion de masse $m = 15\,\text{t} = 15\,000\,\text{kg}$ descend une pente à $\SI{10}{\%}$ (pente $\alpha \approx 5,7^\circ$, $\sin\alpha \approx 0,1$) à vitesse constante $v = \SI{72}{km/h} = \SI{20}{m/s}$. Le moteur est au ralenti (ne fournit pas de travail utile), le frein moteur et les freins de service agissent. On néglige les frottements de l'air et au roulement par rapport au freinage.
\begin{enumerate}
    \item Faire le bilan des forces et classer les forces selon qu'elles sont conservatrices, travail nul, ou non conservatives.
    \item Calculer la puissance dissipée par le freinage (en kW).
    \item Si le camion freine ainsi pendant $t = 2\,\text{min}$, quelle énergie thermique a été dégagée (en MJ) ?
\end{enumerate}
\end{exercice}

\begin{corrige}[Exercice : Freinage d'un camion]
\textbf{1.} Système : \{Camion–Terre\}. Forces :
\begin{itemize}
    \item Poids $\vec{P} = m\vec{g}$ : \textbf{Conservative} ($E_{p,\text{pes}} = mgz$).
    \item Réaction de la route $\vec{R}$ : \textbf{Travail nul} ($\perp \vec{v}$).
    \item Forces de freinage (frein moteur + freins de service) $\vec{F}_{\text{fr}}$ : \textbf{Non conservatives}, travail négatif (opposées à $\vec{v}$).
\end{itemize}
Vitesse constante $\Rightarrow \Delta E_c = 0$. Mais altitude change $\Rightarrow \Delta E_p \neq 0$. Donc $\Delta E_m \neq 0$.

\textbf{2.} Puissance du poids : $P_P = \vec{P} \cdot \vec{v} = mgv\sin\alpha$ (positif, $E_p$ diminue, $E_c$ constant $\Rightarrow$ freins absorbent cette puissance).
$P_P = 15\,000 \times 10 \times 20 \times 0,1 = \SI{300\,000}{W} = \SI{300}{kW}$.
Puissance dissipée par les freins : $P_{\text{fr}} = -P_P = \SI{300}{kW}$ (valeur absolue).

\textbf{3.} Énergie thermique $Q = |W_{\text{fr}}| = P_{\text{fr}} \times t = 300\,000 \times (2 \times 60) = \SI{36\,000\,000}{J} = \SI{36}{MJ}$.
C'est l'énergie potentielle de pesanteur perdue qui a été convertie en chaleur dans les freins (risque de surchauffe, « fading »).
\end{corrige}

\begin{exercice}[Oscillateur vertical amorti : le « saut à l'élastique » simplifié]
Un sauteur de masse $m = 70\,\text{kg}$ est attaché à un élastique de longueur au repos $L_0 = 20\,\text{m}$ et de raideur $k = 100\,\text{N/m}$. Il saute d'un pont (origine $z=0$ au pont, $z$ orienté vers le bas). On néglige la masse de l'élastique et les frottements de l'air.
\begin{enumerate}
    \item Écrire l'expression de l'énergie mécanique $E_m(z)$ pour $z > L_0$ (élastique tendu).
    \item Déterminer l'allongement maximal $x_{\text{max}}$ de l'élastique (point le plus bas).
    \item Calculer l'accélération du sauteur au point le plus bas. Commenter.
\end{enumerate}
\end{exercice}

\begin{corrige}[Exercice : Saut à l'élastique]
\textbf{1.} Système : \{Sauteur–Terre–Élastique\}. Forces : $\vec{P}$ (conservative), $\vec{F}_r$ (conservative pour $z>L_0$). $E_m$ conservée.
Origines : $z=0$ au pont (pesanteur), $x=0$ pour $z=L_0$ (élastique).
Pour $z > L_0$ : $x = z - L_0$.
$E_c = \frac{1}{2}mv^2$, $E_{p,\text{pes}} = -mgz$ (car $z$ vers le bas), $E_{p,\text{él}} = \frac{1}{2}k(z-L_0)^2$.
$E_m = \frac{1}{2}mv^2 - mgz + \frac{1}{2}k(z-L_0)^2 = \text{cste}$.
Au départ ($z=0, v=0$) : $E_{m,0} = 0$.

\textbf{2.} Point le plus bas $B$ : $v_B = 0$, $z_B = L_0 + x_{\text{max}}$.
$E_{m,B} = 0 - mg(L_0 + x_{\text{max}}) + \frac{1}{2}k x_{\text{max}}^2 = 0$.
$\frac{1}{2}k x_{\text{max}}^2 - mg x_{\text{max}} - mgL_0 = 0$.
$50 x_{\text{max}}^2 - 700 x_{\text{max}} - 14\,000 = 0 \implies x_{\text{max}}^2 - 14 x_{\text{max}} - 280 = 0$.
$\Delta = 196 + 1120 = 1316$. $\sqrt{\Delta} \approx 36,28$.
$x_{\text{max}} = \frac{14 + 36,28}{2} \approx \SI{25,1}{m}$ (racine positive).
Altitude max $z_{\text{max}} = 20 + 25,1 = \SI{45,1}{m}$ sous le pont.

\textbf{3.} Au point bas, forces : $\vec{P}$ (vers le bas), $\vec{F}_r = -k x_{\text{max}}$ (vers le haut, magnitude $k x_{\text{max}}$).
$\vec{F}_{\text{rés}} = \vec{P} + \vec{F}_r = mg - k x_{\text{max}}$ (vers le bas si positif).
$F_{\text{rés}} = 700 - 100 \times 25,1 = 700 - 2510 = -\SI{1810}{N}$.
Le résultat est négatif : la force résultante est vers le \textbf{haut} (freinage violent).
$a = F_{\text{rés}}/m \approx -25,9\,\text{m/s}^2 \approx -2,6\,g$.
Le sauteur subit une décélération de $2,6\,g$ vers le haut. C'est supportable mais intense (les élastiques réels sont progressifs pour limiter ce pic).
\end{corrige}

\coursec{Énergie potentielle et énergie mécanique}

\courssub{Notion d'énergie potentielle}

\begin{experience}[L'origine de l'énergie hydroélectrique : les barrages de Lom Pangar et Nachtigal]
Au Cameroun, l'essentiel de l'électricité provient de l'hydroélectricité (barrages d'\textbf{Edéa}, de \textbf{Lom Pangar}, de \textbf{Nachtigal}, de \textbf{Song Loulou}). L'eau retenue en amont du barrage possède une énergie du fait de sa \textbf{position en hauteur}. En tombant dans la galerie de conduite forcée, cette énergie se transforme en énergie cinétique qui fait tourner les turbines des alternateurs (ENEO/SONATREL). Cette énergie « de position » est un exemple paradigmatique d'\textbf{énergie potentielle}.
\end{experience}

\begin{definition}[Force conservative et énergie potentielle]
Une force $\vec{F}$ est dite \cle{conservative} si son travail ne dépend que des positions initiale et finale du système, et non du chemin suivi. Équivalentement, le travail d'une force conservative sur un parcours \textbf{fermé} est nul : $\oint \vec{F} \cdot d\vec{l} = 0$.
À toute force conservative, on peut associer une grandeur scalaire, l'\textbf{énergie potentielle} $E_p$, définie par la relation :
\[
W_{A \to B}(\vec{F}_{\text{cons}}) = -\Delta E_p = -(E_{p,B} - E_{p,A})
\]
L'énergie potentielle est définie \textbf{à une constante additive près} (choix de l'origine / référence nulle).
\end{definition}

\begin{aretenir}[Principales énergies potentielles en Première]
\begin{itemize}
    \item \textbf{Énergie potentielle de pesanteur} ($E_{p,p}$) : associée au poids $\vec{P} = m\vec{g}$ (champ de pesanteur uniforme).
    \item \textbf{Énergie potentielle élastique} ($E_{p,e}$) : associée à la force de rappel d'un ressort linéaire $\vec{F}_e = -k\vec{x}$ (loi de Hooke).
    \item \textit{(Culture générale) Énergie potentielle électrostatique} : associée à la force de Coulomb entre charges.
    \item \textit{(Culture générale) Énergie potentielle chimique} : énergie de liaison dans les molécules (piles, carburants, biomasse).
\end{itemize}
\end{aretenir}

\courssub{Énergie potentielle de pesanteur}

\paragraph{Modélisation.}
On considère un système de masse $m$ (point matériel ou solide en translation) au voisinage de la Terre. Le champ de pesanteur est supposé uniforme : $\vec{g}$ constant, vertical, orienté vers le bas. On choisit un repère $(O, \vec{x}, \vec{y}, \vec{z})$ avec $\vec{z}$ vertical \textbf{orienté vers le haut}. L'altitude est $z$.

\paragraph{Travail du poids.}
Pour un déplacement de $A(z_A)$ à $B(z_B)$ :
\[
W_{A \to B}(\vec{P}) = \int_A^B m\vec{g} \cdot d\vec{l} = \int_{z_A}^{z_B} (-mg)\, dz = -mg(z_B - z_A) = -mg\Delta z
\]

\paragraph{Définition de l'énergie potentielle de pesanteur.}
D'après $W = -\Delta E_{p,p}$, on pose :
\[
\boxed{E_{p,p}(z) = mgz + Cte}
\]
On choisit généralement l'origine des énergies potentielles ($E_{p,p}=0$) pour $z=0$ (niveau de la mer, sol, fond de cuvette...), ce qui donne $Cte = 0$ :
\[
\boxed{E_{p,p} = mgz}
\]

\begin{propriete}[Variation de l'énergie potentielle de pesanteur]
\[
\boxed{\Delta E_{p,p} = mg\Delta z = mg(z_f - z_i)}
\]
\begin{itemize}
    \item Si le corps \textbf{monte} ($\Delta z > 0$) : $\Delta E_{p,p} > 0$, le travail du poids est négatif ($W_P < 0$), l'énergie cinétique diminue.
    \item Si le corps \textbf{descend} ($\Delta z < 0$) : $\Delta E_{p,p} < 0$, le travail du poids est positif ($W_P > 0$), l'énergie cinétique augmente.
\end{itemize}
\emph{Analyse dimensionnelle :} $[E_{p,p}] = [M][L][T]^{-2}[L] = [M][L]^2[T]^{-2} = \text{Joule (J)}$.
\end{propriete}

\begin{exemplebox}[Centrale de Nachtigal : de l'énergie potentielle à l'électricité]
L'eau du fleuve Sanaga arrive au barrage de Nachtigal (hauteur de chute nette $H \approx 35\,\text{m}$). Un débit $Q = 100\,\text{m}^3\!\cdot\!\text{s}^{-1}$ (masse volumique $\rho = 1000\,\text{kg}\!\cdot\!\text{m}^{-3}$) traverse les turbines.
Puissance potentielle disponible : $P = \dot{m}gH = (\rho Q)gH = 1000 \times 100 \times 9,8 \times 35 \approx 34,3\,\text{MW}$.
Après rendement turbines/alternateurs ($\eta \approx 0,9$), puissance électrique $\approx 31\,\text{MW}$ injectée sur le réseau interconnecté Sud (Yaoundé, Douala, Edéa).
\end{exemplebox}

\courssub{Énergie potentielle élastique}

\paragraph{Modélisation : ressort à réponse linéaire.}
Un ressort de raideur $k$ (en $\text{N}\!\cdot\!\text{m}^{-1}$), de longueur au repos $L_0$, est fixé à un mur. Un bloc de masse $m$ est attaché à l'autre extrémité. On repère la position du bloc par l'abscisse $x$ (origine $x=0$ au repos du ressort, axe orienté dans le sens de l'allongement). La force de rappel exercée par le ressort sur le bloc est :
\[
\vec{F}_e = -k x\,\vec{u}_x \quad \text{(Loi de Hooke, valide pour petites déformations)}
\]

\paragraph{Travail de la force élastique.}
\[
W_{A \to B}(\vec{F}_e) = \int_{x_A}^{x_B} (-kx)\, dx = -\frac{1}{2}k(x_B^2 - x_A^2)
\]

\paragraph{Définition de l'énergie potentielle élastique.}
D'après $W = -\Delta E_{p,e}$, on définit :
\[
\boxed{E_{p,e}(x) = \frac{1}{2}k x^2}
\]
L'origine ($E_{p,e}=0$) est choisie pour $x=0$ (ressort au repos, non déformé).

\begin{propriete}[Variation de l'énergie potentielle élastique]
\[
\boxed{\Delta E_{p,e} = \frac{1}{2}k(x_f^2 - x_i^2)}
\]
\begin{itemize}
    \item L'énergie potentielle élastique est \textbf{toujours positive ou nulle} ($E_{p,e} \ge 0$).
    \item Elle est maximale pour l'allongement (ou compression) maximal.
    \item \emph{Analyse dimensionnelle :} $[k] = \text{N}\!\cdot\!\text{m}^{-1} = [M][T]^{-2}$, $[E_{p,e}] = [M][T]^{-2}[L]^2 = \text{J}$.
\end{itemize}
\end{propriete}

\begin{attention}[Conditions d'application de $E_{p,e} = \frac{1}{2}kx^2$]
\begin{itemize}
    \item Ressort \textbf{idéal} (masse négligeable, pas de frottements internes).
    \item Réponse \textbf{linéaire} (loi de Hooke respectée) : $k$ constant. Au-delà de la limite d'élasticité, la formule n'est plus valable (déformation plastique, hystérésis).
    \item $x$ est l'allongement \textbf{algébrique} (positif si étiré, négatif si compressé) ; $x^2$ rend l'énergie positive dans les deux cas.
\end{itemize}
\end{attention}

\begin{experience}[Mesure de l'énergie potentielle élastique (TP possible)]
\textbf{Matériel :} Ressort ($k$ connu), masse $m$, règle graduée, support.
\textbf{Protocole :} On suspend la masse, on mesure l'allongement $x$ à l'équilibre ($mg = kx$). On tire la masse vers le bas d'un $x_0$ supplémentaire et on lâche sans vitesse initiale. On mesure la hauteur max $h$ atteinte au-dessus de l'équilibre.
\textbf{Interprétation :} $\frac{1}{2}k x_0^2 = mgh$ (si frottements négligeables). Permet de vérifier la conservation de $E_m$ et de valider $E_{p,e} = \frac{1}{2}kx^2$.
\end{experience}

\courssub{Énergie mécanique et sa conservation}

\begin{definition}[Énergie mécanique d'un système]
L'énergie mécanique $E_m$ d'un système (point matériel ou solide en translation) est la somme de son énergie cinétique $E_c$ et de ses énergies potentielles (pesanteur $E_{p,p}$ et élastique $E_{p,e}$ le cas échéant) :
\[
\boxed{E_m = E_c + E_{p,p} + E_{p,e} = \frac{1}{2}mv^2 + mgz + \frac{1}{2}kx^2}
\]
\end{definition}

\paragraph{Démonstration du principe de conservation (cas idéal).}
On applique le Théorème de l'Énergie Cinétique (TEC) : $\Delta E_c = \sum W_{\text{ext}}$.
On sépare les forces extérieures en conservatrices ($\vec{P}, \vec{F}_e$) et non conservatrices ($\vec{f}, \vec{F}_{\text{moteur}}, \dots$) :
\[
\Delta E_c = W(\vec{P}) + W(\vec{F}_e) + \sum W(\vec{F}_{\text{nc}})
\]
Par définition des énergies potentielles : $W(\vec{P}) = -\Delta E_{p,p}$ et $W(\vec{F}_e) = -\Delta E_{p,e}$.
\[
\Delta E_c = -\Delta E_{p,p} - \Delta E_{p,e} + \sum W(\vec{F}_{\text{nc}})
\]
En regroupant : $\Delta E_c + \Delta E_{p,p} + \Delta E_{p,e} = \sum W(\vec{F}_{\text{nc}})$, soit :
\[
\boxed{\Delta E_m = \sum W(\vec{F}_{\text{nc}})}
\]

\begin{propriete}[Principe de conservation de l'énergie mécanique — Exigible au Probatoire]
Si le système ne subit que des forces \textbf{conservatrices} (ou si la somme des travaux des forces non conservatrices est nulle), alors l'énergie mécanique est \textbf{constante} :
\[
\boxed{E_{m,i} = E_{m,f} \quad \text{soit} \quad E_c + E_{p,p} + E_{p,e} = \text{cte}}
\]
\emph{Justification :} Issue directe du TEC et de la définition $W_{\text{cons}} = -\Delta E_p$.
\end{propriete}

\begin{aretenir}[Situations types de conservation de $E_m$ (modèles idéaux)]
\begin{itemize}
    \item Chute libre sans frottements de l'air.
    \item Glissement sur plan incliné \textbf{sans frottement} (pente de glace, rail lisse).
    \item Oscillateur masse-ressort horizontal \textbf{sans frottement}.
    \item Pendule simple \textbf{sans frottement} (fil inextensible, masse du fil négligeable, pas de frottement d'air).
    \item Mouvement d'un satellite sur orbite (force centrale conservative).
\end{itemize}
\end{aretenir}

\begin{methode}[Résolution d'un problème par conservation de l'énergie mécanique]
\begin{enumerate}
    \item \textbf{Système / Référentiel :} Définir le système étudié (souvent un point matériel), référentiel galiléen (terrestre).
    \item \textbf{Bilan des forces / Schéma :} Dessiner les forces. Vérifier l'absence de forces non conservatrices travailleuses (frottements, choc, moteur).
    \item \textbf{Choix des origines :} Fixer $z=0$ pour $E_{p,p}$ et $x=0$ pour $E_{p,e}$. Noter les positions/états initial (i) et final (f).
    \item \textbf{Écrire $E_{m,i} = E_{m,f}$ :} $\frac{1}{2}mv_i^2 + mgz_i + \frac{1}{2}kx_i^2 = \frac{1}{2}mv_f^2 + mgz_f + \frac{1}{2}kx_f^2$.
    \item \textbf{Résoudre} l'inconnue (vitesse, hauteur, raideur, allongement...).
    \item \textbf{Vérifier} : Homogénéité, ordre de grandeur, signe des grandeurs.
\end{enumerate}
\end{methode}

\begin{exoresolu}[Pendule balistique : vitesse d'une balle de chasse]
\textbf{Énoncé :} Une balle de masse $m = 10\,\text{g}$ est tirée horizontalement dans un bloc de bois (masse $M = 2,0\,\text{kg}$) suspendu par un fil inextensible de longueur $L = 1,5\,\text{m}$. Le bloc (avec la balle incrustée) s'élève jusqu'à une hauteur $h = 12\,\text{cm}$ au-dessus de sa position d'équilibre. On néglige les frottements d'air et du pivot. Déterminer la vitesse $v_0$ de la balle avant l'impact. $g = 9,8\,\text{N}\!\cdot\!\text{kg}^{-1}$.
\tcblower
\textbf{Solution :}
\textbf{1. Choc (inélastique) :} Conservation de la quantité de mouvement (forces intérieures $\gg$ forces extérieures pendant le choc bref).
$m v_0 = (M+m) V \Rightarrow V = \frac{m}{M+m}v_0$.
\textbf{2. Montée du pendule (après choc) :} Forces : Poids (conservative), Tension (travail nul). $\Rightarrow$ Conservation de $E_m$.
État A (juste après choc, bas) : $E_{m,A} = \frac{1}{2}(M+m)V^2$ ($z=0$).
État B (point haut, $v=0$) : $E_{m,B} = (M+m)gh$.
$\frac{1}{2}(M+m)V^2 = (M+m)gh \Rightarrow V = \sqrt{2gh} = \sqrt{2 \times 9,8 \times 0,12} \approx 1,53\,\text{m}\!\cdot\!\text{s}^{-1}$.
\textbf{3. Vitesse initiale :} $v_0 = \frac{M+m}{m} V = \frac{2,01}{0,010} \times 1,53 \approx \boxed{308\,\text{m}\!\cdot\!\text{s}^{-1}}$.
\end{exoresolu}

\coursec{Non conservation de l'énergie mécanique}

Dans la section précédente, nous avons découvert une situation idéale : lorsque seules des forces conservatrices (pesanteur, force élastique) travaillent, l'énergie mécanique $E_m$ d'un système reste constante. C'est une propriété puissante, mais elle ne s'applique qu'à des systèmes isolés des frottements ou des chocs. Or, dans la vie quotidienne — que ce soit un taxi qui freine sur le boulevard du 20 Mai à Yaoundé, un mototaxi qui dérape sur une piste boueuse en saison des pluies, ou un ballon de football qui rebondit sur le terrain — l'énergie mécanique \emph{ne se conserve pas}. Elle diminue. Où va-t-elle ? C'est l'objet de cette section.

\courssub{Forces non conservatrices et variation de l'énergie mécanique}

\paragraph{Constat expérimental.}
Prenons un galet que l'on lance sur une table horizontale rugueuse. Il ralentit et s'arrête. Son énergie cinétique initiale $E_{c,i} = \frac{1}{2}mv_i^2$ devient nulle. Aucune variation d'énergie potentielle de pesanteur (hauteur constante) ni d'énergie potentielle élastique (pas de ressort) n'a lieu. L'énergie mécanique finale $E_{m,f}$ est donc nulle : $\Delta E_m = E_{m,f} - E_{m,i} < 0$. L'énergie mécanique a \emph{diminué}. Elle n'a pas « disparu » : elle a été transformée en énergie thermique (échauffement de la table et du galet) et en énergie sonore (bruit du frottement).

\paragraph{Forces conservatrices vs non conservatrices.}
Rappelons la définition opératoire : une force est \cle{conservatrice} si son travail ne dépend que des positions initiale et finale, et non du chemin suivi. Le travail d'une force conservative sur un parcours fermé est nul. La pesanteur $\vec{P}$ et la force élastique $\vec{F}_e = -k\vec{x}$ sont conservatrices.
Une force est \cle{non conservative} si son travail dépend du chemin, ou si son travail sur un parcours fermé n'est pas nul. Les principales forces non conservatrices rencontrées en Première sont :
\begin{itemize}
    \item Les \textbf{forces de frottement} (solides ou fluides) : $\vec{f}$, toujours opposées à la vitesse, travail toujours négatif ($W_f < 0$).
    \item Les \textbf{forces de choc} (impact, collision) : forces de contact intenses et brèves qui déforment les corps de manière irréversible.
    \item Les \textbf{forces musculaires} ou \textbf{moteurs} : forces extérieures qui apportent de l'énergie au système.
    \item Certaines \textbf{réactions de liaison} : bien que souvent perpendiculaires au déplacement (travail nul), une réaction de liaison mobile (ex : un bloc sur un tapis roulant) peut effectuer un travail non nul.
\end{itemize}

\begin{popfigure}
\begin{tikzpicture}[scale=1.1, >=Stealth]
    % Route
    \fill[popmuted!30] (-4,-0.5) rectangle (4,0);
    \draw[thick, popdark] (-4,0) -- (4,0);
    % Lignes route
    \foreach \x in {-3.5,-2.5,-1.5,-0.5,0.5,1.5,2.5,3.5}
        \draw[white, thick] (\x, -0.05) -- (\x, 0.05);
    
    % Voiture (simplifiée)
    \def\carx{0}
    \draw[popblue, thick, fill=popblue!20, rounded corners=3pt] (\carx-1.2,0) -- (\carx-1.2,1) -- (\carx+1.2,1) -- (\carx+1.2,0.6) -- (\carx+0.8,0.6) -- (\carx+0.8,0.3) -- (\carx-0.8,0.3) -- (\carx-0.8,0.6) -- (\carx-1.2,0.6) -- cycle;
    % Roues
    \draw[popdark, thick, fill=popdark] (\carx-0.9,0) circle (0.25);
    \draw[popdark, thick, fill=popdark] (\carx+0.9,0) circle (0.25);
    % Flèche vitesse
    \draw[poporange, very thick, ->] (\carx+1.6, 0.8) -- (\carx+2.6, 0.8) node[right] {$\vec{v}$};
    % Force frottement
    \draw[poppink, very thick, ->] (\carx-1.6, 0.3) -- (\carx-2.6, 0.3) node[left] {$\vec{f}$};
    % Poids et Normale
    \draw[popgreen, very thick, ->] (\carx, 1) -- (\carx, -0.5) node[below] {$\vec{P}=m\vec{g}$};
    \draw[popgreen, very thick, ->] (\carx, -0.5) -- (\carx, 0) node[above, pos=0.5, right=2pt] {$\vec{N}$};
    
    % Bilan énergétique (encadré à droite)
    \begin{scope}[shift={(5.5,0.5)}]
        \node[draw=popblue, thick, fill=popcream, rounded corners, minimum width=5cm, minimum height=4.5cm] (box) {};
        \node[popblue, font=\bfseries] at (box.north) [yshift=-0.3cm] {Bilan énergétique du freinage};
        
        % Initiale
        \node[align=left, anchor=west] at (-2.2, 1.2) {\textbf{État initial :}};
        \draw[popblue, thick, fill=popblue!30] (-2, 0.8) rectangle (1.5, 1.1) node[midway, black] {$E_{m,i} = E_{c,i} = \frac{1}{2}mv_i^2$};
        \draw[poporange, ->, thick] (-2, 0.8) -- (-2.5, 0.2) node[left] {$W_{f} < 0$};
        
        % Frottement travail
        \node[align=left, anchor=west] at (-2.2, 0.0) {\textbf{Travail des frottements :}};
        \draw[poppink, thick, fill=poppink!20] (-2, -0.2) rectangle (1.5, 0.1) node[midway, black] {$W_f = -f \cdot d = -\Delta E_m$};
        
        % Thermique
        \node[align=left, anchor=west] at (-2.2, -1.0) {\textbf{Énergie dissipée (thermique) :}};
        \draw[poporange, thick, fill=poporange!30] (-2, -1.4) rectangle (1.5, -1.1) node[midway, black] {$Q = |W_f| = f \cdot d$};
        \draw[poporange, ->, thick] (1.5, -1.25) -- (2.2, -1.25) node[right] {Échauffement freins/pneus/route};
        
        % Finale
        \node[align=left, anchor=west] at (-2.2, -2.2) {\textbf{État final (arrêt) :}};
        \draw[popgreen, thick, fill=popgreen!30] (-2, -2.6) rectangle (1.5, -2.3) node[midway, black] {$E_{m,f} = 0$};
    \end{scope}
\end{tikzpicture}
\legende{Freinage d'un véhicule sur route horizontale : l'énergie cinétique est dissipée par le travail négatif des forces de frottement, se transformant en énergie thermique (échauffement).}
\end{popfigure}

\paragraph{Théorème fondamental : Variation de l'énergie mécanique.}
On démontre ce résultat à partir du théorème de l'énergie cinétique (TEC) appliqué au système.
Le TEC s'écrit : $\Delta E_c = W_{\text{ext}} = \sum W(\text{forces extérieures})$.
Séparons les forces extérieures en forces conservatrices ($\vec{F}_{\text{cons}}$ : pesanteur, élastique) et forces non conservatrices ($\vec{F}_{\text{nc}}$ : frottements, choc, moteur...).
\[
\Delta E_c = \sum W(\vec{F}_{\text{cons}}) + \sum W(\vec{F}_{\text{nc}})
\]
Par définition de l'énergie potentielle, pour une force conservative, $W(\vec{F}_{\text{cons}}) = -\Delta E_p$. En sommant sur toutes les forces conservatrices :
\[
\sum W(\vec{F}_{\text{cons}}) = -\Delta E_p
\]
On obtient alors :
\[
\Delta E_c = -\Delta E_p + \sum W(\vec{F}_{\text{nc}})
\]
Or $\Delta E_m = \Delta E_c + \Delta E_p$. En substituant :
\[
\boxed{\Delta E_m = \sum W(\vec{F}_{\text{nc}})}
\]

\begin{propriete}[Théorème de la variation de l'énergie mécanique — Admis au Probatoire]
Pour un système en translation (ou assimilé à un point matériel), la variation de son énergie mécanique entre deux instants $t_i$ et $t_f$ est égale à la somme des travaux des forces \textbf{non conservatrices} extérieures (et intérieures non conservatrices s'il y a lieu) :
\[
\boxed{\Delta E_m = E_{m,f} - E_{m,i} = \sum W(\text{forces non conservatrices})}
\]
\emph{Justification esquissée :} Issue directe du TEC et de la définition de l'énergie potentielle pour les forces conservatrices ($W_{\text{cons}} = -\Delta E_p$).
\end{propriete}

\begin{aretenir}[Conséquences immédiates]
\begin{itemize}
    \item Si $\sum W(\text{forces nc}) = 0$ (aucune force non conservative, ou travaux nuls) $\Rightarrow \Delta E_m = 0$ : \textbf{conservation de $E_m$} (cas de la section précédente).
    \item Si $\sum W(\text{forces nc}) < 0$ (frottements, choc inélastique) $\Rightarrow \Delta E_m < 0$ : \textbf{diminution de $E_m$} (dissipation).
    \item Si $\sum W(\text{forces nc}) > 0$ (moteur, force musculaire propulsive) $\Rightarrow \Delta E_m > 0$ : \textbf{augmentation de $E_m$} (apport d'énergie).
\end{itemize}
\end{aretenir}

\begin{attention}[Piège majeur : « Énergie perdue » $\neq$ « Énergie détruite »]
Il est \textbf{interdit} d'écrire « l'énergie est perdue » ou « l'énergie a disparu » dans une copie de Probatoire.
\begin{itemize}
    \item L'énergie mécanique $E_m$ \emph{diminue} (elle n'est plus sous forme mécanique macroscopique).
    \item L'énergie \textbf{totale} de l'Univers (mécanique + thermique + chimique + nucléaire...) se \textbf{conserve toujours} (Premier principe de la thermodynamique).
    \item On dit : « L'énergie mécanique est \textbf{dissipée} en énergie thermique (chaleur) et/ou sonore » ou « Elle est \textbf{convertie} en énergie interne des corps en contact ».
\end{itemize}
\emph{Exemple :} Dans les descentes de la côte de Buea, les taxis sentent le brûlé. Ce n'est pas l'énergie qui « part en fumée », c'est l'énergie mécanique du véhicule qui se transforme en énergie thermique dans les plaquettes et disques de frein (température $\uparrow$), puis se diffuse dans l'air.
\end{attention}

\courssub{Application 1 : Analyse énergétique d'un freinage}

C'est la situation type du programme : « analyse d'une situation de freinage ». Un véhicule de masse $m$ roule à la vitesse $v_i$ sur une route horizontale. Le conducteur freine ; les forces de frottement (freins + pneus/route) sont modélisées par une force moyenne constante $\vec{f}$, opposée au mouvement. Le véhicule s'arrête après une distance $d$.

\paragraph{Modélisation.}
Système : $\{ \text{Véhicule} \}$. Référentiel terrestre galiléen.
Forces extérieures : $\vec{P}$, $\vec{N}$ (travail nul car $\perp \vec{v}$), $\vec{f}$ (frottements, non conservative, travail $W_f = -f d$).
Énergie potentielle de pesanteur constante ($\Delta E_p = 0$).
Énergie mécanique $E_m = E_c = \frac{1}{2}mv^2$.

\paragraph{Application du théorème.}
\[
\Delta E_m = W_f \quad \Rightarrow \quad E_{c,f} - E_{c,i} = -f d
\]
À l'arrêt, $v_f = 0 \Rightarrow E_{c,f} = 0$.
\[
0 - \frac{1}{2}mv_i^2 = -f d \quad \Rightarrow \quad \boxed{f = \frac{m v_i^2}{2d}}
\]
Cette relation lie la \textbf{distance de freinage} $d$ à la \textbf{vitesse initiale} $v_i$ (au carré) et à la \textbf{force de frottement moyenne} $f$. Elle montre que doubler la vitesse \textbf{quadruple} la distance de freinage (si $f$ constant), un fait crucial pour la sécurité routière.

\begin{exoresolu}[Freinage d'un taxi à Douala]
\textbf{Énoncé :} Un taxi de masse $m = 1\,200\,\text{kg}$ roule à $v_i = 72\,\text{km}\!\cdot\!\text{h}^{-1}$ sur le pont Wouri. Le conducteur freine brutalement. Le véhicule s'immobilise après $d = 45\,\text{m}$. On néglige les frottements de l'air. On prend $g = 9,8\,\text{N}\!\cdot\!\text{kg}^{-1}$.
\begin{enumerate}
    \item Convertir la vitesse en $\text{m}\!\cdot\!\text{s}^{-1}$.
    \item Calculer l'énergie cinétique initiale $E_{c,i}$.
    \item En déduire le travail des forces de frottement $W_f$.
    \item Calculer la valeur moyenne $f$ de la force de frottement.
    \item Comparer $f$ au poids du véhicule $P$. Commenter.
\end{enumerate}
\tcblower
\textbf{Solution détaillée :}
\begin{enumerate}
    \item $v_i = 72\,\text{km}\!\cdot\!\text{h}^{-1} = \frac{72}{3,6} = 20,0\,\text{m}\!\cdot\!\text{s}^{-1}$.
    \item $E_{c,i} = \frac{1}{2} m v_i^2 = \frac{1}{2} \times 1200 \times (20,0)^2 = 600 \times 400 = 2,40 \times 10^5\,\text{J} = 240\,\text{kJ}$.
    \item $\Delta E_m = E_{m,f} - E_{m,i} = 0 - 2,40 \times 10^5 = -2,40 \times 10^5\,\text{J}$.
        D'après le théorème : $W_f = \Delta E_m = \boxed{-2,40 \times 10^5\,\text{J}}$. Le travail est négatif (force opposée au déplacement).
    \item $W_f = -f d \Rightarrow f = -\frac{W_f}{d} = \frac{2,40 \times 10^5}{45} \approx 5\,333\,\text{N} \approx \boxed{5,33\,\text{kN}}$.
    \item $P = mg = 1200 \times 9,8 = 11\,760\,\text{N} \approx 11,8\,\text{kN}$.
        Ratio : $\frac{f}{P} \approx \frac{5,33}{11,8} \approx 0,45$.
        \emph{Commentaire :} La force de frottement moyenne nécessaire pour arrêter le taxi en $45\,\text{m}$ représente environ $45\%$ du poids du véhicule. C'est une force considérable que les pneus et le système de freinage doivent fournir sans glisser (adhérence). Si la route est mouillée, le coefficient de frottement maximal diminue, $f_{\max}$ diminue, et la distance de freinage $d$ augmente d'autant pour la même vitesse initiale.
\end{enumerate}
\end{exoresolu}

\begin{popfigure}
\begin{tikzpicture}
    \begin{axis}[
        ybar, bar width=12pt,
        width=\linewidth, height=7cm,
        ylabel={Énergie (kJ)}, ymin=0, ymax=260,
        xtick={1,2,3,4}, xticklabels={$E_{c,i}$ (initiale), $E_{c,f}$ (finale), $|W_f|$ (dissipée), $Q$ (thermique)},
        ymajorgrids=true, grid style={popline},
        nodes near coords, nodes near coords align={vertical},
        every node near coord/.append style={font=\small, popdark},
        legend style={at={(0.5,-0.15)}, anchor=north, legend columns=-1, fill=popcream, draw=popblue},
        title={Bilan énergétique du freinage du taxi (exemple résolu)},
        title style={font=\bfseries, popblue},
        bar cycle list={
            {fill=popblue!70, draw=popblue},
            {fill=popgreen!70, draw=popgreen},
            {fill=poppink!70, draw=poppink, pattern=north east lines, pattern color=poppink},
            {fill=poporange!70, draw=poporange},
        }
    ]
    \addplot coordinates {(1,240)};
    \addplot coordinates {(2,0)};
    \addplot coordinates {(3,240)};
    \addplot coordinates {(4,240)};
    \legend{$E_c$, $E_c$, $|W_f|$, $Q$}
    \end{axis}
\end{tikzpicture}
\legende{Diagramme en barres : l'énergie cinétique initiale ($240\,\text{kJ}$) est entièrement convertie en travail des frottements ($|W_f| = 240\,\text{kJ}$) qui devient énergie thermique $Q$. L'énergie mécanique finale est nulle.}
\end{popfigure}

\courssub{Application 2 : Analyse énergétique d'un choc}

Un choc est une interaction brève et intense entre deux corps. Pendant le choc, les forces de contact sont très grandes. Si les corps se déforment de manière \textbf{irréversible} (choc inélastique, plastique), une partie de l'énergie mécanique initiale est transformée en énergie de déformation permanente, en chaleur (échauffement au contact) et en énergie sonore (bruit de l'impact). L'énergie mécanique \textbf{ne se conserve pas}.

\paragraph{Modélisation simplifiée.}
On considère souvent le choc comme instantané. On ne peut pas calculer le travail des forces de choc directement (force inconnue, déplacement inconnu). On utilise le théorème sous forme globale :
\[
\Delta E_m = E_{m,f} - E_{m,i} = W_{\text{choc}} < 0
\]
L'énergie mécanique « perdue » $\Delta E_{\text{diss}} = -\Delta E_m = E_{m,i} - E_{m,f} > 0$ correspond à l'énergie dissipée.

\begin{popfigure}
\begin{tikzpicture}
    \begin{axis}[
        ybar stacked, bar width=20pt,
        width=\linewidth, height=7cm,
        ylabel={Énergie (J)}, ymin=0, ymax=110,
        xtick={1,2}, xticklabels={Avant le choc, Après le choc (rebond)},
        ymajorgrids=true, grid style={popline},
        legend style={at={(0.5,-0.18)}, anchor=north, legend columns=-1, fill=popcream, draw=popblue},
        title={Évolution de l'énergie mécanique lors d'un choc inélastique (balle de tennis)},
        title style={font=\bfseries, popblue},
        bar cycle list={
            {fill=popblue!70, draw=popblue}, % Ec
            {fill=popgreen!70, draw=popgreen}, % Ep
            {fill=poporange!70, draw=poporange, pattern=north west lines, pattern color=poporange}, % Dissipée
        }
    ]
    % Avant : Ec=80, Ep=20, Diss=0
    \addplot coordinates {(1,80)};
    \addplot coordinates {(1,20)};
    \addplot coordinates {(1,0)};
    % Après : Ec=30, Ep=10, Diss=60 (total 100, mais Em = 40)
    \addplot coordinates {(2,30)};
    \addplot coordinates {(2,10)};
    \addplot coordinates {(2,60)};
    \legend{$E_c$, $E_p$, $E_{\text{dissipée}}$}
    \end{axis}
\end{tikzpicture}
\legende{Lors d'un choc inélastique (rebond d'une balle de tennis sur le sol), l'énergie mécanique $E_m = E_c + E_p$ diminue (barres bleues+vertes). La partie hachurée orange représente l'énergie dissipée en chaleur, déformation permanente et son.}
\end{popfigure}

\begin{exemplebox}[Coefficient de restitution et énergie]
Pour un choc unidimensionnel entre une balle (masse $m$) et un mur massif, on définit le coefficient de restitution $e = \frac{v_{\text{après}}}{v_{\text{avant}}}$ ($0 \le e \le 1$), rapport des vitesses relatives (normes) après et avant le choc.
L'énergie cinétique après le choc vaut $E_{c,f} = \frac{1}{2}m (e v_i)^2 = e^2 E_{c,i}$.
L'énergie dissipée est $\Delta E_{\text{diss}} = E_{c,i} - E_{c,f} = (1-e^2) E_{c,i}$.
\begin{itemize}
    \item $e=1$ : Choc \textbf{parfaitement élastique}, $\Delta E_{\text{diss}}=0$, $E_m$ conservée (modèle idéal, bille d'acier sur acier).
    \item $e=0$ : Choc \textbf{parfaitement inélastique}, corps soudés, $\Delta E_{\text{diss}} = E_{c,i}$ (perte max d'énergie cinétique compatible avec la conservation de la quantité de mouvement).
    \item $0<e<1$ : Choc \textbf{inélastique réel} (balle de tennis, $e \approx 0,75$ ; ballon de foot, $e \approx 0,6$). Une fraction $(1-e^2)$ de l'énergie cinétique initiale est dissipée.
\end{itemize}
\end{exemplebox}

\courssub{Méthode de résolution : Choisir la bonne approche}

Face à un problème de mécanique énergétique, la première étape est \textbf{l'inventaire des forces} et leur classification.

\begin{methode}[Analyse énergétique d'un système réel]
\begin{enumerate}
    \item \textbf{Choisir le système} et le référentiel (généralement terrestre, galiléen).
    \item \textbf{Dresser le bilan des forces} (schéma obligatoire). Identifier chaque force.
    \item \textbf{Classer les forces} :
        \begin{itemize}
            \item \textbf{Conservatrices} : Pesanteur $\vec{P}$ $\rightarrow$ $E_{p,p} = mgz$ ; Force élastique $\vec{F}_e$ $\rightarrow$ $E_{p,e} = \frac{1}{2}kx^2$. Leur travail s'écrit $-\Delta E_p$.
            \item \textbf{Non conservatrices (travail non nul)} : Frottements $\vec{f}$, forces de choc, moteur, réaction de liaison mobile. Leur travail $W_{\text{nc}}$ doit être calculé (souvent $W_f = -f d$ ou inconnu $\rightarrow$ c'est la question).
            \item \textbf{Non conservatrices (travail nul)} : Réaction normale $\vec{N}$ fixe, tension de fil inextensible (si point d'application fixe). Ne rentrent pas dans $\sum W_{\text{nc}}$.
        \end{itemize}
    \item \textbf{Écrire l'énergie mécanique} : $E_m = E_c + E_{p,p} + E_{p,e}$.
    \item \textbf{Appliquer le bon théorème} :
        \begin{itemize}
            \item Si $\sum W_{\text{nc}} = 0$ : \textbf{Conservation} $E_{m,i} = E_{m,f}$.
            \item Si $\sum W_{\text{nc}} \neq 0$ : \textbf{Variation} $\Delta E_m = \sum W_{\text{nc}}$.
        \end{itemize}
    \item \textbf{Résoudre l'équation} pour l'inconnue (vitesse, distance, coefficient de frottement, raideur $k$, énergie dissipée...).
    \item \textbf{Vérifier} : Homogénéité dimensionnelle, ordre de grandeur, signe des travaux (frottements $\rightarrow W<0$).
\end{enumerate}
\end{methode}

\begin{attention}[Erreurs fréquentes à ne pas commettre]
\begin{itemize}
    \item \textbf{Oublier une force non conservative} (ex: frottements de l'air, frottements dans l'axe d'une poulie) et appliquer la conservation de $E_m$ à tort. $\rightarrow$ Toujours vérifier : « Y a-t-il des frottements ? Un choc ? Un moteur ? »
    \item \textbf{Confondre travail de la pesanteur et variation d'énergie potentielle.} Rappel : $W(\vec{P}) = -\Delta E_{p,p}$. Ne pas ajouter $W(\vec{P})$ \textbf{ET} $\Delta E_{p,p}$ dans le bilan ! On utilise \textbf{soit} les travaux de toutes les forces (TEC pur), \textbf{soit} $E_m$ et travaux des forces non conservatrices seulement.
    \item \textbf{Calculer $W_f = +f d$ au lieu de $-f d$.} La force de frottement s'oppose au déplacement : $\vec{f} \cdot d\vec{l} = -f\,dl$. Le travail est \textbf{négatif}.
    \item \textbf{Ne pas convertir les unités} (km/h en m/s, cm en m, g en kg) avant de calculer les énergies en Joules.
\end{itemize}
\end{attention}

\courssub{Exemple complet : Descente sur pente avec frottements}

Un bloc de masse $m = 2,0\,\text{kg}$ glisse depuis le repos au point A sur une piste inclinée de $\alpha = 30^\circ$. La longueur de la pente est $L = 5,0\,\text{m}$. Le coefficient de frottement solide est $\mu = 0,20$. On prend $g = 9,8\,\text{N}\!\cdot\!\text{kg}^{-1}$. Déterminer la vitesse $v_B$ du bloc au point B (bas de la pente).

\begin{exoresolu}[Bloc sur pente rugueuse]
\textbf{Solution :}
\begin{enumerate}
    \item \textbf{Système / Référentiel :} $\{ \text{Bloc} \}$, référentiel terrestre.
    \item \textbf{Forces :} $\vec{P}=m\vec{g}$ (conservative), $\vec{N}$ (travail nul), $\vec{f}$ (frottement, non conservative, $f = \mu N = \mu mg\cos\alpha$).
    \item \textbf{Énergies :} Origine $E_{p,p}$ au point B ($z_B=0$).
        $z_A = L\sin\alpha = 5,0 \times \sin 30^\circ = 2,5\,\text{m}$.
        $E_{m,A} = E_{c,A} + E_{p,p,A} = 0 + mgz_A = 2,0 \times 9,8 \times 2,5 = 49,0\,\text{J}$.
        $E_{m,B} = E_{c,B} + E_{p,p,B} = \frac{1}{2}mv_B^2 + 0$.
    \item \textbf{Travail force non conservative (frottement) :}
        $W_f = -f L = -\mu mg\cos\alpha \times L$.
        $W_f = -0,20 \times 2,0 \times 9,8 \times \cos 30^\circ \times 5,0$.
        $\cos 30^\circ = \frac{\sqrt{3}}{2} \approx 0,866$.
        $W_f \approx -0,20 \times 19,6 \times 0,866 \times 5,0 \approx -16,97\,\text{J}$.
    \item \textbf{Théorème $\Delta E_m = W_f$ :}
        $E_{m,B} - E_{m,A} = W_f$
        $\frac{1}{2}mv_B^2 - 49,0 = -16,97$
        $\frac{1}{2} \times 2,0 \times v_B^2 = 49,0 - 16,97 = 32,03$
        $v_B^2 = 32,03 \Rightarrow v_B \approx \sqrt{32,03} \approx \boxed{5,66\,\text{m}\!\cdot\!\text{s}^{-1}}$.
    \item \textbf{Vérification :} Sans frottement ($\mu=0$), $v_B = \sqrt{2gz_A} = \sqrt{2\times9,8\times2,5} = 7,0\,\text{m}\!\cdot\!\text{s}^{-1}$. Avec frottements, $v_B$ est plus faible ($5,66 < 7,0$). Cohérent.
\end{enumerate}
\end{exoresolu}

\begin{savaistu}[Le freinage par récupération d'énergie : l'avenir au Cameroun ?]
Dans les véhicules thermiques classiques (taxis, cars de transport interurbain), l'énergie cinétique est \textbf{perdue} en chaleur lors du freinage (disques, tambours). C'est un gaspillage énergétique énorme.
Les \textbf{véhicules hybrides et électriques} (de plus en plus présents à Yaoundé et Douala) utilisent le \textbf{freinage régénératif} : le moteur électrique fonctionne en \textbf{générateur} lors du freinage. La force électromagnétique (force de Laplace) s'oppose au mouvement (force non conservative, travail négatif), mais au lieu de chauffer des plaquettes, elle \textbf{recharge la batterie} (conversion $E_c \rightarrow$ énergie électrique $\rightarrow$ énergie chimique). L'énergie mécanique n'est pas « perdue », elle est \textbf{récupérée} ! C'est une application directe du théorème $\Delta E_m = W_{\text{nc}}$, où $W_{\text{nc}}$ (moteur/générateur) est négatif mais l'énergie est stockée utilement.
\end{savaistu}

\courssub{Bilan de la leçon : Énergie potentielle et énergie mécanique}

\begin{aretenir}[Synthèse des savoirs et savoir-faire clés]
\begin{enumerate}
    \item \textbf{Énergie potentielle de pesanteur} (près de la Terre) : $E_{p,p} = mgz$ (choix de l'origine $z=0$ libre). Variation : $\Delta E_{p,p} = mg\Delta z$.
    \item \textbf{Énergie potentielle élastique} (ressort linéaire, raideur $k$) : $E_{p,e} = \frac{1}{2}k x^2$ (origine $x=0$ au repos). Variation : $\Delta E_{p,e} = \frac{1}{2}k(x_f^2 - x_i^2)$.
    \item \textbf{Énergie mécanique} : $E_m = E_c + E_{p,p} + E_{p,e}$.
    \item \textbf{Conservation de $E_m$} : Si \textbf{seules des forces conservatrices} travaillent (ou si $\sum W_{\text{nc}} = 0$), alors $E_{m,i} = E_{m,f} = \text{cte}$.
    \item \textbf{Variation de $E_m$ (cas général)} : $\Delta E_m = \sum W(\text{forces non conservatrices})$.
        \begin{itemize}
            \item Frottements : $W_f = -f d < 0 \Rightarrow \Delta E_m < 0$ (dissipation thermique).
            \item Choc inélastique : $W_{\text{choc}} < 0 \Rightarrow \Delta E_m < 0$ (déformation, chaleur, son).
            \item Moteur/Propulsion : $W_{\text{moteur}} > 0 \Rightarrow \Delta E_m > 0$ (apport d'énergie).
        \end{itemize}
    \item \textbf{Analyse dimensionnelle} : $[E] = [M][L]^2[T]^{-2} = \text{Joule (J)}$. $[k] = \text{N}\!\cdot\!\text{m}^{-1}$. $[\mu] = \text{sans dimension}$.
\end{enumerate}
\end{aretenir}

\begin{methode}[Check-list pour le Probatoire (Problème « Énergie »)]
\begin{enumerate}
    \item \textbf{Lire} : Identifier le système, le référentiel, les points A (initial) et B (final).
    \item \textbf{Dessiner} : Schéma des forces au point A et/ou B. Vecteurs $\vec{v}_A$, $\vec{v}_B$.
    \item \textbf{Lister} : Forces conservatrices ($E_p$ associée) / Forces non conservatrices ($W$ à calculer ou inconnu).
    \item \textbf{Choisir} : Conservation ($E_{m,A}=E_{m,B}$) \textbf{OU} Variation ($\Delta E_m = \sum W_{\text{nc}}$). \emph{Justifier le choix en une phrase.}
    \item \textbf{Écrire} : L'équation symbolique complète.
    \item \textbf{Calculer} : Conversions d'unités $\rightarrow$ Application numérique $\rightarrow$ Résultat encadré avec unité.
    \item \textbf{Commenter} : Sens du signe, ordre de grandeur, comparaison cas idéal/réel.
\end{enumerate}
\end{methode}

\begin{exercice}[Entraînement : Le saut à l'élastique au barrage de Lom Pangar]
Un sauteur de masse $m = 70\,\text{kg}$ saute du haut du barrage de Lom Pangar (hauteur $H = 50\,\text{m}$ au-dessus de l'eau). Il est attaché à un élastique de longueur naturelle $L_0 = 20\,\text{m}$ et de raideur $k = 150\,\text{N}\!\cdot\!\text{m}^{-1}$. On néglige la masse de l'élastique et les frottements de l'air. Le sauteur part du repos ($v=0$).
\begin{enumerate}
    \item Choisir l'origine de l'énergie potentielle de pesanteur. Exprimer $E_{m}$ en fonction de l'altitude $z$ (mesurée vers le haut depuis l'eau) et de l'allongement $x$ de l'élastique.
    \item Déterminer l'altitude $z_{\text{min}}$ la plus basse atteinte par le sauteur (élastique tendu au maximum).
    \item Calculer sa vitesse $v$ au moment où l'élastique commence à se tendre ($z = H - L_0$).
    \item \textbf{Question bonus (non conservation) :} En réalité, l'élastique s'échauffe (hystérésis) et l'air freine. Si le sauteur ne descend qu'à $z_{\text{réel}} = 15\,\text{m}$ au-dessus de l'eau, calculer l'énergie mécanique dissipée lors de la descente.
\end{enumerate}
\end{exercice}

\begin{corrige}[Exercice : Saut à l'élastique]
\begin{enumerate}
    \item Origine $E_{p,p}$ au niveau de l'eau ($z=0$).
        $E_c = \frac{1}{2}mv^2$.
        $E_{p,p} = mgz$.
        $E_{p,e} = \begin{cases} 0 & \text{si } z > H-L_0 \text{ (élastique détendu)} \\ \frac{1}{2}k(H - L_0 - z)^2 & \text{si } z \le H-L_0 \text{ (élastique tendu, } x = H-L_0-z \ge 0) \end{cases}$.
        $E_m = E_c + mgz + E_{p,e}$.
    \item Au point le plus bas ($z_{\text{min}}$), $v=0 \Rightarrow E_c=0$. Conservation de $E_m$ (pas de frottements) :
        $E_{m,\text{initial}}(z=H) = E_{m,\text{final}}(z=z_{\text{min}})$.
        $mgH = mgz_{\text{min}} + \frac{1}{2}k(H - L_0 - z_{\text{min}})^2$.
        $70 \times 9,8 \times 50 = 70 \times 9,8 \times z_{\text{min}} + \frac{1}{2} \times 150 \times (30 - z_{\text{min}})^2$.
        $34300 = 686 z_{\text{min}} + 75 (900 - 60 z_{\text{min}} + z_{\text{min}}^2)$.
        $75 z_{\text{min}}^2 + (686 - 4500) z_{\text{min}} + 67500 - 34300 = 0$.
        $75 z_{\text{min}}^2 - 3814 z_{\text{min}} + 33200 = 0$.
        $\Delta = 3814^2 - 4 \times 75 \times 33200 \approx 1,455 \times 10^7 - 9,96 \times 10^6 = 4,59 \times 10^6$.
        $\sqrt{\Delta} \approx 2142$.
        $z_{\text{min}} = \frac{3814 \pm 2142}{150}$.
        Deux solutions : $z_1 \approx 39,7\,\text{m}$ (point de départ, élastique détendu) et $z_2 \approx \frac{1672}{150} \approx \boxed{11,1\,\text{m}}$.
        Le sauteur s'arrête à $\approx 11\,\text{m}$ au-dessus de l'eau.
    \item Au début de la tension ($z_1 = H - L_0 = 30\,\text{m}$), $E_{p,e}=0$. Conservation $E_m$ entre $H$ et $30\,\text{m}$ :
        $mgH = \frac{1}{2}mv_1^2 + mgz_1 \Rightarrow v_1 = \sqrt{2g(H-z_1)} = \sqrt{2 \times 9,8 \times 20} \approx \boxed{19,8\,\text{m}\!\cdot\!\text{s}^{-1}}$.
    \item Énergie dissipée réelle : $\Delta E_{\text{diss}} = E_{m,\text{initial}} - E_{m,\text{final réel}}$.
        $E_{m,i} = mgH = 34300\,\text{J}$.
        $E_{m,f} = mgz_{\text{réel}} + \frac{1}{2}k(H-L_0-z_{\text{réel}})^2$ (car $v=0$ au plus bas).
        $z_{\text{réel}} = 15\,\text{m}$. Allongement $x = 30 - 15 = 15\,\text{m}$.
        $E_{m,f} = 70 \times 9,8 \times 15 + \frac{1}{2} \times 150 \times 15^2 = 10290 + 16875 = 27165\,\text{J}$.
        $\Delta E_{\text{diss}} = 34300 - 27165 = \boxed{7135\,\text{J}} \approx 7,1\,\text{kJ}$.
        Cette énergie a été convertie en chaleur dans l'élastique (hystérésis) et en travail contre les frottements de l'air.
\end{enumerate}
\end{corrige}

\coursec{Bilan et méthodes de résolution}

Après avoir analysé en détail les cas où l'énergie mécanique se conserve et ceux où elle ne se conserve pas, nous disposons maintenant de tous les outils pour aborder sereinement n'importe quel problème d'énergie mécanique au Probatoire. Cette section synthétise les formules clés, te propose un organigramme de décision infaillible, détaille les trois méthodes types qui reviennent systématiquement, et t'entraîne à lire les graphiques d'énergies — une compétence très prisée à l'examen.

\courssub{Synthèse des expressions d'énergie}

Avant de te lancer dans un calcul, assure-toi de maîtriser parfaitement les trois énergies en jeu. Le tableau ci-dessous regroupe les expressions, leurs conditions de validité et leurs unités SI. C'est ta « fiche de révision » à connaître par cœur.

\begin{center}
\begin{tabularx}{\linewidth}{|l|X|X|X|}\hline
\rowcolor{popblueL}
\textbf{Énergie} & \textbf{Expression} & \textbf{Conditions / Précisions} & \textbf{Unité SI} \\ \hline
Énergie cinétique $E_c$ & $E_c = \frac{1}{2} m v^2$ & $m$ en \SI{}{kg}, $v$ en \SI{}{m/s}. Toujours positive ou nulle. & \SI{}{J} \\ \hline
Énergie potentielle de pesanteur $E_{pp}$ & $E_{pp} = m g h$ & $h$ : altitude par rapport à un \textbf{référentiel d'altitude choisi et fixé} ($h=0$ au niveau de référence). $g \approx \SI{9,81}{N/kg}$ (ou \SI{10}{N/kg} si autorisé). & \SI{}{J} \\ \hline
Énergie potentielle élastique $E_{pe}$ & $E_{pe} = \frac{1}{2} k x^2$ & $k$ : raideur du ressort (\SI{}{N/m}). $x$ : allongement ou compression \textbf{par rapport à la position de repos} ($x=0$ au repos). Toujours positive ou nulle. & \SI{}{J} \\ \hline
\end{tabularx}
\end{center}

\begin{aretenir}[Énergie mécanique]
L'énergie mécanique d'un système est la somme de son énergie cinétique et de ses énergies potentielles :
\[
E_m = E_c + E_{pp} + E_{pe}
\]
Si plusieurs ressorts ou plusieurs masses sont en jeu, on somme toutes les contributions.
\end{aretenir}

\begin{attention}[Origine des altitudes — Piège n°1 au Probatoire]
L'énergie potentielle de pesanteur $E_{pp} = mgh$ dépend \textbf{entièrement} du choix de l'origine des altitudes ($h=0$).
\begin{itemize}
    \item \textbf{Règle d'or :} Tu choisis \textbf{une seule fois} ton niveau $h=0$ au début de l'exercice (souvent le point le plus bas de la trajectoire, ou le niveau du sol, ou la position d'équilibre du ressort) et tu \textbf{ne changes jamais} en cours de calcul.
    \item Si tu calcules une \textbf{variation} $\Delta E_{pp} = mg\Delta h$, le choix de l'origine n'importe plus, car $\Delta h = h_f - h_i$ est invariant. Mais si tu écris $E_{m,i} = E_{m,f}$ (conservation), tu dois utiliser les $h_i$ et $h_f$ mesurés par rapport à la \textbf{même} origine.
\end{itemize}
\end{attention}

\courssub{Organigramme de décision : Conservation ou non conservation ?}

Devant tout énoncé, ta première action est d'identifier les forces qui travaillent et leur nature. L'arbre de décision ci-dessous te guide pas à pas.

\begin{popfigure}
\begin{tikzpicture}[
    node distance=1.5cm and 2.5cm,
    decision/.style={diamond, draw=popblue, thick, fill=popblue!10, text width=3.5cm, align=center, inner sep=4pt},
    process/.style={rectangle, draw=popdark, thick, fill=popcream, text width=3.5cm, align=center, inner sep=4pt, rounded corners=4pt},
    start/.style={ellipse, draw=popgreen, thick, fill=popgreen!10, text width=2.5cm, align=center, inner sep=4pt},
    arrow/.style={-Stealth, thick, draw=popdark},
    yesno/.style={midway, fill=white, inner sep=2pt, font=\small}
]
\node[start] (start) {Début : Analyser le système};
\node[decision, below of=start, align=center] (q1){Y a-t-il des forces \textbf{non conservatrices} dont le travail n'est pas nul ?\\(Frottements, choc plastique, moteur, freinage...)};
\node[process, below left=1.5cm and 1cm of q1, align=center] (yes){\textbf{Cas NON CONSERVATIF}\\$\Delta E_m = \sum W_{\text{ext, nc}}$\\(Travail des forces non conservatrices)};
\node[process, below right=1.5cm and 1cm of q1, align=center] (no){\textbf{Cas CONSERVATIF}\\$\Delta E_m = 0 \iff E_{m,i} = E_{m,f}$\\(Seules forces conservatrices : Poids, Ressort)};
\node[process, below=2.5cm of yes, align=center] (calc1){Calculer $W = \vec{F} \cdot \vec{AB} = F \cdot AB \cos\theta$\\Souvent $W < 0$ (frottements)};
\node[process, below=2.5cm of no, align=center] (calc2){Écrire $E_{c,i} + E_{pp,i} + E_{pe,i} = E_{c,f} + E_{pp,f} + E_{pe,f}$\\Résoudre pour l'inconnue ($v$, $h$, $x$...)};

\draw[arrow] (start) -- (q1);
\draw[arrow] (q1) -- node[yesno, left] {OUI} (yes);
\draw[arrow] (q1) -- node[yesno, right] {NON} (no);
\draw[arrow] (yes) -- (calc1);
\draw[arrow] (no) -- (calc2);
\end{tikzpicture}
\legende{Arbre de décision : identifier la nature des forces pour choisir la bonne relation énergétique.}
\end{popfigure}

\begin{propriete}[Critère rapide]
Une force est \textbf{conservative} si son travail ne dépend que des positions initiale et finale (pas du chemin). Le poids $\vec{P}$ et la force de rappel d'un ressort idéal $\vec{F} = -k\vec{x}$ sont conservatives. Les frottements $\vec{f}$, les forces de freinage, les chocs plastiques sont \textbf{non conservatifs} (leur travail est négatif et dépend du chemin).
\end{propriete}

\courssub{Méthode type 1 : Conservation de l'énergie mécanique ($E_{m,i} = E_{m,f}$)}

C'est la méthode « reine » : chute libre, plan incliné sans frottement, pendule simple, oscillateur masse-ressort horizontal sans frottement. L'inconnue est souvent une vitesse, une hauteur ou une compression maximale.

\begin{methode}[Résolution type « Conservation de $E_m$ » — 5 étapes]
\begin{enumerate}
    \item \textbf{Choisir le système} et lister les forces extérieures. Vérifier qu'aucune force non conservative ne travaille (pas de frottement, pas de choc).
    \item \textbf{Choisir l'origine des altitudes ($h=0$)} et la position de repos du ressort ($x=0$). Les fixer \textbf{définitivement}.
    \item \textbf{Identifier deux états} : Initial (i) et Final (f). Noter pour chaque état : $v$, $h$, $x$ (connus ou inconnus).
    \item \textbf{Écrire la conservation} : $E_{c,i} + E_{pp,i} + E_{pe,i} = E_{c,f} + E_{pp,f} + E_{pe,f}$.
    \item \textbf{Résoudre l'équation} pour l'inconnue. Vérifier l'homogénéité (J = J) et le sens physique ($v^2 \ge 0$, $x^2 \ge 0$).
\end{enumerate}
\end{methode}

\begin{exemplebox}[Bille sur plan incliné puis ressort horizontal]
Une bille de masse $m = \SI{50}{g} = \SI{0,050}{kg}$ est lâchée sans vitesse initiale au point A, à une hauteur $h_A = \SI{0,80}{m}$ au-dessus d'un plan horizontal. Elle glisse sans frottement sur un plan incliné, puis sur le plan horizontal où elle rencontre un ressort de raideur $k = \SI{200}{N/m}$ fixé à un mur. On cherche la compression maximale $x_{\text{max}}$ du ressort.
\begin{center}
\begin{tikzpicture}[scale=0.8, >=Stealth]
\draw[thick, popdark] (0,0) -- (6,0) node[midway, below=2mm] {Plan horizontal};
\draw[thick, popdark] (0,0) -- (-2,2) node[midway, left=2mm] {Plan incliné};
\draw[fill=poporange] (-2,2) circle (4pt) node[above left] {A ($h=\SI{0,80}{m}$, $v=0$)};
\draw[thick, decorate, decoration={coil, aspect=0.3, segment length=3mm, amplitude=2mm}] (4,0) -- (6,0) node[midway, above=2mm] {Ressort $k$};
\draw[fill=popblue] (4,0) circle (4pt) node[below] {Contact ressort};
\draw[<->] (-2,2) -- (-2,0) node[midway, left] {$h_A$};
\draw[<->] (4,0) -- (6,0) node[midway, above] {$x_{\text{max}}$};
\draw[->, thick, popgreen] (-1.5, 2.5) -- (-0.5, 1.5) node[midway, above right] {Trajectoire};
\end{tikzpicture}
\end{center}
\textbf{Solution :}
\begin{enumerate}
    \item Système = \{Bille\}. Forces : Poids (conservative), Réaction du plan (perpendiculaire au déplacement $\to W=0$), Force du ressort (conservative). \cle{Conservation de $E_m$ applicable}.
    \item Origine $h=0$ : Plan horizontal. Repos ressort $x=0$ : position contact bille/ressort.
    \item État initial (A) : $v_i=0$, $h_i=\SI{0,80}{m}$, $x_i=0$ (ressort non touché). $\Rightarrow E_{m,i} = 0 + mgh_A + 0$.
    \item État final (B, compression max) : $v_f=0$ (arrêt momentané), $h_f=0$, $x_f=x_{\text{max}}$. $\Rightarrow E_{m,f} = 0 + 0 + \frac{1}{2}k x_{\text{max}}^2$.
    \item $mgh_A = \frac{1}{2}k x_{\text{max}}^2 \implies x_{\text{max}} = \sqrt{\frac{2mgh_A}{k}} = \sqrt{\frac{2 \times 0,050 \times 9,81 \times 0,80}{200}} = \sqrt{0,003924} \approx \SI{0,0626}{m} = \textbf{\SI{6,3}{cm}}$.
\end{enumerate}
\end{exemplebox}

\courssub{Méthode type 2 : Non conservation — $\Delta E_m = W_{\text{nc}}$}

Ici, des frottements (sol, air), un freinage, ou un choc inélastique dissipent de l'énergie. L'inconnue est souvent un coefficient de frottement, une distance de freinage, ou une vitesse après un passage rugueux.

\begin{methode}[Résolution type « $\Delta E_m = W_{\text{frott}}$ » — 5 étapes]
\begin{enumerate}
    \item \textbf{Identifier la force non conservative} (frottement $\vec{f}$, force de freinage $\vec{F}_f$...). Déterminer son expression ($f = \mu N = \mu mg\cos\alpha$ sur plan incliné).
    \item \textbf{Calculer son travail} $W_{\text{nc}} = \vec{F}_{\text{nc}} \cdot \vec{AB} = F_{\text{nc}} \cdot AB \cos\theta$. \textbf{Attention au signe :} $\theta = 180^\circ \Rightarrow \cos\theta = -1 \Rightarrow W_{\text{nc}} < 0$.
    \item \textbf{Choisir origine $h=0$ et $x=0$} (identique aux deux états).
    \item \textbf{Écrire le théorème de l'énergie} : $E_{m,f} - E_{m,i} = W_{\text{nc}}$.
    \item \textbf{Résoudre} pour l'inconnue ($\mu$, $d$, $v_f$...).
\end{enumerate}
\end{methode}

\begin{exoresolu}[Freinage d'un mototaxi à Douala]
Un mototaxi de masse totale $m = \SI{180}{kg}$ (conducteur + passager + moto) roule à $v_i = \SI{72}{km/h} = \SI{20}{m/s}$ sur une route plane et horizontale. Le conducteur freine brutalement ; les roues se bloquent. Le coefficient de frottement cinétique pneus/bitume est $\mu = 0,60$. On néglige la résistance de l'air. Calculer la distance de freinage $d$.
\tcblower
\textbf{Solution :}
\begin{enumerate}
    \item Système = \{Mototaxi\}. Forces : Poids $\vec{P}$, Réaction $\vec{R}$ ($N=mg$), Frottement $\vec{f}$ (non conservative, oppose au mouvement). $f = \mu N = \mu mg$.
    \item Travail du frottement sur distance $d$ : $\vec{f} \perp \vec{P}, \vec{R}$ ; $\vec{f}$ antiparallèle à $\vec{d}$. $W_f = -f d = -\mu mg d$.
    \item Origine $h=0$ : niveau de la route (inutile car $h$ constant). États : Initial (début freinage, $v_i=\SI{20}{m/s}$), Final (arrêt, $v_f=0$).
    \item $\Delta E_m = E_{m,f} - E_{m,i} = \frac{1}{2}m v_f^2 - \frac{1}{2}m v_i^2 = 0 - \frac{1}{2}m v_i^2$.
    \item Relation : $-\frac{1}{2}m v_i^2 = -\mu mg d \implies d = \frac{v_i^2}{2\mu g} = \frac{20^2}{2 \times 0,60 \times 9,81} = \frac{400}{11,772} \approx \textbf{\SI{34,0}{m}}$.
\end{enumerate}
\textbf{Analyse dimensionnelle :} $[v^2] = \SI{}{m^2/s^2}$, $[\mu g] = \SI{}{m/s^2}$ $\Rightarrow$ $[d] = \SI{}{m}$. \textbf{Correct.}
\end{exoresolu}

\courssub{Méthode type 3 : Couplage Plan incliné + Ressort (Transformation $E_{pp} \leftrightarrow E_{pe}$)}

C'est un grand classique du Probatoire C : un mobile glisse sur un plan incliné (avec ou sans frottement) et vient comprimer un ressort placé en bas. L'énergie potentielle de pesanteur se transforme en énergie potentielle élastique (et éventuellement en travail de frottement).

\begin{exemplebox}[Bloc sur plan incliné rugueux avec ressort]
Un bloc de masse $m = \SI{2,0}{kg}$ est lâché depuis le repos au point A, situé à une distance $L = \SI{1,50}{m}$ (mesurée le long du plan) d'un ressort de raideur $k = \SI{500}{N/m}$ non comprimé. Le plan est incliné à $\alpha = 30^\circ$. Le coefficient de frottement est $\mu = 0,15$. On cherche la compression maximale $x_{\text{max}}$ du ressort.
\begin{center}
\begin{tikzpicture}[scale=0.9, >=Stealth]
\draw[thick, popdark] (0,0) -- (5,0) node[midway, below] {Plan incliné $\alpha=30^\circ$};
\draw[thick, decorate, decoration={coil, aspect=0.3, segment length=3mm, amplitude=2mm}] (5,0) -- (6.5,0) node[midway, above] {Ressort $k$};
\draw[fill=poporange] (0,0) circle (4pt) node[above left] {A (Départ)};
\draw[fill=popblue] (5,0) circle (4pt) node[below] {Contact ressort};
\draw[<->] (0,0) -- (5,0) node[midway, above=2mm, sloped] {$L$};
\draw[<->] (5,0) -- (6.5,0) node[midway, above] {$x_{\text{max}}$};
\draw[->, thick, popgreen] (0,0.5) -- (4,0.5) node[midway, above] {Glissement};
% Forces at a generic point
\coordinate (M) at (2,0);
\draw[->, thick] (M) -- ++(-30:1) node[midway, right] {$\vec{P}=m\vec{g}$};
\draw[->, thick] (M) -- ++(60:0.7) node[midway, above left] {$\vec{R}$};
\draw[->, thick] (M) -- ++(210:0.7) node[midway, below left] {$\vec{f}$};
\end{tikzpicture}
\end{center}
\textbf{Résolution guidée :}
\begin{enumerate}
    \item \textbf{Forces non conservatrices :} Frottements $\vec{f}$ sur la distance totale $L + x_{\text{max}}$. Travail $W_f = -f(L+x_{\text{max}}) = -\mu mg\cos\alpha (L+x_{\text{max}})$.
    \item \textbf{Origines :} $h=0$ au niveau du point de contact initial avec le ressort (bas du plan incliné). $x=0$ ressort au repos.
    \item \textbf{États :}
    \begin{itemize}
        \item Initial (A) : $v_i=0$, $h_i = L\sin\alpha$, $x_i=0$. $E_{m,i} = mgL\sin\alpha$.
        \item Final (B, compression max) : $v_f=0$, $h_f = -x_{\text{max}}\sin\alpha$ (le bloc descend encore de $x_{\text{max}}\sin\alpha$ en comprimant le ressort), $x_f=x_{\text{max}}$. $E_{m,f} = -mg x_{\text{max}}\sin\alpha + \frac{1}{2}k x_{\text{max}}^2$.
    \end{itemize}
    \item \textbf{Équation :} $E_{m,f} - E_{m,i} = W_f$
    \[
    \left(-mg x_{\text{max}}\sin\alpha + \frac{1}{2}k x_{\text{max}}^2\right) - mgL\sin\alpha = -\mu mg\cos\alpha (L+x_{\text{max}})
    \]
    \item \textbf{Résolution :} On regroupe les termes en $x_{\text{max}}^2$, $x_{\text{max}}$ et constant.
    \[
    \frac{1}{2}k x_{\text{max}}^2 + mg(\mu\cos\alpha - \sin\alpha)x_{\text{max}} - mgL(\sin\alpha - \mu\cos\alpha) = 0
    \]
    C'est une équation du second degré $ax^2+bx+c=0$ avec $a>0, c<0$ (donc une racine positive).
    Numériquement : 
    \begin{align*}
    a &= \frac{k}{2} = 250 \\
    b &= mg(\mu\cos\alpha - \sin\alpha) = 2,0 \times 9,81 \times (0,15 \times 0,866 - 0,5) \approx 19,62 \times (-0,370) \approx -7,26 \\
    c &= -mgL(\sin\alpha - \mu\cos\alpha) = -2,0 \times 9,81 \times 1,50 \times (0,5 - 0,130) \approx -29,43 \times 0,370 \approx -10,89
    \end{align*}
    $\Delta = b^2 - 4ac \approx 52,7 + 10890 \approx 10943 > 0$.
    $x_{\text{max}} = \frac{-b + \sqrt{\Delta}}{2a} \approx \frac{7,26 + 104,6}{500} \approx \textbf{0,224 m} = \textbf{\SI{22,4}{cm}}$.
\end{enumerate}
\end{exemplebox}

\begin{attention}[Erreur fréquente : Altitude du point final sur le ressort]
Dans l'exemple ci-dessus, au point final (compression max), le bloc est \textbf{plus bas} que le niveau de contact initial avec le ressort. Son altitude est $h_f = -x_{\text{max}}\sin\alpha$ (négative par rapport à l'origine choisie). Oublier ce terme $-mg x_{\text{max}}\sin\alpha$ dans $E_{m,f}$ est une erreur majeure qui fausse l'équation du second degré. \textbf{Dessine toujours un schéma avec les altitudes $h_i$ et $h_f$ repérées.}
\end{attention}

\courssub{Complément Probatoire : Lecture de graphes $E_c(t)$, $E_{pp}(t)$, $E_m(t)$}

À l'examen, on te donne souvent les courbes temporelles des énergies pour un oscillateur (masse-ressort horizontal ou vertical, pendule). Voici comment les décrypter.

\begin{propriete}[Oscillateur non amorti (idéal)]
\begin{itemize}
    \item $E_m(t)$ : \textbf{Constante} (ligne horizontale). Preuve visuelle de la conservation.
    \item $E_c(t)$ et $E_{pp}(t)$ ou $E_{pe}(t)$ : \textbf{Courbes sinusoïdales} (ou $\cos^2/\sin^2$) \textbf{en opposition de phase}. Quand l'une est max, l'autre est min (nulle).
    \item Période des énergies : $T_E = T_0/2$ (fréquence doublée) car $E_c \propto v^2 \propto \cos^2(\omega t)$.
    \item Somme instantanée : $E_c(t) + E_p(t) = E_m$ (constant) à tout instant.
\end{itemize}
\end{propriete}

\begin{popfigure}
\begin{tikzpicture}
\begin{axis}[
    width=\linewidth, height=7cm,
    axis lines=middle,
    xlabel={$t\ (\si{s})$}, ylabel={Énergie (\si{J})},
    xmin=0, xmax=4, ymin=-0.2, ymax=2.2,
    xtick={0,1,2,3,4}, ytick={0,0.5,1,1.5,2},
    domain=0:4, samples=200,
    legend style={at={(0.5,1.02)}, anchor=south, legend columns=3, font=\small},
    grid=major, grid style={popline},
]
% Em = 1.5 J constant
\addplot[popcurve, popdark, very thick] {1.5};
\addlegendentry{$E_m$ (constante)}

% Ec = 1.5 * cos^2(pi*t)  -> T0=2s, TE=1s
\addplot[popcurve, popblue, thick] {1.5*cos(deg(3.14159*x))^2};
\addlegendentry{$E_c$}

% Ep = 1.5 * sin^2(pi*t)
\addplot[popcurve, poporange, thick] {1.5*sin(deg(3.14159*x))^2};
\addlegendentry{$E_p$ (pesanteur ou élastique)}

% Annotations
\node[popblue, anchor=south] at (axis cs:0,1.5) {$E_{c,\text{max}}$};
\node[poporange, anchor=north] at (axis cs:0,0) {$E_{p,\text{min}}=0$};
\node[poporange, anchor=south] at (axis cs:0.5,1.5) {$E_{p,\text{max}}$};
\node[popblue, anchor=north] at (axis cs:0.5,0) {$E_{c,\text{min}}=0$};
\draw[dashed, popmuted] (axis cs:0.5,0) -- (axis cs:0.5,1.5);
\draw[dashed, popmuted] (axis cs:1,0) -- (axis cs:1,1.5);
\end{axis}
\end{tikzpicture}
\legende{Évolution temporelle des énergies pour un oscillateur non amorti ($T_0=\SI{2}{s}$). $E_c$ et $E_p$ oscillent à la fréquence $2f_0$ ; leur somme $E_m$ est constante.}
\end{popfigure}

\begin{propriete}[Oscillateur amorti (frottements visqueux)]
\begin{itemize}
    \item $E_m(t)$ : \textbf{Décroît exponentiellement} (ou linéairement par morceaux si frottements solides constants).
    \item $E_c(t)$ et $E_p(t)$ : Oscillations \textbf{amorties} (enveloppe décroissante).
    \item La surface entre la courbe $E_m(t)$ et l'axe du temps sur un intervalle $\Delta t$ représente l'énergie dissipée (travail des frottements).
\end{itemize}
\end{propriete}

\begin{methode}[Analyse d'un graphe $E(t)$ au Probatoire]
\begin{enumerate}
    \item \textbf{Identifier $E_m$ :} C'est la courbe la plus haute (somme des autres). Vérifier si elle est plate (conservation) ou descend (dissipation).
    \item \textbf{Lire les valeurs extrêmes :} $E_{c,\text{max}} = E_{m}$ (si $E_{p,\text{min}}=0$). $E_{p,\text{max}} = E_{m}$ (si $E_{c,\text{min}}=0$).
    \item \textbf{Calculer $v_{\text{max}}$ ou $x_{\text{max}}$ :} $v_{\text{max}} = \sqrt{2E_{c,\text{max}}/m}$, $x_{\text{max}} = \sqrt{2E_{p,\text{max}}/k}$.
    \item \textbf{Déterminer la période $T_0$ :} Mesurer la période $T_E$ d'oscillation de $E_c$ (ou $E_p$) : $T_0 = 2 T_E$.
    \item \textbf{Quantifier la dissipation :} Si $E_m$ baisse, $\Delta E_m = W_{\text{frott}}$ sur l'intervalle considéré.
\end{enumerate}
\end{methode}

\courssub{Erreurs classiques aux examens — Check-list finale}

Voici les 5 erreurs qui font perdre le plus de points au Probatoire. Relis-les avant chaque entraînement.

\begin{attention}[Top 5 des pièges à éviter]
\begin{enumerate}
    \item \textbf{Unités incohérentes :} Masse en grammes au lieu de kg ($m=\SI{50}{g} \neq \SI{50}{kg}$), $v$ en km/h au lieu de m/s, $k$ en N/cm au lieu de N/m. \textbf{Convertis TOUT en SI dès la lecture.}
    \item \textbf{Signe du travail :} $W = F d \cos\theta$. Frottement $\Rightarrow \theta=180^\circ \Rightarrow W < 0$. Oublier le signe moins dans $\Delta E_m = W_f$ conduit à $v^2 < 0$ (impossible) ou des distances négatives.
    \item \textbf{Origine des altitudes changeante :} Voir l'encadré \textbf{Attention} page 1. Choisis $h=0$ une fois pour toutes.
    \item \textbf{Confusion $W(\vec{P})$ et $\Delta E_{pp}$ :} $W(\vec{P}) = -\Delta E_{pp}$. Si tu utilises le théorème de l'énergie cinétique ($\Delta E_c = \sum W$), tu dois inclure $W(\vec{P})$. Si tu utilises l'énergie mécanique ($\Delta E_m = W_{\text{nc}}$), tu n'inclus \textbf{pas} $W(\vec{P})$ car il est déjà « rangé » dans $\Delta E_{pp}$. \textbf{Ne compte jamais les deux en même temps !}
    \item \textbf{Oublier l'énergie potentielle élastique au point final :} Dans les problèmes « plan incliné + ressort », à la compression max, le bloc a encore une $E_{pp}$ (souvent négative si $h=0$ au contact initial) ET une $E_{pe}$. Les deux doivent apparaître dans $E_{m,f}$.
\end{enumerate}
\end{attention}

\begin{savaistu}[Le Probatoire C et l'Énergie Mécanique]
Au Cameroun, les sujets de Probatoire série C comportent presque systématiquement un exercice complet (3 à 4 points sur 20) sur l'énergie mécanique.
\begin{itemize}
    \item \textbf{Thème favori :} « Mobile sur plan incliné (avec/sans frottement) $\to$ compression d'un ressort ». C'est la combinaison type qui teste \textbf{les deux énergies potentielles} et le \textbf{travail des frottements} en une seule question.
    \item \textbf{Autre classique :} « Lecture de courbes $E(t)$ pour un oscillateur » (souvent 1,5 à 2 points faciles si tu sais lire un graphique).
    \item \textbf{Conseil stratégique :} Maîtrise parfaitement la \textbf{Méthode type 3} (couplage plan/ressort). C'est le « grand exercice » qui distingue les notes moyennes des excellentes. Entraîne-toi à poser l'équation du second degré en $x_{\text{max}}$ sans faute de signe.
\end{itemize}
\end{savaistu}

\begin{exercice}[Entraînement Probatoire — Bille dans un tube courbe]
Une bille de masse $m = \SI{10}{g}$ est lâchée sans vitesse au point A ($h_A = \SI{50}{cm}$) dans un tube vertical courbe sans frottement. Au point B ($h_B = \SI{20}{cm}$), la bille entre en contact avec un ressort de raideur $k = \SI{100}{N/m}$ fixé horizontalement. La bille comprime le ressort jusqu'à l'arrêt momentané au point C.
\begin{enumerate}
    \item En choisissant l'origine des altitudes au niveau de B, exprimer $E_{m,A}$ et $E_{m,C}$.
    \item En appliquant la conservation de l'énergie mécanique, calculer la compression maximale $x_{\text{max}}$ du ressort.
    \item Si le tube est rugueux entre A et B (frottement constant $f = \SI{0,02}{N}$) et sans frottement entre B et C, retrouver $x_{\text{max}}$.
\end{enumerate}
\end{exercice}

\begin{corrige}[Exercice : Bille dans un tube courbe]
\begin{enumerate}
    \item \textbf{Origine $h=0$ au niveau B.} $m=\SI{0,010}{kg}$, $g=\SI{9,81}{N/kg}$.
    \begin{itemize}
        \item A : $v_A=0$, $h_A = \SI{0,30}{m}$ (au-dessus de B), $x_A=0$. $E_{m,A} = mgh_A = 0,010 \times 9,81 \times 0,30 = \mathbf{\SI{0,0294}{J}}$.
        \item C : $v_C=0$, $h_C = 0$ (même niveau que B, ressort horizontal), $x_C = x_{\text{max}}$. $E_{m,C} = \frac{1}{2}k x_{\text{max}}^2$.
    \end{itemize}
    \item \textbf{Sans frottement :} $E_{m,A} = E_{m,C} \implies mgh_A = \frac{1}{2}k x_{\text{max}}^2$.
    $x_{\text{max}} = \sqrt{\frac{2mgh_A}{k}} = \sqrt{\frac{2 \times 0,0294}{100}} = \sqrt{5,88 \times 10^{-4}} \approx \mathbf{\SI{0,0243}{m} = \SI{2,43}{cm}}$.
    \item \textbf{Avec frottement A$\to$B :} Travail du frottement $W_f = -f \cdot AB$. Distance $AB$ inconnue ? \textbf{Piège !} On ne connaît pas la longueur du tube $AB$, seulement la dénivelée.
    \textbf{Correction :} Le travail du frottement dépend du chemin. Si $AB$ n'est pas donné, on ne peut pas calculer $W_f$ numériquement.
    \emph{Hypothèse implicite probable :} Le tube est un quart de cercle de rayon $R = h_A - h_B = \SI{0,30}{m}$. Alors $AB = \frac{\pi}{2}R \approx \SI{0,47}{m}$.
    $W_f = -0,02 \times 0,47 = -\SI{0,0094}{J}$.
    $\Delta E_m = W_f \implies E_{m,C} - E_{m,A} = -0,0094$.
    $\frac{1}{2}k x_{\text{max}}^2 - 0,0294 = -0,0094 \implies \frac{1}{2}k x_{\text{max}}^2 = 0,0200$.
    $x_{\text{max}} = \sqrt{\frac{2 \times 0,0200}{100}} = \sqrt{4 \times 10^{-4}} = \mathbf{\SI{0,020}{m} = \SI{2,0}{cm}}$.
    \textbf{Moralité :} Vérifie toujours si la distance de frottement est donnée ou déductible (longueur de piste, rayon de courbure).
\end{enumerate}
\end{corrige}

\courssub{Bilan des compétences validées}

À l'issue de cette leçon complète (Sections 1 à 6), tu dois être capable de :
\begin{itemize}
    \item \cle{Savoir} : Définir $E_c$, $E_{pp}$, $E_{pe}$, $E_m$ ; énoncer le principe de conservation de $E_m$ et la relation $\Delta E_m = W_{\text{nc}}$.
    \item \cle{Savoir-faire} : Choisir judicieusement l'origine des altitudes ; appliquer la conservation de $E_m$ pour trouver $v$, $h$ ou $x$ ; calculer le travail des frottements et utiliser $\Delta E_m = W_f$ ; modéliser un système « plan incliné + ressort » (équation du 2nd degré) ; lire et exploiter des graphes $E(t)$.
    \item \cle{Savoir-être} : Rigueur dans les unités (SI), cohérence des signes, clarté du schéma initial (forces, altitudes, référentiels), analyse dimensionnelle systématique.
\end{itemize}

Tu es maintenant prêt(e) pour les épreuves du Probatoire. Travaille les exercices types, chronomètre-toi, et vise l'excellence ! \textbf{Courage pour la suite du programme.}

\newpage

\coursec{Activité d'intégration}

\coursec{Leçon N05 : Énergie potentielle et énergie mécanique}

\courssub{Objectifs de l'activité}
Cette activité d'intégration vise à mobiliser l'ensemble des savoirs et savoir-faire de la leçon sur l'énergie potentielle et l'énergie mécanique à travers deux situations concrètes ancrées dans le quotidien camerounais. Elle permet de :
\begin{itemize}
    \item Appliquer le \cle{principe de conservation de l'énergie mécanique} dans un cas idéal (sans frottements) et le \cle{théorème de l'énergie mécanique} dans un cas réel (avec frottements).
    \item Manipuler les expressions de l'\cle{énergie potentielle de pesanteur}, de l'\cle{énergie cinétique} et du \cle{travail d'une force}.
    \item Analyser l'influence de la vitesse sur la distance de freinage d'un véhicule et en tirer des conséquences concrètes pour la sécurité routière.
    \item Développer une argumentation scientifique rigoureuse à destination d'un public non spécialisé (message de sensibilisation).
\end{itemize}

\courssub{Ressources à mobiliser}
\begin{propriete}[Formules utiles]
\begin{itemize}
    \item \cle{Énergie potentielle de pesanteur} (voisinage de la Terre) : $E_p = m g h$ avec $g = \SI{9,81}{N/kg}$ (ou $\SI{9,81}{m/s^2}$).
    \item \cle{Énergie cinétique} : $E_c = \frac{1}{2} m v^2$.
    \item \cle{Énergie mécanique} : $E_m = E_c + E_p$.
    \item \cle{Travail d'une force constante} : $W = \vec{F} \cdot \vec{AB} = F\, AB \cos(\vec{F}, \vec{AB})$.
    \item \cle{Conservation de l'énergie mécanique} (système isolé, forces intérieures conservatrices) : $E_{m,i} = E_{m,f}$.
    \item \cle{Théorème de l'énergie mécanique} (système non isolé ou forces non conservatrices) : $\Delta E_m = W_{\text{ext, non cons}}$.
\end{itemize}
\end{propriete}

\begin{attention}[Pièges à éviter]
\begin{itemize}
    \item Ne pas confondre la hauteur de chute $h$ (verticale) et la longueur du toboggan $L$ (inclinée) pour le calcul du travail des frottements : $W_f = -F_f \times L$.
    \item Vérifier l'homogénéité des unités : convertir systématiquement km/h en m/s ($1\,\text{km/h} = \frac{1}{3,6}\,\text{m/s}$).
    \item Le travail d'une force de frottement (ou freinage) est \textbf{négatif} car la force s'oppose au déplacement ($\cos 180^\circ = -1$).
    \item L'énergie cinétique est proportionnelle au \textbf{carré} de la vitesse : $E_c \propto v^2$.
\end{itemize}
\end{attention}

\begin{exoresolu}[Activité d'intégration : Toboggan (Yaoundé) et Freinage taxi-brousse (N3)]
\textbf{Situation}

\noindent \textbf{Première partie : Le toboggan du parc de la Réunification (Yaoundé).}
Un enfant de masse $m = \SI{30,0}{kg}$ se trouve au point le plus haut d'un toboggan droit, à une hauteur $h = \SI{4,00}{m}$ au-dessus du sol. La longueur du toboggan est $L = \SI{5,00}{m}$. L'enfant se laisse glisser sans élan initial ($v_A = 0$).
\begin{enumerate}
    \item \textbf{Cas idéal :} On néglige dans un premier temps les frottements entre l'enfant et le toboggan. On modélise l'enfant comme un point matériel et le système « enfant + Terre » comme isolé sur le plan mécanique.
    \item \textbf{Cas réel :} On prend en compte les frottements. La force de frottement moyenne exercée par le toboggan sur l'enfant est estimée à $\|\vec{F}_f\| = \SI{40,0}{N}$, orientée opposée au mouvement.
\end{enumerate}

\noindent \textbf{Deuxième partie : Freinage d'un taxi-brousse sur la route nationale N3 (Douala--Yaoundé).}
Un taxi-brousse de masse $M = \SI{1500}{kg}$ (chargé) roule à la vitesse réglementaire de $v_1 = \SI{72,0}{km/h}$ ($\SI{20,0}{m/s}$) sur une portion rectiligne et plane. Le conducteur aperçoit un obstacle et freine brutalement. Le véhicule s'immobilise après avoir parcouru une distance de freinage $d_1 = \SI{50,0}{m}$. On suppose que la force de freinage moyenne $\|\vec{F}_f\|$ est constante pendant le freinage et que la route est horizontale.

\textbf{Consignes}

\noindent \textbf{Première partie : Le toboggan}
\begin{enumerate}
    \item Choix du référentiel et du système. Justifier que le système « enfant + Terre » peut être considéré comme isolé dans le cas idéal.
    \item Écrire l'expression de l'énergie mécanique $E_m$ de l'enfant en A (haut) et en B (bas) pour le cas idéal. En déduire la vitesse $v_B$ en bas du toboggan.
    \item Pour le cas réel, identifier les forces non conservatrices dont le travail n'est pas nul. Écrire la relation fondamentale $\Delta E_m = W_{\text{ext}}$ (théorème de l'énergie mécanique) entre A et B.
    \item Calculer la vitesse $v'_B$ de l'enfant en bas du toboggan avec frottements. Comparer $v'_B$ et $v_B$ et interpréter physiquement la différence.
\end{enumerate}

\noindent \textbf{Deuxième partie : Le taxi-brousse}
\begin{enumerate}
    \item[5.] Exprimer l'énergie cinétique initiale $E_{c,i}$ du taxi-brousse. Calculer sa valeur numérique.
    \item[6.] Appliquer le théorème de l'énergie mécanique au système « taxi-brousse » (forces extérieures : freinage, frottements air/route regroupés en une force de freinage moyenne $\vec{F}_f$) entre le début du freinage (point I) et l'arrêt (point F). En déduire l'expression de la magnitude de la force de freinage moyenne $\|\vec{F}_f\|$ en fonction de $M$, $v_1$ et $d_1$. Calculer sa valeur.
    \item[7.] Le conducteur double sa vitesse initiale : $v_2 = 2v_1 = \SI{144}{km/h}$ ($\SI{40,0}{m/s}$). En supposant la même force de freinage moyenne $\|\vec{F}_f\|$, déterminer la nouvelle distance de freinage $d_2$. Comparer $d_2$ à $d_1$.
    \item[8.] \textbf{Message de sensibilisation :} Rédiger un court message (5 à 7 lignes) destiné aux conducteurs de taxi-brousse et aux usagers de la route, expliquant \emph{physiquement} pourquoi « doubler la vitesse quadruple la distance de freinage » et les risques encourus. Utiliser les termes « énergie cinétique », « travail des freins », « distance de freinage ».
\end{enumerate}

\tcblower

\textbf{Corrigé détaillé}

\noindent \textbf{Première partie : Le toboggan}

\noindent \textbf{1. Système et référentiel.}
On choisit le référentiel terrestre (supposé galiléen à l'échelle de l'expérience). Le système étudié est l'enfant (modèle du point matériel) en interaction avec la Terre. Les forces extérieures sont : le poids $\vec{P} = m\vec{g}$ (force conservative, dérivée de $E_p$) et la réaction du toboggan $\vec{R}$ (perpendiculaire au déplacement, travail nul). Dans le cas idéal (sans frottements), il n'y a pas de forces extérieures non conservatrices dont le travail soit non nul $\rightarrow$ le système « enfant + Terre » est \textbf{mécaniquement isolé}. L'énergie mécanique se conserve.

\noindent \textbf{2. Cas idéal : Conservation de $E_m$.}
Point A (haut) : $v_A = 0$, $h_A = h = \SI{4,00}{m}$.
\[ E_{m,A} = E_{c,A} + E_{p,A} = 0 + m g h \]
Point B (bas) : $h_B = 0$, $v_B = ?$
\[ E_{m,B} = E_{c,B} + E_{p,B} = \frac{1}{2} m v_B^2 + 0 \]
Conservation : $E_{m,A} = E_{m,B}$
\[ m g h = \frac{1}{2} m v_B^2 \quad \Rightarrow \quad v_B = \sqrt{2 g h} \]
Calcul numérique : $v_B = \sqrt{2 \times 9,81 \times 4,00} = \sqrt{78,48} \approx \textbf{\SI{8,86}{m/s}}$.
\emph{Analyse dimensionnelle :} $[g] = \text{L}\text{T}^{-2}$, $[h] = \text{L}$ $\rightarrow$ $[2gh] = \text{L}^2\text{T}^{-2}$ $\rightarrow$ $[\sqrt{2gh}] = \text{L}\text{T}^{-1}$ (vitesse). Correct.

\noindent \textbf{3. Cas réel : Forces non conservatrices.}
La force de frottement $\vec{F}_f$ (magnitude $F_f = \SI{40,0}{N}$) est une force extérieure non conservative (ou force intérieure dissipative si on élargit le système au toboggan, mais ici on garde « enfant » comme système, $\vec{F}_f$ est extérieure). Son travail est non nul et négatif.
Travail de $\vec{F}_f$ sur la longueur $L = \SI{5,00}{m}$ :
\[ W_f = \vec{F}_f \cdot \vec{AB} = - F_f L = -40,0 \times 5,00 = \SI{-200}{J} \]
Théorème de l'énergie mécanique entre A et B :
\[ \Delta E_m = E_{m,B} - E_{m,A} = W_f \]
\[ \left( \frac{1}{2} m v'_B{}^2 \right) - (m g h) = - F_f L \]

\noindent \textbf{4. Calcul de $v'_B$ et comparaison.}
\[ \frac{1}{2} m v'_B{}^2 = m g h - F_f L \]
\[ v'_B{}^2 = 2 g h - \frac{2 F_f L}{m} \]
\[ v'_B = \sqrt{ 2 \times 9,81 \times 4,00 - \frac{2 \times 40,0 \times 5,00}{30,0} } \]
\[ v'_B = \sqrt{ 78,48 - \frac{400}{30,0} } = \sqrt{ 78,48 - 13,33 } = \sqrt{ 65,15 } \approx \textbf{\SI{8,07}{m/s}} \]
\textbf{Comparaison :} $v'_B = \SI{8,07}{m/s} < v_B = \SI{8,86}{m/s}$.
\textbf{Interprétation :} Le travail négatif des frottements ($\SI{-200}{J}$) a prélevé de l'énergie mécanique au système. Cette énergie a été transformée en énergie thermique (échauffement du pantalon et du toboggan) et en énergie sonore (bruit de glissement), conformément au premier principe de la thermodynamique. L'enfant arrive moins vite en bas.

\vspace{0.5cm}
\noindent \textbf{Deuxième partie : Le taxi-brousse}

\noindent \textbf{5. Énergie cinétique initiale.}
Conversion de la vitesse : $v_1 = \SI{72,0}{km/h} = \frac{72,0}{3,6} = \SI{20,0}{m/s}$.
\[ E_{c,i} = \frac{1}{2} M v_1^2 = \frac{1}{2} \times 1500 \times (20,0)^2 = 750 \times 400 = \textbf{\SI{300\,000}{J}} = \textbf{\SI{300}{kJ}} \]

\noindent \textbf{6. Force de freinage moyenne.}
Système : Taxi-brousse (point matériel). Forces extérieures travaillantes : Force de freinage/frottements $\vec{F}_f$ (constante, opposée au mouvement). Poids et réaction normale : travail nul (perpendiculaires au déplacement horizontal).
État initial I : $v_i = v_1 = \SI{20,0}{m/s}$, $E_{c,i} = \SI{300}{kJ}$.
État final F : $v_f = 0$, $E_{c,f} = 0$.
Variation d'énergie mécanique (ici seulement cinétique car pas de variation d'altitude) :
\[ \Delta E_m = E_{c,f} - E_{c,i} = -\SI{300\,000}{J} \]
Travail de la force de freinage sur la distance $d_1 = \SI{50,0}{m}$ :
\[ W_f = \vec{F}_f \cdot \vec{d_1} = - \|\vec{F}_f\| d_1 \]
Théorème : $\Delta E_m = W_f$
\[ -300\,000 = - \|\vec{F}_f\| \times 50,0 \]
\[ \|\vec{F}_f\| = \frac{300\,000}{50,0} = \textbf{\SI{6\,000}{N}} \]

\noindent \textbf{7. Distance de freinage à vitesse doublée.}
Nouvelle vitesse $v_2 = 2 v_1 = \SI{40,0}{m/s}$ ($\SI{144}{km/h}$).
Nouvelle énergie cinétique initiale :
\[ E'_{c,i} = \frac{1}{2} M v_2^2 = \frac{1}{2} M (2v_1)^2 = 4 \times \frac{1}{2} M v_1^2 = 4 \times 300\,000 = \textbf{\SI{1\,200\,000}{J}} = \textbf{\SI{1,2}{MJ}} \]
La force de freinage maximale (adhérence pneus/route, puissance freins) est supposée inchangée : $\|\vec{F}_f\| = \SI{6\,000}{N}$.
Théorème : $\Delta E'_m = W'_f$
\[ 0 - 1\,200\,000 = - 6\,000 \times d_2 \]
\[ d_2 = \frac{1\,200\,000}{6\,000} = \textbf{\SI{200}{m}} \]
\textbf{Comparaison :} $d_2 = \SI{200}{m} = 4 \times d_1$ ($d_1 = \SI{50}{m}$).
\emph{Relation générale :} $d = \frac{\frac{1}{2} M v^2}{\|\vec{F}_f\|} \propto v^2$. Doubler $v$ quadruple $d$.

\noindent \textbf{8. Message de sensibilisation routière.}
\begin{quote}
\textbf{« Attention conducteurs ! Doubler votre vitesse ne double pas le danger, elle le quadruple.
Physiquement, l'énergie cinétique de votre taxi-brousse vaut $\frac{1}{2}mv^2$ : à \SI{144}{km/h}, elle est \textbf{4 fois plus grande} qu'à \SI{72}{km/h}. Or, pour vous arrêter, les freins doivent fournir un travail égal à cette énergie ($W = F \times d$). Comme la force de freinage maximale $F$ ne change pas, la distance de freinage $d$ devient \textbf{4 fois plus longue} (\SI{200}{m} au lieu de \SI{50}{m}). Sur la route, ces \SI{150}{m} supplémentaires coûtent des vies. Respectez les limites de vitesse : votre énergie cinétique ne pardonne pas l'excès. »}
\end{quote}
\end{exoresolu}

\courssub{Bilan de l'activité}
Cette activité a permis de mettre en évidence les points fondamentaux suivants :
\begin{itemize}
    \item \textbf{Conversion d'énergies :} Dans le toboggan idéal, l'énergie potentielle de pesanteur se convertit intégralement en énergie cinétique ($mgh = \frac{1}{2}mv^2$). La masse se simplifie : tous les objets tombent avec la même accélération (équivalence galiléenne).
    \item \textbf{Rôle des forces dissipatives :} Les frottements (toboggan) ou le freinage (taxi-brousse) effectuent un travail négatif ($W < 0$) qui diminue l'énergie mécanique du système : $\Delta E_m = W_{\text{frott}} < 0$. L'énergie mécanique « perdue » se transforme en énergie thermique (chaleur) et sonore, conformément au principe de conservation de l'énergie \emph{totale} (1er principe de la thermodynamique).
    \item \textbf{Théorème de l'énergie mécanique :} C'est l'outil universel : $\Delta E_m = \sum W_{\text{ext, non cons}}$. Il gère aussi bien la conservation ($W=0$) que la dissipation ($W<0$) ou l'apport d'énergie ($W>0$, ex: moteur).
    \item \textbf{Sécurité routière et physique :} La relation $d \propto v^2$ (démontrée via $E_c = \frac{1}{2}mv^2$ et $W = Fd$) est une loi physique rigoureuse, non une simple statistique. Elle explique pourquoi les excès de vitesse sont la première cause de gravité des accidents sur nos routes (N3, N10, autoroute Yaoundé-Nsimalen). La distance de réaction (non calculée ici, proportionnelle à $v$) s'ajoute à la distance de freinage, aggravant encore le risque.
    \item \textbf{Méthodologie de résolution :} 1) Définir système/référentiel. 2) Lister les forces et classer leur travail (nul, conservatif, non conservatif). 3) Choisir la relation adaptée (Conservation ou $\Delta E_m = W$). 4) Projeter/Calculer avec unités SI. 5) Interpréter le signe et l'ordre de grandeur.
\end{itemize}
\begin{aretenir}[À retenir pour le Probatoire]
\begin{itemize}
    \item Savoir écrire $E_p = mgh$, $E_c = \frac{1}{2}mv^2$, $E_m = E_c + E_p$.
    \item Savoir énoncer : « Si le système est isolé et les forces intérieures conservatrices, $E_m$ est constante. » (Démonstration exigible : $\Delta E_m = W_{\text{int}} + W_{\text{ext}} = 0$).
    \item Savoir appliquer $\Delta E_m = W_{\text{frott}}$ ou $W_{\text{freinage}}$.
    \item Maîtriser la conversion km/h $\leftrightarrow$ m/s (diviser/multiplier par 3,6).
    \item Savoir expliquer le lien $d \propto v^2$ pour le freinage.
\end{itemize}
\end{aretenir}

\newpage

\coursec{Méthodes et exercices résolus}

\courssub{Méthodes et exercices résolus}

\begin{methode}[Exprimer l'énergie potentielle de pesanteur]
    \textbf{Objectif :} Calculer $E_p$ d'un système dans le champ de pesanteur terrestre uniforme.
    \begin{enumerate}
        \item \textbf{Choisir un référentiel galiléen} lié à la Terre (souvent le sol ou le niveau de la mer).
        \item \textbf{Fixer l'origine des altitudes ($z=0$)} : c'est le niveau de référence où $E_p = 0$. Ce choix est \emph{arbitraire} mais doit rester \textbf{constant} pendant tout le problème.
        \item \textbf{Repérer l'altitude $z$} du centre d'inertie $G$ du solide par rapport à ce niveau de référence ($z > 0$ au-dessus, $z < 0$ en dessous).
        \item \textbf{Appliquer la formule} : $E_p = m g z$.
        \item \textbf{Vérifier l'homogénéité} : $[E_p] = \text{kg} \cdot \text{m}\cdot\text{s}^{-2} \cdot \text{m} = \text{J}$.
    \end{enumerate}
    \begin{aretenir}[Formule]
        Pour un solide de masse $m$ au voisinage de la Terre, l'énergie potentielle de pesanteur vaut :
        \[ E_p = m g z \]
        où $g = \SI{9,81}{N/kg}$ (ou $\SI{9,81}{m/s^2}$) et $z$ est l'altitude du centre d'inertie par rapport au niveau de référence choisi.
    \end{aretenir}
\end{methode}

\begin{exoresolu}[Chute d'eau au barrage de Nachtigal]
    \textbf{Énoncé :} Au barrage de Nachtigal, une masse d'eau $m = 2,5 \times 10^5\ \SI{}{kg}$ tombe d'une hauteur de $50,0\ \SI{}{m}$ pour actionner les turbines. On choisit le niveau de la turbine comme référence des altitudes ($z=0$).
    \begin{enumerate}
        \item Quelle est l'énergie potentielle de pesanteur de cette masse d'eau au départ de la chute ?
        \item Quelle est sa variation d'énergie potentielle $\Delta E_p$ entre le départ et l'arrivée ?
    \end{enumerate}
    \tcblower
    \textbf{Solution détaillée :}
    \begin{enumerate}
        \item \textbf{Données :} $m = 2,5 \times 10^5\ \SI{}{kg}$ ; $g = \SI{9,81}{N/kg}$ ; $z_i = 50,0\ \SI{}{m}$ (niveau de la retenue) ; $z_f = 0\ \SI{}{m}$ (niveau turbine, référence).
        \item \textbf{Calcul de $E_{p,i}$ :}
        \[ E_{p,i} = m g z_i = (2,5 \times 10^5) \times 9,81 \times 50,0 = 1,22625 \times 10^8\ \SI{}{J} \]
        On arrondit à 2 chiffres significatifs (donnée $m$ à 2 chiffres) : $\boxed{E_{p,i} \approx 1,2 \times 10^8\ \SI{}{J}}$.
        \item \textbf{Calcul de $\Delta E_p$ :}
        \[ \Delta E_p = E_{p,f} - E_{p,i} = m g z_f - m g z_i = 0 - 1,22625 \times 10^8 = -1,22625 \times 10^8\ \SI{}{J} \]
        $\boxed{\Delta E_p \approx -1,2 \times 10^8\ \SI{}{J}}$.
        \item \textbf{Interprétation :} L'énergie potentielle diminue (signe négatif), elle est convertie en énergie cinétique de l'eau puis en énergie électrique via les alternateurs d'ENEO.
    \end{enumerate}
    \textbf{Conclusion :} L'eau perd $1,2 \times 10^8\ \SI{}{J}$ d'énergie potentielle de pesanteur lors de sa chute.
\end{exoresolu}

\begin{methode}[Exprimer l'énergie potentielle élastique]
    \textbf{Objectif :} Calculer $E_{p,\text{el}}$ stockée dans un ressort à réponse linéaire (loi de Hooke).
    \begin{enumerate}
        \item \textbf{Identifier la raideur $k$} du ressort (en $\SI{}{N/m}$) et sa longueur à vide $L_0$.
        \item \textbf{Déterminer la déformation $x$} (allongement ou raccourcissement) par rapport à la position d'équilibre (sans force) : $x = L - L_0$.
        \item \textbf{Vérifier la linéarité} : la loi de Hooke $F = k x$ doit être valide (domaine élastique).
        \item \textbf{Appliquer la formule} : $E_{p,\text{el}} = \frac{1}{2} k x^2$.
        \item \textbf{Noter que $E_{p,\text{el}} \ge 0$} toujours (énergie stockée, scalaire positif).
    \end{enumerate}
    \begin{aretenir}[Formule]
        Pour un ressort de raideur $k$ déformé de $x$ par rapport à sa longueur à vide :
        \[ E_{p,\text{el}} = \frac{1}{2} k x^2 \]
        Unité : Joule ($\SI{}{J}$). $k$ en $\SI{}{N/m}$, $x$ en $\SI{}{m}$.
    \end{aretenir}
\end{methode}

\begin{exoresolu}[Amortisseur de moto-taxi]
    \textbf{Énoncé :} L'amortisseur arrière d'un moto-taxi Yamaha DT 125 (très courant à Douala) comporte un ressort de raideur $k = 3,2 \times 10^4\ \SI{}{N/m}$. Au repos, sous la charge du pilote ($m = 80\ \SI{}{kg}$), le ressort se raccourcit de $2,5\ \SI{}{cm}$ par rapport à sa longueur à vide. On néglige la masse du ressort.
    \begin{enumerate}
        \item Calculer l'énergie potentielle élastique stockée dans le ressort à l'équilibre statique.
        \item Si le moto-taxi franchit un dos-d'âne et que le ressort se compresse de $5,0\ \SI{}{cm}$ au total (par rapport à $L_0$), quelle est l'énergie élastique stockée à cet instant ?
    \end{enumerate}
    \tcblower
    \textbf{Solution détaillée :}
    \begin{enumerate}
        \item \textbf{Données :} $k = 3,2 \times 10^4\ \SI{}{N/m}$ ; $x_1 = 2,5\ \SI{}{cm} = 0,025\ \SI{}{m}$.
        \[ E_{p,\text{el},1} = \frac{1}{2} k x_1^2 = 0,5 \times 3,2 \times 10^4 \times (0,025)^2 \]
        \[ E_{p,\text{el},1} = 1,6 \times 10^4 \times 6,25 \times 10^{-4} = 10,0\ \SI{}{J} \]
        $\boxed{E_{p,\text{el},1} = 10,0\ \SI{}{J}}$.
        \item \textbf{Données :} $x_2 = 5,0\ \SI{}{cm} = 0,050\ \SI{}{m}$.
        \[ E_{p,\text{el},2} = \frac{1}{2} k x_2^2 = 0,5 \times 3,2 \times 10^4 \times (0,050)^2 \]
        \[ E_{p,\text{el},2} = 1,6 \times 10^4 \times 25 \times 10^{-4} = 40,0\ \SI{}{J} \]
        $\boxed{E_{p,\text{el},2} = 40,0\ \SI{}{J}}$.
        \item \textbf{Remarque :} L'énergie est proportionnelle au \emph{carré} de la déformation. Doubler la déformation quadruple l'énergie stockée ($40 = 4 \times 10$).
    \end{enumerate}
    \textbf{Conclusion :} L'amortisseur stocke $10\ \SI{}{J}$ à l'arrêt et $40\ \SI{}{J}$ en compression maximale sur le dos-d'âne.
\end{exoresolu}

\begin{methode}[Exprimer l'énergie mécanique d'un système]
    \textbf{Objectif :} Calculer $E_m = E_c + E_p$ à un instant donné.
    \begin{enumerate}
        \item \textbf{Définir le système} étudié (ex: \{bloc + Terre\}, \{bloc + ressort\}).
        \item \textbf{Identifier les énergies potentielles présentes} : pesanteur ($m g z$), élastique ($\frac{1}{2}k x^2$), ou les deux.
        \item \textbf{Calculer l'énergie cinétique} $E_c = \frac{1}{2} m v^2$ (vitesse du centre d'inertie dans le référentiel galiléen).
        \item \textbf{Faire la somme} : $E_m = E_c + \sum E_p$.
        \item \textbf{Soigner les unités} : tout en SI ($m$ en kg, $v$ en m/s, $z$ en m, $k$ en N/m, $x$ en m) $\rightarrow$ $E_m$ en Joules.
    \end{enumerate}
    \begin{aretenir}[Définition]
        L'énergie mécanique d'un système est la somme de son énergie cinétique et de ses énergies potentielles :
        \[ E_m = E_c + E_p = \frac{1}{2} m v^2 + m g z + \frac{1}{2} k x^2 + \dots \]
    \end{aretenir}
\end{methode}

\begin{exoresolu}[Lancer de balle au stade Omnisport]
    \textbf{Énoncé :} Un footballeur lance un ballon de masse $m = 0,43\ \SI{}{kg}$ verticalement vers le haut avec une vitesse initiale $v_0 = 18,0\ \SI{}{m/s}$. On choisit le niveau du sol comme référence ($z=0$). Le ballon est lancé depuis $z_0 = 1,80\ \SI{}{m}$.
    \begin{enumerate}
        \item Calculer l'énergie mécanique $E_{m,0}$ du système \{ballon + Terre\} au lancement.
        \item Déterminer l'altitude maximale $z_{\max}$ atteinte par le ballon (on suppose l'énergie mécanique conservée).
        \item Calculer l'énergie mécanique à $z = 10,0\ \SI{}{m}$.
    \end{enumerate}
    \tcblower
    \textbf{Solution détaillée :}
    \begin{enumerate}
        \item \textbf{Au lancement ($t=0$) :}
        $E_{c,0} = \frac{1}{2} m v_0^2 = 0,5 \times 0,43 \times (18,0)^2 = 0,215 \times 324 = 69,66\ \SI{}{J}$.
        $E_{p,0} = m g z_0 = 0,43 \times 9,81 \times 1,80 = 7,59\ \SI{}{J}$.
        $E_{m,0} = E_{c,0} + E_{p,0} = 69,66 + 7,59 = 77,25\ \SI{}{J}$.
        $\boxed{E_{m,0} \approx 77,3\ \SI{}{J}}$.
        \item \textbf{Au point haut ($v=0$) :} $E_c = 0$, donc $E_m = E_p = m g z_{\max}$.
        Conservation : $E_{m,0} = m g z_{\max}$.
        \[ z_{\max} = \frac{E_{m,0}}{m g} = \frac{77,25}{0,43 \times 9,81} = \frac{77,25}{4,2183} \approx 18,31\ \SI{}{m} \]
        $\boxed{z_{\max} \approx 18,3\ \SI{}{m}}$.
        \item \textbf{À $z = 10,0\ \SI{}{m}$ :} Si conservation, $E_m$ est constante.
        $\boxed{E_m(10\ \text{m}) = 77,3\ \SI{}{J}}$.
        \emph{Vérification :} $E_p = 0,43 \times 9,81 \times 10 = 42,18\ \SI{}{J}$. $E_c = 77,25 - 42,18 = 35,07\ \SI{}{J}$. $v = \sqrt{2E_c/m} \approx 12,8\ \SI{}{m/s}$. Cohérent.
    \end{enumerate}
    \textbf{Conclusion :} L'énergie mécanique vaut $77,3\ \SI{}{J}$ tout au long de la montée (hypothèse sans frottements).
\end{exoresolu}

\begin{methode}[Exploiter la conservation de l'énergie mécanique ($\Delta E_m = 0$)]
    \textbf{Objectif :} Résoudre un problème où seules des forces conservatrices travaillent (pesanteur, force élastique).
    \begin{enumerate}
        \item \textbf{Vérifier les conditions} : Le système ne subit que des forces conservatrices (poids, force de rappel du ressort). \textbf{Aucun frottement, aucun choc inélastique, aucune force extérieure non conservative (moteur, frein, main qui pousse)}.
        \item \textbf{Choisir deux états} : Initial (i) et Final (f) pertinents pour la question.
        \item \textbf{Écrire l'égalité} : $E_{m,i} = E_{m,f}$ $\Leftrightarrow$ $E_{c,i} + E_{p,i} = E_{c,f} + E_{p,f}$.
        \item \textbf{Exprimer chaque terme} avec les formules ($ \frac{1}{2}mv^2, mgz, \frac{1}{2}kx^2 $).
        \item \textbf{Résoudre l'équation} pour l'inconnue (vitesse, hauteur, compression, raideur...).
        \item \textbf{Conclure avec unité et sens physique}.
    \end{enumerate}
    \begin{attention}[Piège]
        Ne pas écrire "$E_c = E_p$" mais "$E_{c,i} + E_{p,i} = E_{c,f} + E_{p,f}$". L'énergie mécanique est \emph{constante}, pas les énergies individuelles.
    \end{attention}
\end{methode}

\begin{exoresolu}[Toboggan du parc aquatique de Yaoundé]
    \textbf{Énoncé :} Un enfant de masse $m = 30\ \SI{}{kg}$ glisse sans frottement sur un toboggan de hauteur $h = 4,50\ \SI{}{m}$. Il part du repos au sommet ($v_A = 0$). On note B l'arrivée au niveau du sol ($z=0$).
    \begin{enumerate}
        \item En appliquant le principe de conservation de l'énergie mécanique entre A et B, déterminer la vitesse $v_B$ de l'enfant à l'arrivée.
        \item L'enfant freine ensuite sur une piste horizontale rugueuse et s'arrête en $d = 6,0\ \SI{}{m}$. Quelle est la force de frottement moyenne $\bar{f}$ ? (Utiliser le théorème de l'énergie cinétique).
    \end{enumerate}
    \tcblower
    \textbf{Solution détaillée :}
    \begin{enumerate}
        \item \textbf{Conservation entre A (sommet) et B (bas) :}
        Système : \{Enfant + Terre\}. Forces : Poids (conservative), Réaction du toboggan (perpendiculaire au déplacement $\rightarrow$ travail nul). $\Rightarrow E_m$ conservée.
        État A : $v_A = 0 \Rightarrow E_{c,A} = 0$ ; $z_A = h = 4,50\ \SI{}{m} \Rightarrow E_{p,A} = mgh$.
        État B : $z_B = 0 \Rightarrow E_{p,B} = 0$ ; $v_B = ? \Rightarrow E_{c,B} = \frac{1}{2}m v_B^2$.
        \[ E_{m,A} = E_{m,B} \Rightarrow m g h = \frac{1}{2} m v_B^2 \]
        La masse $m$ se simplifie (galiléen) :
        \[ v_B = \sqrt{2 g h} = \sqrt{2 \times 9,81 \times 4,50} = \sqrt{88,29} \approx 9,40\ \SI{}{m/s} \]
        $\boxed{v_B \approx 9,40\ \SI{}{m/s}}$ (soit $\approx 33,8\ \SI{}{km/h}$).
        \item \textbf{Freinage sur partie horizontale (B $\to$ C) :}
        Système : \{Enfant\}. Forces : Poids + Normale (travail nul) + Frottement $\vec{f}$ (non conservative).
        Théorème de l'énergie cinétique : $W_{f} = \Delta E_c = E_{c,C} - E_{c,B}$.
        $E_{c,C} = 0$ (arrêt). $E_{c,B} = \frac{1}{2} m v_B^2 = m g h$ (d'après question 1).
        Travail du frottement (force constante opposée au déplacement) : $W_f = - \bar{f} \times d$.
        \[ -\bar{f} d = - m g h \Rightarrow \bar{f} = \frac{m g h}{d} = \frac{30 \times 9,81 \times 4,50}{6,0} \]
        \[ \bar{f} = \frac{1324,35}{6,0} = 220,725\ \SI{}{N} \]
        $\boxed{\bar{f} \approx 221\ \SI{}{N}}$.
    \end{enumerate}
    \textbf{Conclusion :} La vitesse au bas du toboggan ne dépend pas de la masse ni de la forme de la piste (si pas de frottement), seulement de la hauteur. Le freinage mobilise une force moyenne de $221\ \SI{}{N}$.
\end{exoresolu}

\begin{methode}[Identifier conservation / non-conservation de $E_m$]
    \textbf{Objectif :} Décider si $E_m$ est constante ou non dans une situation donnée.
    \begin{enumerate}
        \item \textbf{Lister toutes les forces} extérieures agissant sur le système (poids, frottements, réactions de contact, forces de liaison, forces motrices, forces électriques/magnétiques...).
        \item \textbf{Classer chaque force} :
            \begin{itemize}
                \item \textbf{Conservative} : Poids (champ de pesanteur uniforme), Force de rappel d'un ressort parfait, Force électrostatique, Force magnétique (Lorentz, ne travaille pas).
                \item \textbf{Non conservative (dissipative ou motrice)} : Frottements solides/fluides (travail négatif), Force de freinage, Force motrice (moteur, main qui pousse), Force de tension d'un fil non élastique (si elle modifie $E_m$), Choc inélastique (perte d'énergie interne).
            \end{itemize}
        \item \textbf{Regarder les travaux} : Si la somme des travaux des forces non conservatrices est nulle ($W_{\text{nc}} = 0$), alors $\Delta E_m = 0$ (conservation). Sinon $\Delta E_m = W_{\text{nc}} \neq 0$.
        \item \textbf{Attention aux forces de liaison} : Une réaction de contact idéale (sans frottement) est perpendiculaire à la vitesse $\rightarrow$ travail nul $\rightarrow$ ne change pas $E_m$. Mais si le point d'application se déplace (ex: piston mobile), elle peut travailler.
    \end{enumerate}
    \begin{aretenir}[Critère]
        \textbf{Conservation de $E_m$} $\Leftrightarrow$ \textbf{Somme des travaux des forces non conservatrices nulle} ($W_{\text{nc}} = 0$).
        \textbf{Non conservation} $\Leftrightarrow$ $W_{\text{nc}} \neq 0$ $\Rightarrow$ $\Delta E_m = W_{\text{nc}}$.
    \end{aretenir}
\end{methode}

\begin{exoresolu}[Freinage d'urgence sur la route Yaoundé-Douala]
    \textbf{Énoncé :} Un camion de masse $M = 12\ \SI{}{t} = 12\,000\ \SI{}{kg}$ roule à $v_0 = 90\ \SI{}{km/h} = 25,0\ \SI{}{m/s}$ sur une portion plane. Le conducteur freine brutalement. Les forces de frottement (freins + air + route) effectuent un travail total $W_{\text{frot}} = -3,75 \times 10^6\ \SI{}{J}$ jusqu'à l'arrêt complet.
    \begin{enumerate}
        \item L'énergie mécanique du système \{Camion + Terre\} est-elle conservée pendant le freinage ? Justifier.
        \item Vérifier que le travail des frottements est bien égal à la variation d'énergie mécanique.
        \item Calculer la distance de freinage $d$ si la force de frottement moyenne est constante et vaut $F = 150\ \SI{}{kN}$.
    \end{enumerate}
    \tcblower
    \textbf{Solution détaillée :}
    \begin{enumerate}
        \item \textbf{Analyse des forces :} Poids $\vec{P}$ et Normale $\vec{N}$ (travail nul car $\perp$ au déplacement). Forces de frottement $\vec{f}$ (parallèles au déplacement, sens opposé). Les frottements sont des forces \textbf{non conservatrices} (dissipatives).
        \textbf{Conclusion :} $W_{\text{nc}} = W_{\text{frot}} \neq 0 \Rightarrow \Delta E_m \neq 0$. \textbf{L'énergie mécanique n'est PAS conservée.} Elle diminue.
        \item \textbf{Bilan énergétique :}
        État initial : $E_{c,i} = \frac{1}{2} M v_0^2 = 0,5 \times 12000 \times (25,0)^2 = 6000 \times 625 = 3,75 \times 10^6\ \SI{}{J}$. $E_{p,i} = \text{cste}$ (route plane).
        État final : $v_f = 0 \Rightarrow E_{c,f} = 0$. $E_{p,f} = E_{p,i}$.
        $\Delta E_m = \Delta E_c = 0 - 3,75 \times 10^6 = -3,75 \times 10^6\ \SI{}{J}$.
        Travail des forces non conservatrices : $W_{\text{nc}} = W_{\text{frot}} = -3,75 \times 10^6\ \SI{}{J}$.
        \textbf{Vérification :} $\Delta E_m = W_{\text{nc}} = -3,75 \times 10^6\ \SI{}{J}$. \textbf{Relation vérifiée.}
        \item \textbf{Distance de freinage :} $W_{\text{frot}} = - F \times d$ (force constante, angle $\pi$).
        \[ -3,75 \times 10^6 = - (150 \times 10^3) \times d \]
        \[ d = \frac{3,75 \times 10^6}{1,5 \times 10^5} = 25,0\ \SI{}{m} \]
        $\boxed{d = 25,0\ \SI{}{m}}$.
    \end{enumerate}
    \textbf{Conclusion :} L'énergie mécanique perdue ($3,75\ \SI{}{MJ}$) est transformée en énergie thermique (échauffement des freins, turbulence d'air) par le travail des forces de frottement.
\end{exoresolu}

\begin{methode}[Exprimer $\Delta E_m$ par le travail des forces non conservatrices]
    \textbf{Objectif :} Appliquer la relation fondamentale $\Delta E_m = W_{\text{nc}}$ (Théorème de l'énergie mécanique).
    \begin{enumerate}
        \item \textbf{Définir le système} matériels (ex: \{Bloc\}, \{Bloc + Terre + Ressort\}).
        \item \textbf{Lister les forces extérieures} et identifier celles qui sont \textbf{non conservatrices} (frottements $\vec{f}$, force motrice $\vec{F}$, force de freinage, tension d'un câble tracteur...). Ne pas oublier : le poids et la force de ressort sont \emph{conservatives} (leur travail est déjà dans $\Delta E_p$).
        \item \textbf{Calculer le travail $W_k$} de chaque force non conservative $k$ entre l'état initial (i) et final (f) : $W_k = \int \vec{F}_k \cdot d\vec{l}$. Si force constante : $W = F \Delta l \cos\theta$.
        \item \textbf{Sommer} : $W_{\text{nc}} = \sum W_k$.
        \item \textbf{Écrire le bilan} : $E_{m,f} - E_{m,i} = W_{\text{nc}}$.
        \item \textbf{Déduire l'inconnue} (vitesse finale, hauteur max, compression ressort, coefficient de frottement...).
    \end{enumerate}
    \begin{aretenir}[Théorème de l'énergie mécanique (Formule générale)]
        \[ \Delta E_m = E_{m,f} - E_{m,i} = W_{\text{ext, nc}} = \sum_{\text{forces non cons.}} W_k \]
        \emph{Cas particulier :} Si $W_{\text{nc}} = 0 \Rightarrow \Delta E_m = 0$ (Conservation).
    \end{aretenir}
\end{methode}

\begin{exoresolu}[Remontée mécanique à la station de ski de Ngaoundéré (projet)]
    \textbf{Énoncé :} Un télésiège tracte une benne (masse $m = 300\ \SI{}{kg}$, 4 personnes) sur une pente inclinée à $\alpha = 30^\circ$. Le câble exerce une force de traction $\vec{T}$ parallèle à la pente, de norme constante $T = 2,5\ \SI{}{kN}$. Le frottement sur le câble porteur est modélisé par une force constante $f = 800\ \SI{}{N}$ opposée au mouvement. La benne part du repos ($v_i=0$) et parcourt une distance $L = 500\ \SI{}{m}$ le long du câble.
    \begin{enumerate}
        \item Faire le bilan des forces non conservatrices et calculer leur travail total $W_{\text{nc}}$.
        \item En déduire la vitesse $v_f$ de la benne à l'arrivée.
    \end{enumerate}
    \tcblower
    \textbf{Solution détaillée :}
    \begin{enumerate}
        \item \textbf{Système :} \{Benne + Terre\}. Référentiel : sol (galiléen).
        \textbf{Forces extérieures :} Poids $\vec{P}$ (conservative, $\Delta E_p = mg\Delta z$), Normale $\vec{N}$ (travail nul), Traction $\vec{T}$ (non conservative, moteur), Frottement $\vec{f}$ (non conservative).
        \textbf{Travaux des forces non conservatrices :}
        \begin{itemize}
            \item Traction $\vec{T}$ : même sens que déplacement. $W_T = + T \times L = 2500 \times 500 = 1,25 \times 10^6\ \SI{}{J}$.
            \item Frottement $\vec{f}$ : sens opposé. $W_f = - f \times L = -800 \times 500 = -4,00 \times 10^5\ \SI{}{J}$.
        \end{itemize}
        $W_{\text{nc}} = W_T + W_f = 1,25 \times 10^6 - 0,40 \times 10^6 = 8,50 \times 10^5\ \SI{}{J}$.
        $\boxed{W_{\text{nc}} = 850\ \SI{}{kJ}}$.
        \item \textbf{Application du théorème :} $\Delta E_m = W_{\text{nc}}$.
        $\Delta E_m = \Delta E_c + \Delta E_p$.
        $\Delta E_c = \frac{1}{2} m v_f^2 - 0 = 150 v_f^2$.
        $\Delta E_p = m g \Delta z$. $\Delta z = L \sin\alpha = 500 \times \sin 30^\circ = 500 \times 0,5 = 250\ \SI{}{m}$.
        $\Delta E_p = 300 \times 9,81 \times 250 = 735\,750\ \SI{}{J} \approx 7,36 \times 10^5\ \SI{}{J}$.
        \textbf{Équation :}
        \[ 150 v_f^2 + 7,3575 \times 10^5 = 8,50 \times 10^5 \]
        \[ 150 v_f^2 = 1,1425 \times 10^5 \]
        \[ v_f^2 = \frac{114250}{150} \approx 761,67 \]
        \[ v_f = \sqrt{761,67} \approx 27,6\ \SI{}{m/s} \]
        $\boxed{v_f \approx 27,6\ \SI{}{m/s}\ (\approx 99\ \SI{}{km/h})}$.
        \item \textbf{Vérification ordre de grandeur :} Vitesse élevée pour un télésiège (réel $\approx 5-6\ \SI{}{m/s}$), mais les chiffres de l'exercice (force traction très forte, frottement faible) mènent à ce résultat. L'important est la méthode.
    \end{enumerate}
    \textbf{Conclusion :} Le travail net des forces non conservatrices (moteur - frottements) a permis d'augmenter l'énergie cinétique ET l'énergie potentielle de la benne.
\end{exoresolu}

\begin{methode}[Résolution complète : Système avec ressort, pesanteur et frottements]
    \textbf{Objectif :} Traiter un problème complet mélangeant $E_{p,\text{pes}}$, $E_{p,\text{el}}$, $E_c$ et frottements.
    \begin{enumerate}
        \item \textbf{Schéma} : Représenter le système, les forces, le repère ($x$ ou $z$), l'origine des énergies potentielles ($z=0$, $x=0$).
        \item \textbf{Choisir les états} : Initial (i) et Final (f) (souvent : départ / arrivée / point haut / compression max).
        \item \textbf{Écrire $E_{m,i}$ et $E_{m,f}$} explicitement : $\frac{1}{2}m v_i^2 + m g z_i + \frac{1}{2}k x_i^2$ et idem pour $f$.
        \item \textbf{Calculer $W_{\text{nc}}$} (seulement frottements solides/fluides, forces motrices/freinantes extérieures). \textbf{Ne pas inclure} le travail du poids ni du ressort.
        \item \textbf{Appliquer} : $E_{m,f} - E_{m,i} = W_{\text{nc}}$.
        \item \textbf{Résoudre} l'équation (souvent du second degré en $v$ ou $x$). Garder la solution physiquement acceptable ($v^2 \ge 0$, $x$ dans le bon sens).
        \item \textbf{Phrase de conclusion} avec unité.
    \end{enumerate}
    \begin{aretenir}[Check-list Probatoire]
        \begin{itemize}
            \item Système défini ? Repère choisi ? Origine $E_p$ fixée ?
            \item Forces non conservatrices identifiées ? Travaux calculés (signe !) ?
            \item Équation $\Delta E_m = W_{\text{nc}}$ écrite ?
            \item Calculs numériques avec unités SI ?
            \item Réponse encadrée, unitée, commentée ?
        \end{itemize}
    \end{aretenir}
\end{methode}

\begin{exoresolu}[Bloc sur plan incliné rugueux avec butée à ressort]
    \textbf{Énoncé :} Un bloc de masse $m = 2,0\ \SI{}{kg}$ est lâché sans vitesse initiale ($v_A=0$) au point A d'un plan incliné à $\alpha = 30^\circ$. Il glisse sur une distance $L = 1,00\ \SI{}{m}$ jusqu'au point B où il rencontre un ressort de raideur $k = 500\ \SI{}{N/m}$ (longueur à vide $L_0$). Le coefficient de frottement solide est $\mu = 0,20$. Le bloc comprime le ressort jusqu'à l'arrêt momentané en C. On prend $g = 9,81\ \SI{}{N/kg}$.
    \begin{enumerate}
        \item Choisir l'origine de l'énergie potentielle de pesanteur et écrire $E_{m,A}$.
        \item Exprimer le travail des frottements $W_{\text{frot}}$ entre A et C (compression max $x_{\max}$ inconnue).
        \item Écrire l'équation $\Delta E_m = W_{\text{nc}}$ entre A et C et déterminer $x_{\max}$.
    \end{enumerate}
    \tcblower
    \textbf{Solution détaillée :}
    \begin{enumerate}
        \item \textbf{Choix :} Niveau de référence $z=0$ au niveau de B (contact initial ressort).
        État A : $v_A=0 \Rightarrow E_{c,A}=0$. $z_A = L \sin\alpha = 1,00 \times 0,5 = 0,50\ \SI{}{m}$.
        $E_{p,A} = m g z_A = 2,0 \times 9,81 \times 0,50 = 9,81\ \SI{}{J}$.
        Ressort non déformé $\Rightarrow E_{p,\text{el},A} = 0$.
        $\boxed{E_{m,A} = 9,81\ \SI{}{J}}$.
        \item \textbf{Travail des frottements (force non conservative) :}
        Force de frottement $f = \mu N = \mu m g \cos\alpha$.
        $f = 0,20 \times 2,0 \times 9,81 \times \cos 30^\circ = 0,20 \times 19,62 \times 0,866 \approx 3,40\ \SI{}{N}$.
        Distance parcourue de A à C : $L + x_{\max}$.
        Travail (force opposée au déplacement) : $W_{\text{frot}} = - f (L + x_{\max}) = -3,40 (1,00 + x_{\max})$.
        \item \textbf{État C (compression max) :} $v_C = 0 \Rightarrow E_{c,C} = 0$.
        $z_C = - x_{\max} \sin\alpha = -0,50 x_{\max}$ (sous le niveau B).
        $E_{p,C} = m g z_C = 2,0 \times 9,81 \times (-0,50 x_{\max}) = -9,81 x_{\max}$.
        $E_{p,\text{el},C} = \frac{1}{2} k x_{\max}^2 = 250 x_{\max}^2$.
        $E_{m,C} = 250 x_{\max}^2 - 9,81 x_{\max}$.
        \textbf{Théorème :} $E_{m,C} - E_{m,A} = W_{\text{frot}}$.
        \[ (250 x_{\max}^2 - 9,81 x_{\max}) - 9,81 = -3,40 (1,00 + x_{\max}) \]
        \[ 250 x_{\max}^2 - 9,81 x_{\max} - 9,81 = -3,40 - 3,40 x_{\max} \]
        \[ 250 x_{\max}^2 - 6,41 x_{\max} - 6,41 = 0 \]
        \textbf{Résolution :} $\Delta = (-6,41)^2 - 4 \times 250 \times (-6,41) = 41,09 + 6410 = 6451,09$.
        $\sqrt{\Delta} \approx 80,32$.
        $x_{\max} = \frac{6,41 + 80,32}{2 \times 250} = \frac{86,73}{500} = 0,1735\ \SI{}{m}$ (racine positive).
        $\boxed{x_{\max} \approx 0,174\ \SI{}{m}\ (17,4\ \SI{}{cm})}$.
        \item \textbf{Vérification énergétique :}
        Énergie initiale $E_{m,A} = 9,81\ \SI{}{J}$.
        Énergie finale stockée dans le ressort $E_{p,\text{el},C} \approx 250 \times (0,174)^2 \approx 7,57\ \SI{}{J}$.
        Travail des frottements $|W_{\text{frot}}| \approx 3,40 \times 1,174 \approx 3,99\ \SI{}{J}$.
        Variation énergie potentielle pesanteur $\Delta E_p = E_{p,C} - E_{p,A} = -9,81 \times 0,174 - 9,81 \approx -1,71 - 9,81 = -11,52\ \SI{}{J}$.
        Bilan : Perte $E_p$ ($11,52\ \SI{}{J}$) = Gain $E_{p,\text{el}}$ ($7,57\ \SI{}{J}$) + Énergie dissipée ($3,99\ \SI{}{J}$). $11,56 \approx 11,52$. \textbf{OK}.
    \end{enumerate}
    \textbf{Conclusion :} Le ressort se compresse de $17,4\ \SI{}{cm}$. L'énergie initiale de pesanteur a été dissipée par les frottements et stockée dans le ressort (une partie ayant aussi servi à abaisser le centre de gravité sous le niveau B).
\end{exoresolu}

\begin{attention}[Les 5 erreurs qui coûtent des points au Probatoire]
    \begin{enumerate}
        \item \textbf{Oublier de définir l'origine de l'énergie potentielle de pesanteur ($z=0$).} Sans cela, $E_p = mgz$ n'a pas de sens. Écrire explicitement : "Je choisis $z=0$ au niveau du sol / au point B / au niveau de la mer".
        \item \textbf{Confondre "Énergie mécanique conservée" et "Énergie cinétique conservée".} Dire "Comme il n'y a pas de frottement, $E_c$ est constante" est \textbf{FAUX} (elle s'échange avec $E_p$). Dire "$E_m$ est constante" est \textbf{VRAI}. Toujours écrire $E_{m,i} = E_{m,f}$.
        \item \textbf{Inclure le travail du poids ou du ressort dans $W_{\text{nc}}$.} Le théorème $\Delta E_m = W_{\text{nc}}$ ne concerne que les forces \emph{non conservatrices} (frottements, moteurs, freins, chocs). Le travail du poids est déjà compté dans $\Delta E_p = -W_{\text{poids}}$. Double comptage = 0 point.
        \item \textbf{Erreur de signe sur le travail des frottements.} Frottement $\vec{f}$ toujours opposé au déplacement $\vec{dl} \Rightarrow W_f = -f \times d$ (négatif). Oublier le moins fait gagner de l'énergie au système au lieu d'en perdre $\rightarrow$ résultat absurde ($v^2 < 0$).
        \item \textbf{Négliger l'énergie potentielle élastique dans le bilan final.} Si un ressort est comprimé à l'état final, $E_{p,\text{el},f} = \frac{1}{2}k x^2 \neq 0$. L'oublier fausse tout le calcul de $v_f$ ou $x_{\max}$. Vérifier systématiquement : "Y a-t-il un ressort déformé à l'état final ?"
    \end{enumerate}
\end{attention}

\newpage

\coursec{Exercices}

\courssub{Je vérifie mes connaissances}

\begin{exercice}[QCM : Énergies potentielles et forces conservatrices]
Parmi les affirmations suivantes, une seule est correcte. Laquelle ? Justifiez brièvement votre choix.
\begin{enumerate}
    \item L'énergie potentielle de pesanteur d'un système est toujours positive.
    \item L'énergie potentielle élastique d'un ressort est nulle lorsque le ressort est à sa longueur naturelle.
    \item Une force de frottement sec est une force conservative.
    \item L'énergie mécanique d'un système est conservée en présence de forces de frottement fluides.
\end{enumerate}
\end{exercice}

\begin{exercice}[Vrai ou Faux : Conservation de l'énergie mécanique]
Dire si les affirmations suivantes sont vraies ou fausses. Justifiez chaque réponse par une phrase ou une formule.
\begin{enumerate}
    \item Un mobile soumis uniquement à son poids et à la réaction normale d'un plan incliné parfait conserve son énergie mécanique.
    \item L'énergie potentielle de pesanteur $E_p = mgh$ ne dépend que de la masse $m$ et de l'altitude $h$, pas du chemin parcouru.
    \item Si l'énergie mécanique d'un système diminue, c'est nécessairement que l'énergie cinétique a diminué.
    \item Le travail d'une force conservative dépend uniquement des positions initiale et finale, pas de la trajectoire.
\end{enumerate}
\end{exercice}

\begin{exercice}[Calculs directs : Expressions des énergies potentielles]
\begin{enumerate}
    \item Un bloc de masse $m = 2,5\,\text{kg}$ est hissé à une hauteur $h = 12\,\text{m}$ au-dessus du sol (référence $h=0$). Calculer son énergie potentielle de pesanteur $E_p$. Prendre $g = 9,81\,\text{N/kg}$.
    \item Un ressort de raideur $k = 250\,\text{N/m}$ est comprimé de $x = 8,0\,\text{cm}$ par rapport à sa position d'équilibre. Calculer l'énergie potentielle élastique $E_{p,\text{el}}$ stockée dans le ressort.
    \item Une bille de $m = 50\,\text{g}$ se déplace à la vitesse $v = 4,0\,\text{m/s}$ au sommet d'une bosse de hauteur $h = 1,5\,\text{m}$ (référence au sol). Calculer son énergie mécanique $E_m$ à cet instant.
\end{enumerate}
\end{exercice}

\begin{exercice}[Définitions et relations fondamentales]
\begin{enumerate}
    \item Donnez la définition de l'énergie mécanique $E_m$ d'un système. Précisez les grandeurs qui la composent.
    \item Énoncez la relation liant la variation d'énergie cinétique $\Delta E_c$ et la variation d'énergie potentielle $\Delta E_p$ pour un système soumis uniquement à des forces conservatrices.
    \item Quelle est la condition nécessaire et suffisante pour que l'énergie mécanique d'un système soit conservée au cours de son évolution ?
    \item Exprimez le travail $W_{\text{nc}}$ des forces non conservatrices en fonction de la variation de l'énergie mécanique $\Delta E_m$.
\end{enumerate}
\end{exercice}

\courssub{Je m'entraîne}

\begin{exercice}[Descente sur plan incliné : la côte de Buea]
Un camion de masse $m = 8,0\,\text{t}$ ($8000\,\text{kg}$) stationne au sommet d'une portion rectiligne de la côte de Buea, d'altitude $h_A = 120\,\text{m}$ par rapport au bas de la pente (point B, référence $h=0$). Le camion est lâché sans vitesse initiale ($v_A = 0$). La longueur de la pente est $L = 800\,\text{m}$. On prend $g = 9,81\,\text{N/kg}$.
\begin{enumerate}
    \item \textbf{Sans frottements :} En supposant la route parfaitement lisse, calculez la vitesse $v_B$ du camion au point B.
    \item \textbf{Avec frottements :} En réalité, une force de frottement totale (frottements internes + résistance de l'air) de norme constante $f = 10\,\text{kN}$ s'oppose au mouvement. Calculez la vitesse réelle $v'_B$ au point B.
    \item Quelle quantité d'énergie (en MJ) a été dissipée par les frottements pendant la descente ?
\end{enumerate}
\end{exercice}

\begin{exercice}[Lanceur à ressort : propulsion d'une bille]
Un ressort horizontal de raideur $k = 150\,\text{N/m}$, fixé à un mur, est comprimé de $x_0 = 20\,\text{cm}$. Une bille de masse $m = 40\,\text{g}$ est placée contre le ressort. L'ensemble est sur un plan horizontal parfaitement lisse. Le ressort est relâché.
\begin{enumerate}
    \item Calculez l'énergie potentielle élastique initiale stockée dans le ressort.
    \item En négligeant les frottements, déterminez la vitesse $v_{\text{max}}$ de la bille au moment où elle quitte le ressort (position détendue).
    \item La bille rencontre ensuite une zone rugueuse de longueur $d = 50\,\text{cm}$ où le coefficient de frottement dynamique vaut $\mu = 0,15$. Calculez sa vitesse $v_{\text{sortie}}$ en sortie de cette zone.
    \item La bille gravit ensuite une pente lisse inclinée à $\alpha = 30^\circ$. Quelle hauteur maximale $h_{\text{max}}$ atteint-elle par rapport au plan horizontal ? Quelle distance $D$ parcourt-elle sur la pente ?
\end{enumerate}
\end{exercice}

\begin{exercice}[Pendule simple : étude énergétique]
Un pendule simple est constitué d'une masse $m = 200\,\text{g}$ attachée à un fil inextensible de longueur $L = 1,20\,\text{m}$, sans masse. On écarte la masse de la verticale d'un angle $\theta_0 = 30^\circ$ et on la lâche sans vitesse initiale ($v_0 = 0$). Le point le plus bas de la trajectoire (point B) est choisi comme référence d'altitude ($h=0$). On prend $g = 9,81\,\text{N/kg}$.
\begin{enumerate}
    \item Calculez la hauteur $h_A$ du point de lâcher A par rapport au point B.
    \item Calculez l'énergie mécanique $E_m$ du système en A. Est-elle conservée au cours du mouvement ? Justifiez.
    \item Déterminez la vitesse $v_B$ de la masse au passage au point le plus bas B.
    \item Calculez la vitesse $v_C$ de la masse lorsqu'elle passe par un point C tel que le fil forme un angle $\theta = 15^\circ$ avec la verticale.
    \item Tracez schématiquement l'allure des courbes $E_c(\theta)$, $E_p(\theta)$ et $E_m(\theta)$ pour un demi-période (de $\theta_0$ à $-\theta_0$).
\end{enumerate}
\end{exercice}

\begin{exercice}[Exploitation graphique : Énergies en fonction de l'abscisse]
On étudie le mouvement d'un mobile de masse $m = 0,50\,\text{kg}$ sur un axe horizontal $Ox$. Les courbes ci-dessous représentent l'évolution de l'énergie cinétique $E_c(x)$ et de l'énergie potentielle $E_p(x)$ (en joules) en fonction de l'abscisse $x$ (en mètres). L'énergie potentielle est de type élastique : $E_p(x) = \frac{1}{2}k(x-x_{\text{eq}})^2$.
\begin{popfigure}
\begin{tikzpicture}
\begin{axis}[
    axis lines=middle,
    xlabel={$x\,\text{(m)}$}, ylabel={$E\,\text{(J)}$},
    xmin=0, xmax=10.5, ymin=-0.5, ymax=5.5,
    xtick={0,2,4,5,6,8,10}, ytick={0,1,2,3,4,5},
    width=\linewidth, height=7cm,
    legend style={at={(0.98,0.98)}, anchor=north east, fill=white},
    samples=200,
    domain=0:10
]
\addplot[popblue, thick] {0.2*(x-5)^2};
\addlegendentry{$E_p(x)$}
\addplot[poporange, thick] {5 - 0.2*(x-5)^2};
\addlegendentry{$E_c(x)$}
\addplot[popgreen, dashed, thick] {5};
\addlegendentry{$E_m$}
\end{axis}
\end{tikzpicture}
\legende{Énergies cinétique, potentielle (élastique) et mécanique en fonction de l'abscisse.}
\end{popfigure}
\begin{enumerate}
    \item Quelle est la valeur de l'énergie mécanique $E_m$ du mobile ? Est-elle conservée ? Justifiez en vous basant sur le graphique.
    \item Quelle est la vitesse maximale $v_{\text{max}}$ atteinte par le mobile ? À quelle(s) position(s) $x$ est-elle atteinte ?
    \item Déterminez les positions $x$ où la vitesse du mobile est nulle (points de retour).
    \item Quelle est la nature des forces agissant sur le mobile ? (Conservatrices / Non conservatrices). Justifiez.
    \item Déduisez la raideur $k$ du ressort et la position d'équilibre $x_{\text{eq}}$.
\end{enumerate}
\end{exercice}

\courssub{Je me prépare au Probatoire}

\begin{exercice}[Freinage d'un poids lourd : la descente de la côte de Buea]
\label{exo:probatoire_camion}
Un camion-citerne de masse totale $M = 12\,\text{t}$ ($12\,000\,\text{kg}$) descend une portion rectiligne de la côte de Buea, inclinée d'un angle $\alpha = 12^\circ$ par rapport à l'horizontale. Le conducteur souhaite maintenir une vitesse constante $v = 60\,\text{km/h}$ ($16,7\,\text{m/s}$) sur une longueur $L = 2,5\,\text{km}$. Les forces de frottement de roulement et de l'air (forces non conservatrices hors freinage) sont modélisées par une force constante de norme $f_{\text{ext}} = 8,0\,\text{kN}$, dirigée vers le haut de la pente. On prend $g = 9,81\,\text{N/kg}$.
\begin{enumerate}
    \item Faites le bilan des forces agissant sur le camion (poids, normale, frottements externes, force de freinage $F_{\text{fr}}$ exercée par les freins sur les roues). Projetez les forces sur l'axe de la pente (orienté vers le bas).
    \item Calculez la composante du poids suivant la pente.
    \item En appliquant le théorème de l'énergie cinétique (ou le principe de la variation de l'énergie mécanique), déterminez la norme de la force de freinage $F_{\text{fr}}$ que le conducteur doit appliquer pour maintenir la vitesse constante.
    \item Calculez le travail $W_{\text{fr}}$ de la force de freinage sur la distance $L$. Quelle est l'énergie dissipée (échauffement des freins) en MJ ?
    \item Si le système de freinage a une masse totale $m_{\text{fr}} = 80\,\text{kg}$ (disques + étriers) et une capacité thermique massique $c = 460\,\text{J/(kg$\cdot$K)}$ (acier), et si $80\%$ de l'énergie dissipée chauffe les freins, estimez l'élévation de température $\Delta T$ des freins sur cette portion. Commenter le risque de \emph{fading} (perte d'efficacité par surchauffe).
\end{enumerate}
\end{exercice}

\begin{exercice}[Saut à la perche : conversion d'énergie et record]
\label{exo:probatoire_perche}
Un perchiste camerounais de masse $m = 75\,\text{kg}$ court sur la piste d'élan avec une vitesse $v_0 = 9,5\,\text{m/s}$. Il plante la perche dans la boîte et réalise un saut. On modélise le sauteur comme un point matériel. On néglige dans un premier temps les pertes d'énergie (frottements air, déformation perche, travail musculaire pendant l'envol). On prend $g = 9,81\,\text{N/kg}$.
\begin{enumerate}
    \item En supposant une conversion intégrale de l'énergie cinétique de translation en énergie potentielle de pesanteur au point le plus haut, calculez la hauteur théorique maximale $H_{\text{th}}$ que le centre de gravité (CDG) du sauteur peut atteindre \emph{par rapport à sa position au sol} (référence $h=0$ au niveau du sol).
    \item En réalité, le CDG du sauteur est à $h_{\text{CDG},0} = 1,00\,\text{m}$ du sol en position de course. Quelle est la hauteur théorique maximale $h_{\text{max,th}}$ de la barre qu'il peut franchir (hauteur du CDG au sommet) ?
    \item Le record du monde masculin (Armand Duplantis, 2024) est de $6,24\,\text{m}$. Le record du Cameroun est d'environ $5,00\,\text{m}$. En supposant une vitesse d'élan identique ($9,5\,\text{m/s}$), calculez le "déficit" de hauteur $\Delta h = h_{\text{max,th}} - h_{\text{record}}$ pour le record du monde. Ce déficit est-il positif ? Interprétez physiquement ce résultat (apport énergétique du sauteur sur la perche, effet de la perche, travail musculaire).
    \item La perche agit comme un ressort. Si l'on modélise la perche par un ressort de raideur $k = 2\,500\,\text{N/m}$ comprimé au maximum de $x_{\text{max}} = 1,80\,\text{m}$ lors de l'envol, calculez l'énergie potentielle élastique stockée $E_{p,\text{el}}$. Comparez cette énergie à l'énergie cinétique initiale $E_{c,0}$. Que représente la différence ?
    \item En tenant compte du travail musculaire $W_{\text{mus}}$ fourni par le sauteur pendant la phase de relevé (poussée sur la perche), écrire le bilan énergétique complet entre l'instant de la prise d'appel (vitesse $v_0$, CDG à $1\,\text{m}$) et l'instant du franchissement de la barre (vitesse quasi nulle, CDG à $h_{\text{barre}}$). Exprimez $W_{\text{mus}}$ en fonction des grandeurs connues. Estimez $W_{\text{mus}}$ pour franchir $6,00\,\text{m}$.
\end{enumerate}
\end{exercice}

\begin{exercice}[Centrale hydroélectrique de Nachtigal : de l'énergie potentielle à l'électricité]
\label{exo:probatoire_nachtigal}
La centrale hydroélectrique de Nachtigal (sur la Sanaga) utilise une chute d'eau brute de hauteur $H = 35,5\,\text{m}$. Le débit turbiné maximal est $Q = 350\,\text{m}^3/\text{s}$. La masse volumique de l'eau est $\rho = 1000\,\text{kg/m}^3$. On prend $g = 9,81\,\text{N/kg}$.
\begin{enumerate}
    \item Calculez la masse d'eau $m$ qui traverse les turbines en une seconde (débit massique $\dot{m}$).
    \item Calculez la puissance hydraulique brute $P_{\text{hyd}}$ disponible (puissance de l'énergie potentielle de pesanteur convertie par seconde).
    \item Le rendement global de la chaîne de conversion (turbine + alternateur + transformateur) est $\eta = 0,92$. Calculez la puissance électrique $P_{\text{el}}$ injectée sur le réseau SONATREL/ENEO.
    \item L'énergie électrique produite en une journée de fonctionnement à pleine charge est vendue au tarif moyen de $50\,\text{FCFA/kWh}$. Calculez le chiffre d'affaires journalier (en FCFA et en millions de FCFA).
    \item \textbf{Analyse énergétique détaillée :} L'eau arrive en amont avec une vitesse $v_{\text{amont}} = 2,0\,\text{m/s}$ dans le canal d'amenée. En aval de la turbine, l'eau est rejetée dans le canal de fuite avec une vitesse $v_{\text{aval}} = 3,5\,\text{m/s}$ à la pression atmosphérique. En appliquant le théorème de Bernoulli (ou bilan d'énergie mécanique sur un fluide) entre l'amont et l'aval de la turbine, exprimez le travail spécifique $w_t$ (en J/kg) extrait par la turbine sur l'eau. Vérifiez que la puissance mécanique turbine $P_{\text{meca}} = \dot{m} w_t$ est cohérente avec $P_{\text{hyd}}$ (en négligeant les pertes de charge dans les canaux).
    \item Si la centrale fonctionne $8\,000\,\text{h/an}$, quelle est l'énergie annuelle produite en GWh ? Combien de foyers camerounais (consommation moyenne $150\,\text{kWh/mois}$) peut-elle alimenter ?
\end{enumerate}
\end{exercice}



\newpage

\coursec{Corrigés des exercices}

\begin{corrige}[Exercice 1 — QCM : Énergies potentielles et forces conservatrices]
\emph{Analyse des propositions :}
\begin{enumerate}
    \item \textbf{Fausse.} L'énergie potentielle de pesanteur $E_p = mgh$ dépend du choix de l'origine des altitudes ($h=0$). Elle est positive si le système est au-dessus de la référence, négative en dessous. \cle{Énergie potentielle de pesanteur} n'est pas une valeur absolue.
    \item \textbf{Vraie.} Par définition, $E_{p,\text{el}} = \frac{1}{2}k x^2$ où $x$ est l'allongement (ou compression) par rapport à la longueur naturelle. Pour $x=0$, $E_{p,\text{el}}=0$. C'est le minimum de l'énergie potentielle élastique. \cle{Énergie potentielle élastique}
    \item \textbf{Fausse.} Les forces de frottement (sec ou fluide) sont des forces \emph{non conservatrices} (ou dissipatives) : leur travail dépend du chemin parcouru et est toujours négatif (ou nul), il ne peut pas être récupéré sous forme d'énergie mécanique. \cle{Force non conservative}
    \item \textbf{Fausse.} Les forces de frottement fluides sont non conservatrices ; leur travail diminue l'énergie mécanique du système ($\Delta E_m = W_{\text{nc}} < 0$).
\end{enumerate}
La bonne réponse est la \textbf{proposition 2}.
\end{corrige}

\begin{corrige}[Exercice 2 — Vrai ou Faux : Conservation de l'énergie mécanique]
\begin{enumerate}
    \item \textbf{Vrai.} Le poids est une force \cle{conservative}. La réaction normale est perpendiculaire à la vitesse (donc au déplacement élémentaire), son travail est nul. Le travail des forces non conservatrices est nul $\Rightarrow \Delta E_m = 0$.
    \item \textbf{Vrai.} L'énergie potentielle de pesanteur est une \emph{fonction d'état} : elle ne dépend que de la position finale (altitude $h$) par rapport à la référence, non de l'historique du mouvement.
    \item \textbf{Faux.} $\Delta E_m = \Delta E_c + \Delta E_p$. Si $E_m$ diminue ($\Delta E_m < 0$), on peut avoir $\Delta E_c > 0$ (augmentation de $E_c$) si $\Delta E_p$ est négative et plus grande en valeur absolue (ex: chute avec frottements : $E_p$ baisse vite, $E_c$ augmente mais moins vite, $E_m$ baisse).
    \item \textbf{Vrai.} C'est la définition même d'une force conservative : $W_{A\to B} = -\Delta E_p = E_{p,A} - E_{p,B}$, indépendant du chemin. \cle{Travail d'une force conservative}
\end{enumerate}
\end{corrige}

\begin{corrige}[Exercice 3 — Calculs directs : Expressions des énergies potentielles]
\begin{enumerate}
    \item $E_p = mgh = 2,5 \times 9,81 \times 12 = \SI{294,3}{J}$.
    \item $x = 8,0\,\text{cm} = 0,080\,\text{m}$. $E_{p,\text{el}} = \frac{1}{2}k x^2 = 0,5 \times 250 \times (0,080)^2 = \SI{0,80}{J}$.
    \item $m = 0,050\,\text{kg}$.
    $E_c = \frac{1}{2}mv^2 = 0,5 \times 0,050 \times 4,0^2 = \SI{0,40}{J}$.
    $E_p = mgh = 0,050 \times 9,81 \times 1,5 = \SI{0,736}{J}$.
    $E_m = E_c + E_p = 0,40 + 0,736 = \SI{1,14}{J}$ (arrondi à $10^{-2}$ près).
\end{enumerate}
\end{corrige}

\begin{corrige}[Exercice 4 — Définitions et relations fondamentales]
\begin{enumerate}
    \item L'énergie mécanique $E_m$ d'un système est la somme de son énergie cinétique $E_c$ et de son énergie potentielle $E_p$ : $E_m = E_c + E_p$. $E_p$ regroupe l'énergie potentielle de pesanteur $E_{p,p}$ et l'énergie potentielle élastique $E_{p,\text{el}}$ (et toute autre énergie potentielle associée à une force conservative). \cle{Énergie mécanique}
    \item Pour un système soumis \emph{uniquement} à des forces conservatrices, le travail des forces non conservatrices est nul. D'après le théorème de l'énergie cinétique $\Delta E_c = W_{\text{cons}} + W_{\text{nc}}$ et $W_{\text{cons}} = -\Delta E_p$, on obtient : $\Delta E_c + \Delta E_p = 0$, soit $\Delta E_c = -\Delta E_p$. \cle{Relation $\Delta E_c = -\Delta E_p$} (exigible au Probatoire).
    \item L'énergie mécanique est conservée ($\Delta E_m = 0$) \textbf{si et seulement si} le travail des forces non conservatrices extérieures (et intérieures non conservatrices) est nul : $W_{\text{nc}} = 0$. \cle{Principe de conservation de l'énergie mécanique}
    \item Le théorème de la variation de l'énergie mécanique s'écrit : $\Delta E_m = W_{\text{nc}}$. Le travail des forces non conservatrices est égal à la variation de l'énergie mécanique. \cle{Théorème de la variation de l'énergie mécanique}
\end{enumerate}
\end{corrige}

\begin{corrige}[Exercice 5 — Descente sur plan incliné : la côte de Buea]
\begin{enumerate}
    \item \textbf{Sans frottements :} $W_{\text{nc}} = 0 \Rightarrow \Delta E_m = 0$.
    $E_{m,A} = E_{p,A} = mgh_A$ (car $v_A=0$).
    $E_{m,B} = E_{c,B} = \frac{1}{2}mv_B^2$ (car $h_B=0$).
    $mgh_A = \frac{1}{2}mv_B^2 \Rightarrow v_B = \sqrt{2gh_A} = \sqrt{2 \times 9,81 \times 120} = \sqrt{2354,4} \approx \SI{48,5}{m/s}$ (soit $\approx 175\,\text{km/h}$).
    \item \textbf{Avec frottements :} Forces non conservatrices : frottements $\vec{f}$ (opposée au déplacement). $W_{\text{nc}} = W_f = -fL$.
    $\Delta E_m = W_{\text{nc}} \Rightarrow \frac{1}{2}mv_B'^2 - mgh_A = -fL$.
    $\frac{1}{2}mv_B'^2 = mgh_A - fL$.
    $m = \SI{8000}{kg}$, $h_A = \SI{120}{m}$, $f = \SI{10\,000}{N}$, $L = \SI{800}{m}$.
    $mgh_A = 8000 \times 9,81 \times 120 = 9\,417\,600\,\text{J}$.
    $fL = 10\,000 \times 800 = 8\,000\,000\,\text{J}$.
    $E_{c,B} = 9\,417\,600 - 8\,000\,000 = 1\,417\,600\,\text{J}$.
    $v_B' = \sqrt{\frac{2 \times 1\,417\,600}{8000}} = \sqrt{354,4} \approx \SI{18,8}{m/s}$ (soit $\approx 67,8\,\text{km/h}$).
    \item Énergie dissipée = $|W_f| = fL = 8\,000\,000\,\text{J} = \SI{8,0}{MJ}$.
\end{enumerate}
\end{corrige}

\begin{corrige}[Exercice 6 — Lanceur à ressort : propulsion d'une bille]
\begin{enumerate}
    \item $x_0 = 0,20\,\text{m}$. $E_{p,\text{el},0} = \frac{1}{2}k x_0^2 = 0,5 \times 150 \times 0,20^2 = \SI{3,0}{J}$.
    \item Position détendue : $x=0 \Rightarrow E_{p,\text{el}}=0$. Conservation de $E_m$ (plan lisse horizontal, normale $\perp$ vitesse, poids $\perp$ vitesse).
    $E_{p,\text{el},0} = E_{c,\text{max}} \Rightarrow \frac{1}{2}mv_{\text{max}}^2 = 3,0$.
    $v_{\text{max}} = \sqrt{\frac{2 \times 3,0}{0,040}} = \sqrt{150} \approx \SI{12,2}{m/s}$.
    \item Zone rugueuse : forces non conservatrices = frottements $\vec{f}$, $f = \mu mg = 0,15 \times 0,040 \times 9,81 = \SI{0,0589}{N}$.
    $W_f = -f d = -0,0589 \times 0,50 = -\SI{0,0294}{J}$.
    $\Delta E_m = W_f \Rightarrow E_{c,\text{sortie}} - E_{c,\text{entrée}} = -0,0294$.
    $E_{c,\text{sortie}} = 3,0 - 0,0294 = \SI{2,971}{J}$.
    $v_{\text{sortie}} = \sqrt{\frac{2 \times 2,971}{0,040}} = \sqrt{148,5} \approx \SI{12,2}{m/s}$ (légèrement inférieure à $v_{\text{max}}$).
    \item Pente lisse : conservation de $E_m$ à partir de la sortie de la zone rugueuse (référence $h=0$ au plan horizontal).
    $E_{c,\text{sortie}} = mgh_{\text{max}} \Rightarrow h_{\text{max}} = \frac{E_{c,\text{sortie}}}{mg} = \frac{2,971}{0,040 \times 9,81} \approx \SI{7,57}{m}$.
    Distance sur la pente : $D = \frac{h_{\text{max}}}{\sin 30^\circ} = \frac{7,57}{0,5} = \SI{15,1}{m}$.
\end{enumerate}
\end{corrige}

\begin{corrige}[Exercice 7 — Pendule simple : étude énergétique]
\begin{enumerate}
    \item Géométrie : $h_A = L - L\cos\theta_0 = L(1 - \cos 30^\circ) = 1,20 \times (1 - \frac{\sqrt{3}}{2}) \approx 1,20 \times 0,134 = \SI{0,161}{m}$.
    \item En A : $v_A=0 \Rightarrow E_{c,A}=0$. $E_{p,A} = mgh_A = 0,200 \times 9,81 \times 0,161 \approx \SI{0,316}{J}$.
    $E_m = \SI{0,316}{J}$.
    \textbf{Conservation :} Oui. Forces : Poids (conservative) + Tension du fil. La tension est toujours perpendiculaire à la vitesse (mouvement circulaire), son travail est nul. $W_{\text{nc}}=0 \Rightarrow \Delta E_m=0$.
    \item En B : $h_B=0 \Rightarrow E_{p,B}=0$. $E_m = E_{c,B} = \frac{1}{2}mv_B^2$.
    $v_B = \sqrt{\frac{2E_m}{m}} = \sqrt{\frac{2 \times 0,316}{0,200}} = \sqrt{3,16} \approx \SI{1,78}{m/s}$.
    \item En C : $\theta = 15^\circ$. $h_C = L(1 - \cos 15^\circ) = 1,20 \times (1 - 0,9659) \approx \SI{0,0409}{m}$.
    $E_{p,C} = mgh_C = 0,200 \times 9,81 \times 0,0409 \approx \SI{0,0802}{J}$.
    $E_{c,C} = E_m - E_{p,C} = 0,316 - 0,0802 = \SI{0,236}{J}$.
    $v_C = \sqrt{\frac{2 \times 0,236}{0,200}} = \sqrt{2,36} \approx \SI{1,54}{m/s}$.
    \item \textbf{Allure des courbes (demi-période $\theta_0 \to -\theta_0$) :}
    \begin{itemize}
        \item $E_m(\theta)$ : ligne horizontale à $0,316\,\text{J}$.
        \item $E_p(\theta) = mgL(1-\cos\theta)$ : courbe en "U" symétrique, minimum $0$ en $\theta=0$, maximum $0,316\,\text{J}$ en $\theta=\pm 30^\circ$.
        \item $E_c(\theta) = E_m - E_p(\theta)$ : courbe en "n" symétrique, maximum $0,316\,\text{J}$ en $\theta=0$, nul en $\theta=\pm 30^\circ$.
    \end{itemize}
\end{enumerate}
\end{corrige}

\begin{corrige}[Exercice 8 — Exploitation graphique : Énergies en fonction de l'abscisse]
\begin{enumerate}
    \item Lecture graphique : $E_m = E_c(x) + E_p(x) = \SI{5,0}{J}$ (courbe verte pointillée horizontale). L'énergie mécanique est \textbf{conservée} car sa somme reste constante pour tout $x$ (la courbe $E_m$ est une ligne horizontale).
    \item $v_{\text{max}}$ correspond à $E_c$ maximum. Sur le graphique, $E_{c,\text{max}} = \SI{5,0}{J}$ (au sommet de la courbe orange, en $x=5\,\text{m}$).
    $v_{\text{max}} = \sqrt{\frac{2E_{c,\text{max}}}{m}} = \sqrt{\frac{2 \times 5,0}{0,50}} = \sqrt{20} \approx \SI{4,47}{m/s}$.
    Position : $x = \SI{5,0}{m}$ (position d'équilibre du ressort).
    \item $v=0 \Leftrightarrow E_c=0$. Lecture graphique : intersections de la courbe orange avec l'axe des abscisses : $x = \SI{0}{m}$ et $x = \SI{10}{m}$. Ce sont les points de retour (butées).
    \item Comme $E_m$ est constante ($W_{\text{nc}}=0$), les forces agissant sur le mobile (hors forces normales au déplacement) sont \textbf{conservatrices}. Ici, la force du ressort (élastique) est conservative. L'axe est horizontal $\Rightarrow$ travail du poids et de la normale nuls.
    \item Forme de $E_p(x) = \frac{1}{2}k(x-x_{\text{eq}})^2$.
    Le minimum de $E_p$ est en $x=5\,\text{m}$ ($E_p=0$) $\Rightarrow x_{\text{eq}} = \SI{5,0}{m}$.
    En $x=0$ (ou $x=10$), $E_p = \SI{5,0}{J} = \frac{1}{2}k(5)^2 = \frac{25}{2}k$.
    $k = \frac{2 \times 5,0}{25} = \SI{0,40}{N/m}$.
\end{enumerate}
\end{corrige}

\begin{corrige}[Exercice 9 — Freinage d'un poids lourd : la descente de la côte de Buea]
\begin{enumerate}
    \item \textbf{Bilan des forces (référentiel terrestre, axe $x'$ vers le bas de la pente) :}
    \begin{itemize}
        \item Poids $\vec{P} = M\vec{g}$ : composante $P_{x'} = Mg\sin\alpha$ (vers le bas).
        \item Réaction normale $\vec{N}$ : perpendiculaire à la pente, $N = Mg\cos\alpha$ (pas de projection sur $x'$).
        \item Frottements externes $\vec{f}_{\text{ext}}$ (air, roulement) : norme $f_{\text{ext}}$, dirigée vers le haut $\Rightarrow$ projection $-f_{\text{ext}}$.
        \item Force de freinage $\vec{F}_{\text{fr}}$ : force de frottement exercée par le sol sur les roues (vers le haut) due au couple de freinage, projection $-F_{\text{fr}}$.
    \end{itemize}
    \item $P_{x'} = Mg\sin\alpha = 12\,000 \times 9,81 \times \sin 12^\circ$.
    $\sin 12^\circ \approx 0,2079$.
    $P_{x'} \approx 12\,000 \times 9,81 \times 0,2079 \approx \SI{24\,470}{N}$ (soit $\SI{24,5}{kN}$).
    \item Vitesse constante $\Rightarrow a=0 \Rightarrow \Delta E_c = 0$.
    Théorème de l'énergie cinétique : $\Delta E_c = W_{\text{total}} = W_P + W_{f_{\text{ext}}} + W_{F_{\text{fr}}}$.
    $W_P = P_{x'} L$, $W_{f_{\text{ext}}} = -f_{\text{ext}} L$, $W_{F_{\text{fr}}} = -F_{\text{fr}} L$.
    $0 = (P_{x'} - f_{\text{ext}} - F_{\text{fr}}) L \Rightarrow F_{\text{fr}} = P_{x'} - f_{\text{ext}}$.
    $F_{\text{fr}} = 24\,470 - 8\,000 = \SI{16\,470}{N} \approx \SI{16,5}{kN}$.
    \item $W_{\text{fr}} = -F_{\text{fr}} L = -16\,470 \times 2\,500 = -41\,175\,000\,\text{J} \approx -\SI{41,2}{MJ}$.
    Énergie dissipée (chaleur) = $|W_{\text{fr}}| = \SI{41,2}{MJ}$.
    \item Énergie thermique reçue par les freins : $Q = 0,80 \times 41,2 \times 10^6 = \SI{32,96}{MJ}$.
    $Q = m_{\text{fr}} c \Delta T \Rightarrow \Delta T = \frac{Q}{m_{\text{fr}} c} = \frac{32,96 \times 10^6}{80 \times 460} = \frac{32,96 \times 10^6}{36\,800} \approx \SI{896}{K}$ (soit $\approx \SI{896}{^\circ C}$).
    \textbf{Commentaire :} Cette élévation de température est considérable (l'acier rougit vers $600-700^\circ\text{C}$). Il y a un risque majeur de \emph{fading} : perte d'adhérence plaquettes/disques, vaporisation de l'huile de frein (si freinage hydraulique), déformation thermique des disques. C'est pourquoi les poids lourds utilisent des freins moteurs (retardateurs) et des échappements freins pour limiter l'usage des freins de service sur les longues descentes.
\end{enumerate}
\end{corrige}

\begin{corrige}[Exercice 10 — Saut à la perche : conversion d'énergie et record]
\begin{enumerate}
    \item Conversion $E_{c,0} \to \Delta E_p$ : $\frac{1}{2}mv_0^2 = mgH_{\text{th}} \Rightarrow H_{\text{th}} = \frac{v_0^2}{2g} = \frac{9,5^2}{2 \times 9,81} = \frac{90,25}{19,62} \approx \SI{4,60}{m}$.
    \item $h_{\text{max,th}} = h_{\text{CDG},0} + H_{\text{th}} = 1,00 + 4,60 = \SI{5,60}{m}$.
    \item Pour le record du monde ($h_{\text{record}} = 6,24\,\text{m}$, A. Duplantis 2024) :
    $\Delta h = h_{\text{max,th}} - h_{\text{record}} = 5,60 - 6,24 = \SI{-0,64}{m}$.
    Le déficit est \textbf{négatif}. La hauteur théorique sans apport extérieur ($5,60\,\text{m}$) est inférieure à la barre franchie ($6,24\,\text{m}$).
    \textbf{Interprétation :} Le sauteur ne se contente pas de convertir son $E_c$ initiale. Il effectue un \textbf{travail musculaire} important pendant la phase de compression de la perche (relevé, poussée des bras/jambes) qui injecte de l'énergie supplémentaire dans le système "sauteur + perche". La perche, en se déformant, stocke cette énergie (cinétique + musculaire) sous forme élastique et la restitue pour propulser le sauteur plus haut. L'énergie finale $mg h_{\text{barre}}$ dépasse l'énergie cinétique initiale.
    \item $E_{p,\text{el}} = \frac{1}{2}k x_{\text{max}}^2 = 0,5 \times 2\,500 \times 1,80^2 = 1\,250 \times 3,24 = \SI{4\,050}{J}$.
    $E_{c,0} = \frac{1}{2}mv_0^2 = 0,5 \times 75 \times 90,25 = \SI{3\,384}{J}$.
    $E_{p,\text{el}} > E_{c,0}$. La différence $\Delta E = 4\,050 - 3\,384 = \SI{666}{J}$ correspond au \textbf{travail musculaire} effectué par le sauteur \emph{pendant la compression de la perche} (phase de "planté" et début de relevé) pour cintrer la perche au-delà de ce que permettrait sa seule énergie cinétique d'arrivée.
    \item \textbf{Bilan global (Système = Sauteur, système déformable) :}
    État initial (appel) : $E_{c,0} + mg h_0$.
    État final (barre) : $v_f \approx 0 \Rightarrow E_{c,f} \approx 0$ ; $E_{p,f} = mg h_{\text{barre}}$.
    Forces extérieures : Poids (conservative), Réaction de la perche (conservative, élastique), Forces musculaires (internes non conservatrices).
    Théorème de l'énergie mécanique pour système déformable : $\Delta E_m = W_{\text{mus}}$.
    $(mg h_{\text{barre}} - 0) - (0 + mg h_0) = W_{\text{mus}}$.
    $W_{\text{mus}} = mg(h_{\text{barre}} - h_0) - E_{c,0}$.
    Pour $h_{\text{barre}} = 6,00\,\text{m}$ :
    $W_{\text{mus}} = 75 \times 9,81 \times (6,00 - 1,00) - 3\,384 = 75 \times 9,81 \times 5 - 3\,384 = 3\,679 - 3\,384 = \SI{295}{J}$.
    \emph{Note :} Ce $W_{\text{mus}}$ est le travail \emph{net} fourni par les muscles sur le centre de gravité pendant tout le saut (phase d'appel + relevé). Il est positif mais relativement modeste car l'essentiel de l'énergie vient de la course ($E_{c,0}$). La perche permet de \emph{convertir} efficacement l'énergie horizontale en élévation verticale.
\end{enumerate}
\end{corrige}

\begin{corrige}[Exercice 11 — Centrale hydroélectrique de Nachtigal : de l'énergie potentielle à l'électricité]
\begin{enumerate}
    \item Débit massique : $\dot{m} = \rho Q = 1000 \times 350 = \SI{350\,000}{kg/s}$.
    \item Puissance hydraulique brute : $P_{\text{hyd}} = \dot{m} g H = 350\,000 \times 9,81 \times 35,5$.
    $9,81 \times 35,5 = 348,255$.
    $P_{\text{hyd}} = 350\,000 \times 348,255 = 121\,889\,250\,\text{W} \approx \SI{121,9}{MW}$.
    \item Puissance électrique : $P_{\text{el}} = \eta P_{\text{hyd}} = 0,92 \times 121,9 = \SI{112,1}{MW}$.
    \item Énergie journalière : $E_{\text{jour}} = P_{\text{el}} \times 24\,\text{h} = 112,1 \times 24 = \SI{2\,690,4}{MWh} = 2\,690\,400\,\text{kWh}$.
    Chiffre d'affaires : $CA = 2\,690\,400 \times 50 = \SI{134\,520\,000}{FCFA} \approx \SI{134,5}{millions de FCFA}$.
    \item \textbf{Bernoulli / Bilan énergie mécanique (par unité de masse) :}
    Amont (1) : $z_1$, $v_1 = 2,0\,\text{m/s}$, $P_1 = P_{\text{atm}}$.
    Aval (2) : $z_2 = z_1 - H$, $v_2 = 3,5\,\text{m/s}$, $P_2 = P_{\text{atm}}$.
    Travail spécifique extrait par la turbine $w_t$ (positif si énergie donnée par le fluide) :
    $g z_1 + \frac{v_1^2}{2} = g z_2 + \frac{v_2^2}{2} + w_t$ (pertes de charge négligées).
    $w_t = g(z_1 - z_2) + \frac{v_1^2 - v_2^2}{2} = gH + \frac{2,0^2 - 3,5^2}{2}$.
    $w_t = 9,81 \times 35,5 + \frac{4 - 12,25}{2} = 348,255 - 4,125 = \SI{344,13}{J/kg}$.
    Puissance mécanique turbine : $P_{\text{meca}} = \dot{m} w_t = 350\,000 \times 344,13 = 120\,445\,500\,\text{W} \approx \SI{120,4}{MW}$.
    \textbf{Cohérence :} $P_{\text{hyd}} = \dot{m} g H = \SI{121,9}{MW}$.
    La différence $P_{\text{hyd}} - P_{\text{meca}} = \dot{m} \frac{v_2^2 - v_1^2}{2} = 350\,000 \times 4,125 \approx \SI{1,44}{MW}$ correspond à l'augmentation de l'énergie cinétique de l'eau entre l'amont et l'aval (l'eau accélère dans le canal de fuite). L'énergie potentielle perdue ($mgH$) est convertie en travail turbine \emph{et} en énergie cinétique résiduelle à l'aval. Le bilan est cohérent.
    \item Énergie annuelle : $E_{\text{an}} = P_{\text{el}} \times 8\,000\,\text{h} = 112,1 \times 8\,000 = \SI{896\,800}{MWh} = \SI{896,8}{GWh}$.
    Consommation par foyer/an : $150\,\text{kWh/mois} \times 12 = \SI{1\,800}{kWh/an} = 1,8\,\text{MWh/an}$.
    Nombre de foyers : $N = \frac{896\,800}{1,8} \approx 498\,222$.
    La centrale peut alimenter environ \SI{500\,000}{foyers} (arrondi).
\end{enumerate}
\end{corrige}

\newpage

\coursec{Fiche bilan}

\begin{popfigure}
\centering
\begin{center}\tcbox[colback=poporange,colframe=poporange,coltext=white,arc=9pt,boxsep=4pt,left=6pt,right=6pt]{\bfseries\small Énergie Potentielle \& Mécanique}\end{center}
\vspace{-2pt}
\begin{tcbraster}[raster columns=3,raster equal height=rows,raster column skip=6pt,raster row skip=6pt,enhanced,blankest]
\begin{tcolorbox}[enhanced,colback=poporange,colframe=poporange,coltext=white,boxrule=0.9pt,arc=3pt,left=4pt,right=4pt,top=3pt,bottom=3pt,boxsep=1pt,halign=left]\bfseries\scriptsize Notion Générale (État, Forces conservatrices)\end{tcolorbox}
\begin{tcolorbox}[enhanced,colback=popgreen,colframe=popgreen,coltext=white,boxrule=0.9pt,arc=3pt,left=4pt,right=4pt,top=3pt,bottom=3pt,boxsep=1pt,halign=left]\bfseries\scriptsize Pesanteur $E_p = mgh$ (h: altitude)\end{tcolorbox}
\begin{tcolorbox}[enhanced,colback=poppurple,colframe=poppurple,coltext=white,boxrule=0.9pt,arc=3pt,left=4pt,right=4pt,top=3pt,bottom=3pt,boxsep=1pt,halign=left]\bfseries\scriptsize Élastique $E_p = \frac{1}{2}k(\Delta l)^2$ (Ressort linéaire)\end{tcolorbox}
\begin{tcolorbox}[enhanced,colback=poppink,colframe=poppink,coltext=white,boxrule=0.9pt,arc=3pt,left=4pt,right=4pt,top=3pt,bottom=3pt,boxsep=1pt,halign=left]\bfseries\scriptsize Énergie Mécanique $E_m = E_c + E_p$\end{tcolorbox}
\begin{tcolorbox}[enhanced,colback=popgold,colframe=popgold,coltext=white,boxrule=0.9pt,arc=3pt,left=4pt,right=4pt,top=3pt,bottom=3pt,boxsep=1pt,halign=left]\bfseries\scriptsize Conservation $\Delta E_m = 0$ si $W_{\text{nc}}=0$\end{tcolorbox}
\begin{tcolorbox}[enhanced,colback=popblue,colframe=popblue,coltext=white,boxrule=0.9pt,arc=3pt,left=4pt,right=4pt,top=3pt,bottom=3pt,boxsep=1pt,halign=left]\bfseries\scriptsize Non-Conservation $\Delta E_m = W_{\text{nc}}$ (Frottements, chocs)\end{tcolorbox}
\begin{tcolorbox}[enhanced,colback=popdark,colframe=popdark,coltext=white,boxrule=0.9pt,arc=3pt,left=4pt,right=4pt,top=3pt,bottom=3pt,boxsep=1pt,halign=left]\bfseries\scriptsize Méthodes Choix référentiel Bilan énergétique\end{tcolorbox}
\end{tcbraster}

\legende{Carte mentale : Énergie potentielle et énergie mécanique (Leçon 05 -- 1ère C)}
\end{popfigure}

\begin{bilan}
\courssub{Définitions clés}
\begin{itemize}
    \item \textbf{Énergie potentielle ($E_p$) :} Énergie liée à la \emph{configuration} (position, déformation) d'un système dans un champ de forces \cle{conservatrices}. Elle ne dépend que de l'état initial et final, pas du chemin.
    \item \textbf{Force conservative :} Force dont le travail ne dépend que des positions initiale et finale (pesanteur, force de rappel élastique). Le travail sur un cycle fermé est nul.
    \item \textbf{Énergie mécanique ($E_m$) :} Somme de l'énergie cinétique et de l'énergie potentielle du système : $E_m = E_c + E_p$.
\end{itemize}

\courssub{Formules à connaître par cœur}
\begin{center}
\begin{tabularx}{\linewidth}{|>{\centering\arraybackslash}X|>{\centering\arraybackslash}X|>{\centering\arraybackslash}X|}
\hline
\rowcolor{popblueL}
\textbf{Grandeurs} & \textbf{Expression} & \textbf{Unité SI} \\
\hline
Énergie cinétique & $E_c = \dfrac{1}{2} m v^2$ & Joule (\si{\joule}) \\
\hline
Énergie potentielle de pesanteur & $E_p = m g h$ & Joule (\si{\joule}) \\
\hline
Énergie potentielle élastique & $E_p = \dfrac{1}{2} k (\Delta l)^2$ & Joule (\si{\joule}) \\
\hline
Énergie mécanique & $E_m = E_c + E_p$ & Joule (\si{\joule}) \\
\hline
Variation énergie mécanique & $\Delta E_m = W_{\text{ext}}^{\text{nc}}$ & Joule (\si{\joule}) \\
\hline
Travail force non conservative & $W_{\text{nc}} = \int \vec{F}_{\text{nc}} \cdot d\vec{l}$ & Joule (\si{\joule}) \\
\hline
\end{tabularx}
\end{center}

\courssub{Théorèmes \& Principes (Probatoire : démonstration exigible pour le lien $W = \Delta E_c$ et $\Delta E_p = -W_{\text{cons}}$)}
\begin{itemize}
    \item \textbf{Théorème de l'énergie cinétique :} $W_{\text{ext}} = \Delta E_c = E_{c,f} - E_{c,i}$.
    \item \textbf{Relation $E_p$ -- Forces conservatrices :} $\Delta E_p = -W_{\text{cons}}$.
    \item \textbf{Principe de conservation de l'énergie mécanique :} Si le système ne subit que des forces conservatrices (ou si $W_{\text{nc}} = 0$), alors $E_m$ est constante : $E_{m,i} = E_{m,f}$.
    \item \textbf{Théorème de l'énergie mécanique (cas général) :} $\Delta E_m = W_{\text{nc}}$ (Travail des forces non conservatrices : frottements, moteur, choc inélastique).
\end{itemize}

\courssub{Méthodes express}
\begin{enumerate}
    \item \textbf{Choisir le référentiel d'étude} (galiléen lié au sol) et l'origine des altitudes ($h=0$) pour $E_p^{\text{pes}}$.
    \item \textbf{Lister les forces} : identifier conservatrices (poids, rappel ressort) et non conservatrices (frottements, normal si mobile, tension moteur).
    \item \textbf{Écrire $E_m$ initial et final} : $E_{m,i} = E_{c,i} + E_{p,i}$ ; $E_{m,f} = E_{c,f} + E_{p,f}$.
    \item \textbf{Appliquer le bon théorème} :
    \begin{itemize}
        \item Si $W_{\text{nc}} = 0$ : $E_{m,i} = E_{m,f}$ (Conservation).
        \item Sinon : $E_{m,f} - E_{m,i} = W_{\text{nc}}$ (souvent $W_{\text{frott}} = -f \cdot d$).
    \end{itemize}
    \item \textbf{Résoudre l'inconnue} (vitesse, hauteur, compression, coefficient de frottement).
    \item \textbf{Vérifier} : homogénéité dimensionnelle, ordre de grandeur, signe des travaux.
\end{enumerate}
\end{bilan}

\begin{aretenir}[Checklist avant le Probatoire]
\begin{itemize}
    \item $\square$ Je sais définir une force conservative et donner deux exemples (poids, rappel ressort).
    \item $\square$ Je sais écrire $E_p^{\text{pes}} = mgh$ en précisant le choix de l'origine ($h=0$).
    \item $\square$ Je sais écrire $E_p^{\text{élas}} = \frac{1}{2}k(\Delta l)^2$ pour un ressort de raideur $k$ (en \si{\newton\per\meter}).
    \item $\square$ Je sais calculer l'énergie mécanique $E_m = E_c + E_p$ à un instant donné.
    \item $\square$ Je sais énoncer la condition de conservation de $E_m$ ($W_{\text{nc}}=0$) et l'appliquer ($E_{m,i}=E_{m,f}$).
    \item $\square$ Je sais appliquer le théorème de l'énergie mécanique $\Delta E_m = W_{\text{nc}}$ quand il y a frottements ou choc.
    \item $\square$ Je sais calculer le travail d'une force de frottement constante : $W_{\text{frott}} = -f \cdot d$ (négatif !).
    \item $\square$ Je distingue clairement \emph{énergie potentielle} (énergie de position) et \emph{travail} (transfert d'énergie).
    \item $\square$ Je vérifie toujours l'homogénéité de mes formules : $[E] = \si{\kilogram\meter\squared\per\second\squared} = \si{\joule}$.
    \item $\square$ Je sais analyser un graphique $E_p(x)$, $E_c(x)$, $E_m(x)$ : pentes, intersections, zones interdites ($E_c < 0$).
    \item $\square$ Je sais modéliser un freinage (vélo, moto-taxi) ou un choc (balle, voiture) par un bilan énergétique.
\end{itemize}
\end{aretenir}

\findecours

\end{document}
